% Leçon 5 — Vocabulaire et compréhension de texte : délinquance juvénile — Chinois 1ère A — Cours signé OnBuch+
% Programme officiel MINESEC. Compiler avec : tectonic ZHA05-vocabulaire-delinquance-juvenile.tex
\newcommand{\DOCMATIERE}{Chinois}
\newcommand{\DOCNIVEAU}{1ère A}
\newcommand{\DOCMODULE}{Module 1 : Vie familiale et intégration sociale}
\newcommand{\DOCLECON}{ZHA05}
\newcommand{\DOCTITRE}{Leçon 5 — Vocabulaire et compréhension de texte : délinquance juvénile}
% OnBuch+ — préambule « cours signés » ENRICHI (compilé avec tectonic / XeLaTeX)
% Système visuel « Pop » d'OnBuch+ : fond crème (jamais blanc), bordures encre
% épaisses, ombres « dures » décalées, titres Archivo Black en capitales.
% Version MATHÉMATIQUES : graphes (pgfplots/tikz), boîtes pédagogiques
% (définition, propriété, méthode, attention, le savais-tu, exercice résolu,
% exercice, TP, bilan), page de garde et signature OnBuch+ ancrée en bas.
%
% Macros attendues dans main.tex AVANT \input{preamble} :
%   \DOCMATIERE  \DOCNIVEAU  \DOCMODULE  \DOCLECON (numéro)  \DOCTITRE
\documentclass[11pt]{article}

\usepackage{fontspec}
\usepackage[french]{babel} % français = langue principale ; le chinois via \zh{..} / textebox (xeCJK)
\usepackage{xeCJK}
\usepackage{pifont} % \ding{51} \ding{55} utilisés par les modèles
\usepackage[a4paper,margin=2cm,top=3.4cm,bottom=2.8cm,headheight=1.4cm,footskip=1.3cm]{geometry}
\usepackage{amsmath,amssymb,amsfonts}
\usepackage{mathtools} % \xmapsto, \coloneqq... souvent utilisés par les modèles
\usepackage{cancel} % simplifications barrées (bilans de forces)
% \abs{z} (module, valeur absolue) et \norm{u} (norme), utilisés par les modèles
\providecommand{\abs}{}\let\abs\relax\DeclarePairedDelimiter{\abs}{\lvert}{\rvert}
\providecommand{\norm}{}\let\norm\relax\DeclarePairedDelimiter{\norm}{\lVert}{\rVert}
\usepackage[table]{xcolor} % [table] : fournit aussi \rowcolors (alternance de lignes)
\usepackage{pagecolor}
\usepackage{tikz}
\usetikzlibrary{arrows.meta,positioning,calc,shapes.geometric,shapes.misc,shapes.arrows,decorations.pathmorphing,decorations.pathreplacing,decorations.markings,patterns,shadows,mindmap,trees,fit,backgrounds,matrix,intersections,angles,quotes}
\usepackage{pgfplots}
\usepackage[version=4]{mhchem} % équations nucléaires \ce{^{235}_{92}U}
\usepackage[european,siunitx]{circuitikz} % schémas électriques
\pgfplotsset{compat=1.18}
\usepgfplotslibrary{fillbetween,groupplots} % aires entre courbes (calcul intégral)
\usepackage{enumitem}
\usepackage{fancyhdr}
% [most] charge toutes les bibliothèques tcolorbox (listings, minted, raster,
% documentation...) et gonfle inutilement la mémoire TeX sur un document long
% (~300 boîtes) : on ne charge que ce qui est réellement utilisé.
\usepackage[skins,breakable]{tcolorbox}
\usepackage{graphicx}
\usepackage{colortbl}
\usepackage{booktabs}
\usepackage{tabularx}
\usepackage{diagbox} % case d'en-tête diagonale des tableaux à double entrée
% Colonnes extensibles centrée / gauche / droite (C, L, R), souvent utilisées par les modèles
\newcolumntype{C}{>{\centering\arraybackslash}X}
\newcolumntype{L}{>{\raggedright\arraybackslash}X}
\newcolumntype{R}{>{\raggedleft\arraybackslash}X}
\usepackage{array}
\usepackage{multicol}
\usepackage{multirow}
\usepackage{siunitx}
% parse-numbers=false : accepte aussi \SI{2,5 \times 10^{-3}}{...}
\sisetup{locale=FR,output-decimal-marker={,},group-minimum-digits=4,parse-numbers=false}
\DeclareSIUnit\ohm{\ensuremath{\symup{\Omega}}} % U+2126 absent de la police
\DeclareSIUnit\division{div} % réglages d'oscilloscope (ms/div, V/div)
\DeclareSIUnit\tr{tr} % tours (vitesse de rotation, ex. \SI{1500}{\tr\per\minute})
\DeclareSIUnit\molar{M} % concentration molaire (ex. \SI{3}{\molar}, \SI{1}{\micro\molar})
\DeclareSIUnit\an{an} % année (le modèle confond parfois avec \year, un compteur TeX) — ex. \SI{5}{\kilo\meter\per\an}
\DeclareSIUnit\FCFA{FCFA} % franc CFA (ex. \SI{500}{\FCFA\per\kilo\gram})
\DeclareSIUnit\degreeBrix{\textdegree Brix} % teneur en sucre d'un sirop/jus (réfractométrie)
\DeclareSIUnit\month{mois} % le modèle utilise parfois \month, un compteur TeX, comme unité de durée
\usepackage{qrcode}
% \degree : commande fréquemment utilisée par les modèles pour un angle ou une
% température en dehors de \SI{}{} (non fournie par les paquets chargés).
\providecommand{\degree}{^{\circ}}
% \qed (fin de démonstration), utilisé par les modèles sans amsthm
\providecommand{\qed}{\hfill$\square$}
% \barre : double barre des valeurs interdites dans un tableau de variations (tkz-tab)
\providecommand{\barre}{\ensuremath{\|}}
% environnement proof (amsthm non chargé) utilisé par les modèles
\ifdefined\proof\else\newenvironment{proof}[1][Démonstration]{\par\noindent\textit{#1.}\ }{\qed\par}\fi
\usepackage{lastpage}
\usepackage{hyperref}
\hypersetup{hidelinks}

% ============================================================
% Palette « Pop » OnBuch+ — mêmes hex que la classe P (plus_ui.dart)
% ============================================================
\definecolor{poporange}{HTML}{F59321}
\definecolor{popdark}{HTML}{E07A0C}
\definecolor{popcream}{HTML}{FFF6E8}
\definecolor{popink}{HTML}{1C1714}
\definecolor{poppurple}{HTML}{7A5AE0}
\definecolor{popgreen}{HTML}{2E9E6A}
\definecolor{poppink}{HTML}{E0457B}
\definecolor{popblue}{HTML}{2F7DE1}
\definecolor{popgold}{HTML}{C98A12}
\definecolor{popmuted}{HTML}{8A7F76}
\definecolor{popline}{HTML}{EADFD2}
\definecolor{popgray}{HTML}{8A7F76}
\colorlet{popgrey}{popgray}
% Alias tolérants (le modèle nomme parfois les couleurs différemment)
\colorlet{popwhite}{white}
\colorlet{popblack}{popink}
\colorlet{popred}{poppink}
\colorlet{popyellow}{popgold}
% Teintes claires (fonds de tableaux/encadrés) — le modèle nomme parfois
% les couleurs avec un suffixe L (ex. popblueL) sans les avoir définies.
\colorlet{poporangeL}{poporange!20}
\colorlet{popdarkL}{popdark!20}
\colorlet{poppurpleL}{poppurple!20}
\colorlet{popgreenL}{popgreen!20}
\colorlet{poppinkL}{poppink!20}
\colorlet{popblueL}{popblue!20}
\colorlet{popgoldL}{popgold!20}
\colorlet{popredL}{popred!20}
\colorlet{popgrayL}{popgray!20}
% Teintes claires pour fonds de tableaux / zones
\colorlet{popblueL}{popblue!10!white}
\colorlet{poporangeL}{poporange!14!white}
\colorlet{popgreenL}{popgreen!12!white}
\colorlet{poppurpleL}{poppurple!12!white}
\colorlet{poppinkL}{poppink!12!white}
\colorlet{popgoldL}{popgold!14!white}

\pagecolor{popcream}

% ============================================================
% Typographie
% ============================================================
\defaultfontfeatures{Ligatures=TeX}
\setmainfont{PlusJakartaSans-Medium.ttf}[Path=../../../fonts/,BoldFont=PlusJakartaSans-Bold.ttf]
% Symboles, API et emojis absents de la police principale : repli sur DejaVu Sans (caractères actifs)
\newfontface\symfont{DejaVuSans.ttf}[Path=../../../fonts/]
\makeatletter
\newcommand\symactive[1]{\catcode#1=13 \begingroup\lccode`\~=#1 \lowercase{\endgroup\def~}{{\symfont\char#1}}}
\newcommand\symblank[1]{\catcode#1=13 \begingroup\lccode`\~=#1 \lowercase{\endgroup\def~}{}}
\newcommand\symhyph[1]{\catcode#1=13 \begingroup\lccode`\~=#1 \lowercase{\endgroup\def~}{\mbox{-}}}
\newcount\symc
\newcommand\symrange[2]{\symc=#1 \@whilenum\symc<#2 \do{\expandafter\symactive\expandafter{\the\symc}\advance\symc by 1 }}
\newcommand\symblankrange[2]{\symc=#1 \@whilenum\symc<#2 \do{\expandafter\symblank\expandafter{\the\symc}\advance\symc by 1 }}
\makeatother
\symrange{"0250}{"02FF}\symactive{"2713}\symactive{"2714}\symactive{"2717}\symactive{"2718}\symactive{"2610}\symactive{"2611}\symactive{"2612}
\symrange{"2600}{"27BF}
\catcode"2705=13 \begingroup\lccode`\~="2705 \lowercase{\endgroup\def~}{{\symfont\char"2713}}\symblank{"26BD}\symblank{"26A1}
\symrange{"2460}{"257F}\symactive{"223C}\symactive{"2248}\symactive{"2260}
\symhyph{"2011}\symhyph{"2E1A}\symblankrange{"1F300}{"1FAFF}\symblank{"FE0F}
% Caractères chinois : WenQuanYi Zen Hei (sous-ensemble GB2312 + pinyin), gras/italique simulés
\setCJKmainfont{WenQuanYiZenHei-subset.ttf}[Path=../../../fonts/,AutoFakeBold=3,AutoFakeSlant=0.2]
\xeCJKsetup{CJKspace=false}
% Pinyin : voyelles à caron (ǎ ǐ ǒ ǔ ǖ ǘ ǚ ǜ) absentes de la police principale -> caractères actifs
\newfontface\pyfont{WenQuanYiZenHei-subset.ttf}[Path=../../../fonts/]
\newcommand\pyactive[1]{\catcode#1=13 \begingroup\lccode`\~=#1 \lowercase{\endgroup\def~}{{\pyfont\char#1}}}
\pyactive{"01CD}\pyactive{"01CE}\pyactive{"01CF}\pyactive{"01D0}\pyactive{"01D1}\pyactive{"01D2}\pyactive{"01D3}\pyactive{"01D4}\pyactive{"01D5}\pyactive{"01D6}\pyactive{"01D7}\pyactive{"01D8}\pyactive{"01D9}\pyactive{"01DA}\pyactive{"01DB}\pyactive{"01DC}
% Maths en sans-serif (Fira Math) avec chiffres et lettres droites de Plus
% Jakarta, pour que les formules se fondent dans le texte.
\usepackage[math-style=ISO,bold-style=ISO]{unicode-math}
\setmathfont{FiraMath-Regular.otf}[Path=../../../fonts/]
\setmathfont{PlusJakartaSans-Medium.ttf}[Path=../../../fonts/,range={up/num,up/{Latin,latin},\mathup}]
\usepackage{icomma} % virgule décimale sans espace en mode math
% Caractères mathématiques Unicode absents des polices de texte (Plus Jakarta,
% Archivo Black) : rendus par leur équivalent mathématique, en texte comme en
% formule — sinon ils s'affichent en carré vide (ex. « Arithmétique dans ℤ »).
\usepackage{newunicodechar}
\newunicodechar{ℝ}{\ensuremath{\mathbb{R}}}
\newunicodechar{ℤ}{\ensuremath{\mathbb{Z}}}
\newunicodechar{ℂ}{\ensuremath{\mathbb{C}}}
\newunicodechar{ℕ}{\ensuremath{\mathbb{N}}}
\newunicodechar{ℚ}{\ensuremath{\mathbb{Q}}}
\newunicodechar{𝔻}{\ensuremath{\mathbb{D}}}
\newunicodechar{α}{\ensuremath{\alpha}}
\newunicodechar{β}{\ensuremath{\beta}}
\newunicodechar{γ}{\ensuremath{\gamma}}
\newunicodechar{δ}{\ensuremath{\delta}}
\newunicodechar{ε}{\ensuremath{\varepsilon}}
\newunicodechar{θ}{\ensuremath{\theta}}
\newunicodechar{λ}{\ensuremath{\lambda}}
\newunicodechar{μ}{\ensuremath{\mu}}
\newunicodechar{σ}{\ensuremath{\sigma}}
\newunicodechar{τ}{\ensuremath{\tau}}
\newunicodechar{φ}{\ensuremath{\varphi}}
\newunicodechar{ψ}{\ensuremath{\psi}}
\newunicodechar{ω}{\ensuremath{\omega}}
\newunicodechar{Σ}{\ensuremath{\Sigma}}
% majuscules produites par \MakeUppercase dans les titres (α-aminés -> Α)
\newunicodechar{Α}{\ensuremath{\alpha}}
\newunicodechar{Β}{\ensuremath{\beta}}
\newunicodechar{Γ}{\ensuremath{\Gamma}}
\newunicodechar{Δ}{\ensuremath{\Delta}}
\newunicodechar{Ε}{\ensuremath{\varepsilon}}
\newunicodechar{Θ}{\ensuremath{\Theta}}
\newunicodechar{Λ}{\ensuremath{\Lambda}}
\newunicodechar{Μ}{\ensuremath{\mu}}
\newunicodechar{Τ}{\ensuremath{\tau}}
\newunicodechar{Φ}{\ensuremath{\Phi}}
\newunicodechar{Ψ}{\ensuremath{\Psi}}
\newunicodechar{Ω}{\ensuremath{\Omega}}
\newunicodechar{↦}{\ensuremath{\mapsto}}
\newunicodechar{∈}{\ensuremath{\in}}
\newunicodechar{∉}{\ensuremath{\notin}}
\newunicodechar{∧}{\ensuremath{\wedge}}
\newunicodechar{⇔}{\ensuremath{\Leftrightarrow}}
\newunicodechar{⟺}{\ensuremath{\Longleftrightarrow}}
\newunicodechar{⇒}{\ensuremath{\Rightarrow}}
\newunicodechar{⟹}{\ensuremath{\Longrightarrow}}
\newunicodechar{∘}{\ensuremath{\circ}}
\newunicodechar{∪}{\ensuremath{\cup}}
\newunicodechar{∩}{\ensuremath{\cap}}
\newunicodechar{≡}{\ensuremath{\equiv}}
\newunicodechar{⊂}{\ensuremath{\subset}}
\newunicodechar{⊥}{\ensuremath{\perp}}
\newunicodechar{∃}{\ensuremath{\exists}}
\newunicodechar{∀}{\ensuremath{\forall}}
\newunicodechar{∛}{\ensuremath{\sqrt[3]{\,}}}
\newunicodechar{∜}{\ensuremath{\sqrt[4]{\,}}}
\newunicodechar{⃗}{\ensuremath{{}^{\to}}}
\newunicodechar{ⁿ}{\ifmmode^{n}\else\textsuperscript{\ensuremath{n}}\fi}
\newunicodechar{ˣ}{\ifmmode^{x}\else\textsuperscript{\ensuremath{x}}\fi}
\newunicodechar{ᵇ}{\ifmmode^{b}\else\textsuperscript{\ensuremath{b}}\fi}
\newunicodechar{ʳ}{\ifmmode^{r}\else\textsuperscript{\ensuremath{r}}\fi}
\newunicodechar{ᵘ}{\ifmmode^{u}\else\textsuperscript{\ensuremath{u}}\fi}
\newunicodechar{ʸ}{\ifmmode^{y}\else\textsuperscript{\ensuremath{y}}\fi}
\newunicodechar{ᵃ}{\ifmmode^{a}\else\textsuperscript{\ensuremath{a}}\fi}
\newunicodechar{ᵉ}{\ifmmode^{e}\else\textsuperscript{\ensuremath{e}}\fi}
\newunicodechar{ᵏ}{\ifmmode^{k}\else\textsuperscript{\ensuremath{k}}\fi}
\newunicodechar{ᵖ}{\ifmmode^{p}\else\textsuperscript{\ensuremath{p}}\fi}
\newunicodechar{ᴺ}{\ifmmode^{N}\else\textsuperscript{\ensuremath{N}}\fi}
\newunicodechar{ᶜ}{\ifmmode^{c}\else\textsuperscript{\ensuremath{c}}\fi}
\newunicodechar{ʰ}{\ifmmode^{h}\else\textsuperscript{\ensuremath{h}}\fi}
\newunicodechar{ᵠ}{\ifmmode^{q}\else\textsuperscript{\ensuremath{q}}\fi}
\newunicodechar{ₙ}{\ifmmode_{n}\else\textsubscript{\ensuremath{n}}\fi}
\newunicodechar{ₐ}{\ifmmode_{a}\else\textsubscript{\ensuremath{a}}\fi}
\newunicodechar{ᵢ}{\ifmmode_{i}\else\textsubscript{\ensuremath{i}}\fi}
\newunicodechar{ᵣ}{\ifmmode_{r}\else\textsubscript{\ensuremath{r}}\fi}
\newunicodechar{ₖ}{\ifmmode_{k}\else\textsubscript{\ensuremath{k}}\fi}
\newunicodechar{ᵦ}{\ifmmode_{\beta}\else\textsubscript{\ensuremath{\beta}}\fi}
\newunicodechar{ₓ}{\ifmmode_{x}\else\textsubscript{\ensuremath{x}}\fi}
\newfontfamily\popdisplay{ArchivoBlack-Regular.ttf}[Path=../../../fonts/]
\newfontfamily\poplogo{SpaceGrotesk-Bold.ttf}[Path=../../../fonts/]
\newfontfamily\popsemi{PlusJakartaSans-SemiBold.ttf}[Path=../../../fonts/]
\newfontfamily\popxbold{PlusJakartaSans-ExtraBold.ttf}[Path=../../../fonts/]
\newfontfamily\popxboldls{PlusJakartaSans-ExtraBold.ttf}[Path=../../../fonts/,LetterSpace=3.0]

\setlist[enumerate]{leftmargin=1.7em,itemsep=5pt,topsep=4pt}
\setlist[itemize]{leftmargin=1.7em,itemsep=4pt,topsep=4pt,label={\color{poporange}\textbullet}}
\setlength{\parindent}{0pt}
\setlength{\parskip}{5pt}
\renewcommand{\arraystretch}{1.25}

% pgfplots : style Pop par défaut
\pgfplotsset{
  every axis/.append style={axis line style={popink,line width=1pt},tick style={popink},
    grid style={popline,line width=0.5pt},label style={font=\small},tick label style={font=\footnotesize}},
  popcurve/.style={poporange,line width=1.8pt,no markers,smooth},
}
% Même style disponible dans un \draw TikZ simple (hors pgfplots)
\tikzset{popcurve/.style={poporange,line width=1.8pt}}
\tikzset{popgrid/.style={popline,very thin}}
\colorlet{popcurve}{poporange} % le nom du style est parfois utilisé comme couleur

% ============================================================
% Pill / squiggle
% ============================================================
\newcommand{\poppill}[3]{%
  \tikz[baseline=-0.55ex]\node[draw=popink,line width=1.2pt,rounded corners=2.6pt,%
    fill=#2,inner xsep=6.5pt,inner ysep=3pt]{%
    \popxboldls\fontsize{7}{7}\selectfont\color{#3}\MakeUppercase{#1}};%
}
\newcommand{\popsquiggle}[1]{%
  \tikz[baseline=-0.4ex]{
    \draw[#1,line width=2.6pt,line cap=round]
      plot [smooth] coordinates {(0,0.09) (0.34,0.01) (0.68,0.21) (1.02,0.01) (1.36,0.10)};
  }%
}
% Mot-clé surligné (marqueur orange sous le texte)
\newcommand{\cle}[1]{{\popxbold\color{popdark}#1}}

% ============================================================
% En-tête / pied
% ============================================================
\pagestyle{fancy}
\fancyhf{}
\renewcommand{\headrulewidth}{0pt}
\renewcommand{\footrulewidth}{0pt}
\lhead{\raisebox{-0.15ex}{\poplogo\fontsize{13}{13}\selectfont\color{popink}OnBuch\color{poporange}+}}
\rhead{\poppill{\DOCMATIERE\ \textbullet\ \DOCNIVEAU\ \textbullet\ Leçon \DOCLECON}{white}{popink}}
\lfoot{\popxbold\fontsize{7.6}{7.6}\selectfont\color{popmuted}\MakeUppercase{Cours signé OnBuch+}\ \textbullet\ {\popsemi\DOCTITRE}}
\rfoot{\tikz[baseline=-0.6ex]\node[rounded corners=8.5pt,fill=popink,minimum height=17pt,minimum width=17pt,inner xsep=5pt,inner sep=0]{\popxbold\fontsize{8}{8}\selectfont\color{popcream}\thepage\,/\,\pageref*{LastPage}};}

% ============================================================
% Titres
% ============================================================
\newcommand{\chaptitre}[2]{%
  \begin{tcolorbox}[enhanced,colback=poporange,colframe=popink,boxrule=2.6pt,arc=11pt,
      drop shadow=popink,left=15pt,right=15pt,top=12pt,bottom=11pt,before skip=3pt,after skip=18pt]
    \poppill{#2}{popink}{popcream}\par\vspace{9pt}
    {\raggedright\hyphenpenalty=10000\exhyphenpenalty=10000\popdisplay\fontsize{22}{24}\selectfont\color{popcream}\MakeUppercase{#1}\par}
  \end{tcolorbox}
}
% Section numérotée : pastille orange avec le numéro + titre Archivo
\newcounter{popsec}
\newcommand{\coursec}[1]{%
  \stepcounter{popsec}\setcounter{popsub}{0}%
  \par\vspace{16pt}\noindent
  \begin{minipage}{\linewidth}
  \tikz[baseline=-0.7ex]\node[draw=popink,line width=1.6pt,fill=poporange,rounded corners=4pt,minimum size=22pt,inner sep=0]{\popdisplay\fontsize{12}{12}\selectfont\color{popcream}\thepopsec};%
  \hspace{8pt}{\popdisplay\fontsize{15}{17}\selectfont\color{popink}\MakeUppercase{#1}}\par
  \vspace{4pt}\noindent{\color{poporange}\rule{36pt}{3.6pt}}
  \end{minipage}\par\nopagebreak\vspace{8pt}
}
\newcounter{popsub}
\newcommand{\courssub}[1]{%
  \stepcounter{popsub}%
  \par\vspace{9pt}\noindent{\popxbold\fontsize{11.5}{13}\selectfont\color{popink}{\color{poporange}\thepopsec.\thepopsub}\hspace{6pt}#1}\par\nopagebreak\vspace{4pt}
}

% ============================================================
% Boîtes pédagogiques — même recette : fond blanc, bordure épaisse colorée,
% ombre dure encre, badge plein. Titre optionnel [#1].
% ============================================================
\tcbset{popbase/.style={breakable,enhanced,colback=white,boxrule=2.3pt,arc=9pt,
  shadow={3pt}{-3pt}{0pt}{popink},left=14pt,right=14pt,top=11pt,bottom=11pt,before skip=11pt,after skip=15pt,
  fontupper=\popsemi\fontsize{10.6}{14.5}\selectfont\color{popink},
  fontlower=\popsemi\fontsize{10.6}{14.5}\selectfont\color{popink}}}
% Coche dessinée (le glyphe ✓ n'existe pas dans les polices du document)
\newcommand{\popcheck}[1]{\tikz[baseline=-0.5ex]\draw[#1,line width=1.6pt,line cap=round,line join=round](-0.13,0.0)--(-0.04,-0.09)--(0.13,0.1);}
\providecommand{\coche}{\popcheck{popgreen}}
\providecommand{\resultat}[1]{\par\smallskip\noindent\fcolorbox{popgreen}{popgreenL}{\parbox{\dimexpr\linewidth-2\fboxsep-2\fboxrule\relax}{\textbf{Résultat.} #1}}\par\smallskip}
\newcommand{\popboxhead}[3]{\poppill{#1}{#2}{white}\ifx\relax#3\relax\else\hspace{6pt}{\popxbold\fontsize{10.6}{12}\selectfont\color{popink}#3}\fi\par\vspace{7pt}}

\newtcolorbox{aretenir}[1][]{popbase,colframe=popgreen,before upper={\popboxhead{À retenir}{popgreen}{#1}}}
% Texte en chinois (document de lecture, dialogue, extrait) : cadre
% d'étude. Un titre optionnel (source, auteur) s'affiche dans l'en-tête.
\newtcolorbox{textebox}[1][]{popbase,colframe=popblue,before upper={\popboxhead{文本 · Texte}{popblue}{#1}},after upper={}}
% Marques ✓ / ✗ (forme correcte / fautive) souvent utilisées par les modèles
\providecommand{\cross}{\textcolor{poppink}{\ensuremath{\times}}}
\providecommand{\xmark}{\cross}\providecommand{\crossmark}{\cross}\providecommand{\cmark}{\ensuremath{\checkmark}}
% Ligne de réponse / justification à compléter (inventée par les modèles)
\providecommand{\justification}[1]{\par\noindent\textit{Justification~:} \dotfill\par}
% \textipa (package tipa non chargé) : la police ne couvre pas l'API -> simple passage
\providecommand{\textipa}[1]{#1}
% Soulignement (ulem non chargé) : \uline, \ul
\providecommand{\uline}[1]{\underline{#1}}\providecommand{\ul}[1]{\underline{#1}}
% \glossaire{\item ...} inventée par les modèles : liste de glossaire
\providecommand{\glossaire}[1]{\begin{itemize}#1\end{itemize}}
% \alert (beamer) inventée par les modèles : texte mis en évidence
\providecommand{\alert}[1]{\textbf{\textcolor{poppink}{#1}}}
% variantes ASCII d'ordinaux écrites en macro
\providecommand{\iere}{\textsuperscript{re}}\providecommand{\ieme}{\textsuperscript{e}}\providecommand{\ere}{\textsuperscript{re}}\providecommand{\eme}{\textsuperscript{e}}\providecommand{\er}{\textsuperscript{er}}\providecommand{\ie}{\textsuperscript{e}}
% Ordinaux français écrits en macro par les modèles : 1\ière, 3\ième
\newcommand{\ière}{\textsuperscript{re}}\newcommand{\ième}{\textsuperscript{e}}
% \exemple (inventée par les modèles) : étiquette « Exemple : »
\providecommand{\exemple}{\textbf{Exemple~:}\ }
% \square utilisable aussi hors mode math (cases à cocher)
\AtBeginDocument{\let\squareMath\square\renewcommand{\square}{\ensuremath{\squareMath}}}
% Mot ou phrase en chinois à l'intérieur d'un paragraphe français
\newcommand{\zh}[1]{\ifmmode\text{#1}\else{#1}\fi}
\let\eng\zh \let\esp\zh \let\all\zh % alias tolérés
% flèches d'intonation inventées par les modèles
\providecommand{\textsearrow}{}\renewcommand{\textsearrow}{\ensuremath{\searrow}}\providecommand{\textnearrow}{}\renewcommand{\textnearrow}{\ensuremath{\nearrow}}
% \fr{..} : mot français (inventé par les modèles)
\providecommand{\fr}[1]{\textit{#1}}
% zhbox : encadré de texte chinois inventé par les modèles
\newenvironment{zhbox}{\begin{quote}}{\end{quote}}
% \vs : « versus » inventé par les modèles
\providecommand{\vs}{\textit{vs}\ }
% \blank : case à compléter inventée par les modèles
\providecommand{\blank}[1][]{\underline{\hspace{1.6em}}}
\newtcolorbox{exemplebox}[1][]{popbase,colframe=poporange,before upper={\popboxhead{Exemple}{poporange}{#1}}}
\newtcolorbox{definition}[1][]{popbase,colframe=popblue,before upper={\popboxhead{Définition}{popblue}{#1}}}
\newtcolorbox{propriete}[1][]{popbase,colframe=poppurple,before upper={\popboxhead{Propriété}{poppurple}{#1}}}
\newtcolorbox{methode}[1][]{popbase,colframe=popgold,before upper={\popboxhead{Méthode}{popgold}{#1}}}
\newtcolorbox{attention}[1][]{popbase,colframe=poppink,before upper={\popboxhead{Attention}{poppink}{#1}}}
\newtcolorbox{experience}[1][]{popbase,colframe=popblue,colback=popblueL,before upper={\popboxhead{Expérience}{popblue}{#1}}}
% « Le savais-tu ? » : carte violette pleine (contexte, culture, Cameroun)
\newtcolorbox{savaistu}[1][]{popbase,colback=poppurple,colframe=popink,
  fontupper=\popsemi\fontsize{10.4}{14.2}\selectfont\color{white},
  before upper={\poppill{Le savais-tu\,?}{white}{poppurple}\ifx\relax#1\relax\else\hspace{6pt}{\popxbold\color{white}#1}\fi\par\vspace{7pt}}}
% Exercice résolu : énoncé en haut, \tcblower puis la solution sur fond crème
\newcounter{popexo}\newcounter{popexr}
\newtcolorbox{exoresolu}[1][]{popbase,colframe=poporange,colbacklower=popcream,
  segmentation style={popink,line width=1.2pt,dashed},
  before upper={\stepcounter{popexr}\popboxhead{Exercice résolu \thepopexr}{poporange}{#1}},
  before lower={\poppill{Solution}{popink}{popcream}\par\vspace{6pt}}}
% Exercice (non résolu en place) — la correction est dans la partie Corrigés
\newtcolorbox{exercice}[1][]{popbase,colframe=popink,
  before upper={\stepcounter{popexo}\popboxhead{Exercice \thepopexo}{popink}{#1}}}
% Corrigé
\newtcolorbox{corrige}[1][]{popbase,colframe=popgreen,colback=popgreenL,
  before upper={\popboxhead{Corrigé}{popgreen}{#1}}}
% Objectifs / prérequis / bilan
\newtcolorbox{objectifs}{popbase,colframe=popink,colback=poporangeL,
  before upper={\popboxhead{Objectifs de la leçon}{popink}{}}}
\newtcolorbox{prerequis}{popbase,colframe=popink,colback=popblueL,
  before upper={\popboxhead{Prérequis}{popblue}{}}}
\newtcolorbox{bilan}{popbase,colframe=popink,colback=white,boxrule=2.8pt,
  before upper={\popboxhead{Fiche bilan}{poporange}{L'essentiel à connaître pour l'examen}}}
% Figure encadrée avec légende
\newtcolorbox{popfigure}[1][]{enhanced,breakable=false,colback=white,colframe=popline,boxrule=1.4pt,arc=8pt,
  left=10pt,right=10pt,top=10pt,bottom=8pt,before skip=10pt,after skip=12pt,halign=center,
  lower separated=false,halign lower=center,
  fontlower=\popsemi\fontsize{9}{11.5}\selectfont\color{popmuted},
  #1}
\newcommand{\legende}[1]{\par\vspace{6pt}{\popsemi\fontsize{9}{11.5}\selectfont\color{popmuted}#1}}

% ============================================================
% Signature OnBuch+
% ============================================================
\newcommand{\poplogoinline}[1]{{\poplogo\fontsize{#1}{#1}\selectfont\color{popink}OnBuch\color{poporange}+}}
\newcommand{\signatureonbuch}{%
  \begin{tcolorbox}[enhanced,colback=popink,colframe=popink,boxrule=2.2pt,arc=10pt,
      drop shadow=poporange,left=14pt,right=14pt,top=10pt,bottom=10pt,before skip=0pt,after skip=0pt]
    \tikz[baseline=-0.7ex]\node[circle,fill=popgreen,draw=popcream,line width=1.3pt,minimum size=22pt,inner sep=0]{\popcheck{white}};%
    \hspace{9pt}\begin{minipage}[c]{0.62\linewidth}
      {\popxbold\fontsize{12}{13}\selectfont\color{popcream}Signé\ \textbullet\ {\poplogo OnBuch\color{poporange}+}}\par\vspace{2pt}
      {\raggedright\popsemi\fontsize{8}{10.5}\selectfont\color{popline}\MakeUppercase{Équipe pédagogique OnBuch+}\\\MakeUppercase{Conforme au programme officiel MINESEC}\par}
    \end{minipage}\hfill
    {\poplogo\fontsize{22}{22}\selectfont\color{popcream}OnBuch\color{poporange}+}
  \end{tcolorbox}%
}

% ============================================================
% Page de garde
% #1 titre · #2 matière · #3 niveau · #4 module · #5 n° leçon · #6 durée ·
% #7 description / objectifs (texte ou itemize)
% ============================================================
\newcommand{\pagegarde}[7]{%
  \thispagestyle{empty}
  \newgeometry{a4paper,margin=1.7cm,top=1.6cm,bottom=1.4cm}%
  \begin{tikzpicture}[remember picture,overlay]
    \fill[poporange!18] ([xshift=-3.2cm,yshift=-3.6cm]current page.north east) circle (4.6cm);
    \fill[poppurple!14] ([xshift=2.4cm,yshift=7.6cm]current page.south west) circle (3.8cm);
  \end{tikzpicture}%
  {\poplogo\fontsize{30}{30}\selectfont\color{popink}OnBuch\color{poporange}+}\hfill
  \raisebox{0.6ex}{\poppill{Cours signé \textbullet\ Premium}{popink}{popcream}}\par
  \vspace{1pt}\popsquiggle{poppurple}\par
  \vspace{8pt}
  \begin{tcolorbox}[enhanced,colback=poporange,colframe=popink,boxrule=2.8pt,arc=13pt,
      drop shadow=popink,left=18pt,right=18pt,top=12pt,bottom=13pt,after skip=10pt]
    \poppill{#2 \textbullet\ #3}{popink}{popcream}\hspace{5pt}\poppill{#4}{white}{popink}\par\vspace{8pt}
    {\popdisplay\fontsize{14}{14}\selectfont\color{popink}LEÇON #5}\par\vspace{4pt}
    {\raggedright\hyphenpenalty=10000\exhyphenpenalty=10000\popdisplay\fontsize{25}{27}\selectfont\color{popcream}\MakeUppercase{#1}\par}
    \vspace{8pt}
    {\popxbold\fontsize{9.5}{9.5}\selectfont\color{popink}\MakeUppercase{Durée conseillée : #6}}
  \end{tcolorbox}
  \begin{tcolorbox}[enhanced,colback=white,colframe=popink,boxrule=2.2pt,arc=10pt,
      drop shadow=popink,left=16pt,right=16pt,top=10pt,bottom=10pt,after skip=8pt]
    \poppill{Dans cette leçon}{poppurple}{white}\par\vspace{5pt}
    \popsemi\fontsize{9.3}{12.6}\selectfont\color{popink}#7
  \end{tcolorbox}
  \noindent\begin{minipage}[t]{0.487\linewidth}
    \begin{tcolorbox}[enhanced,colback=popgreen,colframe=popink,boxrule=2.2pt,arc=10pt,
        drop shadow=popink,left=11pt,right=11pt,top=10pt,bottom=10pt,height=3.1cm,valign=center]
      \tcbox[colback=white,colframe=popink,boxrule=1.3pt,arc=5pt,boxsep=0pt,nobeforeafter,
        left=3pt,right=3pt,top=3pt,bottom=3pt,valign=center]{\qrcode[height=1.9cm]{https://play.google.com/store/apps/details?id=com.luvvix.onbuch&hl=fr}}%
      \hspace{8pt}\begin{minipage}[c]{0.5\linewidth}
        {\popdisplay\fontsize{11}{12.5}\selectfont\color{white}\MakeUppercase{Télécharge OnBuch}}\par\vspace{4pt}
        {\raggedright\popsemi\fontsize{8.4}{11}\selectfont\color{white}Gratuit \textbullet\ Google Play\\Tuteur IA, annales, concours, résultats\par}
      \end{minipage}
    \end{tcolorbox}
  \end{minipage}\hfill
  \begin{minipage}[t]{0.487\linewidth}
    \begin{tcolorbox}[enhanced,colback=poppurple,colframe=popink,boxrule=2.2pt,arc=10pt,
        drop shadow=popink,left=11pt,right=11pt,top=10pt,bottom=10pt,height=3.1cm,valign=center]
      \tikz[baseline=-0.7ex]\node[circle,fill=white,draw=popink,line width=1.3pt,minimum size=1.6cm,inner sep=0]{\popdisplay\fontsize{24}{24}\selectfont\color{poppurple}+};%
      \hspace{8pt}\begin{minipage}[c]{0.6\linewidth}
        {\popdisplay\fontsize{11}{12.5}\selectfont\color{white}\MakeUppercase{Passe à OnBuch+}}\par\vspace{4pt}
        {\raggedright\popsemi\fontsize{8.4}{11}\selectfont\color{white}Tous les cours signés, Tuteur IA illimité, fiches PDF — onglet OnBuch+ de l'app\par}
      \end{minipage}
    \end{tcolorbox}
  \end{minipage}
  \vspace{8pt}
  \popcontenu
  \vfill
  \signatureonbuch
  \restoregeometry
  \newpage
}

% Rangée « Ce cours contient » (page de garde)
\newcommand{\poptile}[3]{% #1 couleur #2 gros texte #3 légende
  \begin{tcolorbox}[enhanced,colback=#1,colframe=popink,boxrule=2pt,arc=9pt,shadow={2.5pt}{-2.5pt}{0pt}{popink},
      left=6pt,right=6pt,top=9pt,bottom=9pt,halign=center,valign=center,height=1.75cm,width=\linewidth,nobeforeafter]
    {\popdisplay\fontsize{12.5}{13}\selectfont\color{white}\MakeUppercase{#2}}\par\vspace{3pt}
    {\centering\popsemi\fontsize{7.8}{9.5}\selectfont\color{white}#3\par}
  \end{tcolorbox}}
\newcommand{\popcontenu}{%
  \par\noindent{\popxboldls\fontsize{7.5}{7.5}\selectfont\color{popmuted}CE COURS SIGNÉ CONTIENT}\par\vspace{7pt}
  \noindent\begin{minipage}[t]{0.235\linewidth}\poptile{popblue}{Cours}{complet, illustré, pas à pas}\end{minipage}\hfill
  \begin{minipage}[t]{0.235\linewidth}\poptile{poporange}{Méthodes}{et exercices résolus}\end{minipage}\hfill
  \begin{minipage}[t]{0.235\linewidth}\poptile{poppink}{Exercices}{type examen + corrigés}\end{minipage}\hfill
  \begin{minipage}[t]{0.235\linewidth}\poptile{popgreen}{Fiche bilan}{l'essentiel à retenir}\end{minipage}\par}

% Sommaire
\newtcolorbox{sommaire}{enhanced,colback=white,colframe=popink,boxrule=2.2pt,arc=10pt,
  drop shadow=popink,left=18pt,right=18pt,top=13pt,bottom=13pt,before skip=6pt,after skip=20pt,
  before upper={\poppill{Sommaire}{poppurple}{white}\par\vspace{9pt}\popsemi\fontsize{10.6}{14.5}\selectfont\color{popink}}}

% Signature de fin de cours — poussée tout en bas de la dernière page
\newcommand{\findecours}{%
  \par\vfill\nopagebreak
  \begin{center}\popsquiggle{poporange}\end{center}\vspace{4pt}
  \signatureonbuch
}
\AtBeginDocument{\def\llbracket{\mathopen{[\mkern-3.5mu[}}\def\rrbracket{\mathclose{]\mkern-3.5mu]}}} % intervalles d'entiers (glyphe ⟦ absent de la police)
% Glyphes absents de Fira Math : redéfinis à partir de glyphes disponibles
\AtBeginDocument{%
  \def\setminus{\mathbin{\backslash}}\let\smallsetminus\setminus
  \def\vdots{\mathord{\vbox{\baselineskip3.2pt\lineskiplimit0pt\kern4pt\hbox{.}\hbox{.}\hbox{.}}}}%
  \def\ddots{\mathinner{\mkern1mu\raise6.5pt\hbox{.}\mkern2mu\raise3.6pt\hbox{.}\mkern2mu\raise0.7pt\hbox{.}\mkern1mu}}%
  \def\triangle{\mathord{\scalebox{0.9}{$\symup{\Delta}$}}}%
  \def\nabla{\mathord{\raisebox{\depth}{\scalebox{1}[-1]{$\symup{\Delta}$}}}}%
  \def\checkmark{\popcheck{popdark}}%
  \def\star{\mathbin{\ast}}\def\bigstar{\mathord{\ast}}%
  \def\diamond{\mathbin{\scalebox{0.6}{$\Diamond$}}}%
  \def\bigcup{\mathop{\mathchoice{\raisebox{-0.3em}{\scalebox{1.7}{$\cup$}}}{\scalebox{1.25}{$\cup$}}{\cup}{\cup}}\displaylimits}%
  \def\bigcap{\mathop{\mathchoice{\raisebox{-0.3em}{\scalebox{1.7}{$\cap$}}}{\scalebox{1.25}{$\cap$}}{\cap}{\cap}}\displaylimits}%
  \def\lesseqqgtr{\mathrel{\lessgtr}}\def\bigoplus{\mathop{\oplus}\displaylimits}%
}
\providecommand{\ssi}{si et seulement si} % abréviation fréquente
\AtBeginDocument{\providecommand{\wideparen}{\overparen}} % arc de cercle

%% FIGFIX-BEGIN v5 — correctifs globaux des figures (docs/onbuch-generation/tools/figfix)
\usepackage{adjustbox}\usepackage{etoolbox}\tcbuselibrary{raster}
% toute figure TikZ/pgfplots plus large que la ligne est réduite à la largeur disponible
\BeforeBeginEnvironment{tikzpicture}{\begin{adjustbox}{max width=\linewidth}}
\AfterEndEnvironment{tikzpicture}{\end{adjustbox}}
% légendes pgfplots lisibles par-dessus les courbes
\pgfplotsset{every axis/.append style={legend style={fill=white,fill opacity=0.92,text opacity=1,draw=black!15,inner sep=3pt}}}
% étiquettes d'arêtes (syntaxe edge["..."]) avec fond blanc
\tikzset{every edge quotes/.append style={fill=white,inner sep=1.2pt}}
%% FIGFIX-END
\begin{document}

\pagegarde{Leçon 5 — Vocabulaire et compréhension de texte : délinquance juvénile}{Chinois}{1ère A}{Module 1 : Vie familiale et intégration sociale}{ZHA05}{2 h}{Cette leçon de vocabulaire et compréhension de texte porte sur la délinquance juvénile (青少年犯罪). À travers un texte support original, un glossaire structuré, l'étude des radicaux et de l'ordre des traits, l'élève apprend à lire, prononcer et employer le lexique du thème et à répondre à des questions de compréhension.
\vspace{4pt}
\begin{multicols}{2}\small
\begin{itemize}
\item À la fin de cette leçon, je sais lire à voix haute un court texte chinois sur la délinquance juvénile en respectant les tons.
\item Je sais reconnaître et traduire au moins 15 mots de vocabulaire du thème (犯罪, 青少年, 法律, 原因, 家庭, 学校, 社会, 帮助, 教育, 保护...).
\item Je sais identifier les radicaux (clés) des caractères du thème et les utiliser pour deviner le sens.
\item Je sais écrire correctement quelques caractères clés en respectant l'ordre des traits.
\item Je sais répondre par écrit à des questions de compréhension sur le texte.
\item Je sais employer les collocations du thème (预防犯罪, 违法, 受到惩罚...) dans des phrases complètes.
\end{itemize}\end{multicols}}

\begin{sommaire}
\begin{multicols}{2}
\begin{enumerate}
\item Mise en route et présentation du thème
\item Texte support : lecture guidée
\item Glossaire du thème
\item Clés (radicaux) et ordre des traits
\item Compréhension du texte
\item Production guidée et réinvestissement
\item Activité d'intégration
\item Méthodes et exercices résolus
\item Exercices
\item Corrigés des exercices
\item Fiche bilan
\end{enumerate}
\end{multicols}
\end{sommaire}

\begin{objectifs}
\begin{itemize}
\item À la fin de cette leçon, je sais lire à voix haute un court texte chinois sur la délinquance juvénile en respectant les tons.
\item Je sais reconnaître et traduire au moins 15 mots de vocabulaire du thème (犯罪, 青少年, 法律, 原因, 家庭, 学校, 社会, 帮助, 教育, 保护...).
\item Je sais identifier les radicaux (clés) des caractères du thème et les utiliser pour deviner le sens.
\item Je sais écrire correctement quelques caractères clés en respectant l'ordre des traits.
\item Je sais répondre par écrit à des questions de compréhension sur le texte.
\item Je sais employer les collocations du thème (预防犯罪, 违法, 受到惩罚...) dans des phrases complètes.
\item Je sais produire un court paragraphe guidé donnant les causes et les solutions de la délinquance juvénile.
\item Je sais comparer brièvement la situation au Cameroun et en Chine avec le vocabulaire appris.
\end{itemize}
\end{objectifs}

\begin{prerequis}
\begin{itemize}
\item Vocabulaire de la famille et de l'école vu en Seconde (家, 父母, 学校, 学生, 老师).
\item Structure de base Sujet-Verbe-Objet et particule 的.
\item Lecture du pinyin et maîtrise des quatre tons.
\item Radicaux courants déjà rencontrés (人, 女, 子, 口, 心).
\item Expressions de cause et de conséquence simples (因为...所以...).
\end{itemize}
\end{prerequis}

\newpage

\coursec{Mise en route et présentation du thème}

\courssub{Situation d'accroche et problématique}

À Yaoundé comme à Douala, les faits divers relatent régulièrement des affaires de vols ou de bagarres impliquant des adolescents déscolarisés. En Chine, ce phénomène est pris très au sérieux : la loi sur la prévention de la délinquance juvénile \zh{预防未成年人犯罪法} (\emph{Yùfáng wèichéngniánrén fànzuì fǎ}) encadre la protection et l'éducation des mineurs. Comment un jeune Chinois parle-t-il de ce sujet ? Cette leçon nous donne les mots pour lire, comprendre et discuter d'un texte sur la délinquance juvénile.

\textbf{Problématique de la leçon :} comment nommer, en chinois, un problème de société complexe, ses causes, ses conséquences et ses solutions, afin de comprendre un texte d'information ?

\courssub{Brainstorming : qu'est-ce que la délinquance juvénile ?}

Avant d'ouvrir le chinois, parlons français. En classe entière, répondons à deux questions :

\begin{enumerate}
\item Qu'est-ce que la \cle{délinquance juvénile} ? (Des actes interdits par la loi — vols, violences, trafics — commis par des mineurs, c'est-à-dire des jeunes de moins de 18 ans.)
\item Quelles en sont les \cle{causes possibles} ? La pauvreté, les problèmes familiaux (divorce, absence des parents), la drogue et l'alcool, les mauvaises fréquentations, la déscolarisation, le chômage.
\end{enumerate}

Chaque réponse française trouvera bientôt son équivalent chinois. Observons d'abord la phrase-clé de la leçon :

\begin{textebox}[Phrase-clé du thème]
青少年犯罪是一个严重的社会问题。\\
\emph{Qīngshàonián fànzuì shì yí ge yánzhòng de shèhuì wèntí.}\\
« La délinquance juvénile est un grave problème social. »
\end{textebox}

\begin{attention}[Le classificateur 个 et la liaison de 一]
Dans \zh{一个严重的社会问题}, le classificateur général \zh{个} (\emph{ge}) est obligatoire entre le numéral \zh{一} et le nom \zh{问题} : on ne dit jamais \zh{一问题}. Remarquez aussi la prononciation : \zh{一} seul se lit \emph{yī}, mais devant un quatrième ton il devient \emph{yí} : \emph{yí ge}. C'est la règle de changement de ton de \zh{一}.
\end{attention}

\courssub{Le titre chinois : 青少年犯罪}

Le thème du jour se dit \zh{青少年犯罪} (\emph{qīngshàonián fànzuì}). Découpons ce mot composé, car en chinois, comprendre les morceaux, c'est comprendre le tout :

\begin{center}
\begin{tabularx}{\linewidth}{|X|X|X|X|}
\hline
\rowcolor{popblueL} Caractères & Pinyin & Nature & Traduction \\
\hline
青少年 & qīngshàonián & nom & adolescent, jeunesse \\
\hline
犯罪 & fànzuì & verbe / nom & commettre un crime / crime, délit \\
\hline
青少年犯罪 & qīngshàonián fànzuì & nom & délinquance juvénile \\
\hline
严重 & yánzhòng & adjectif & grave, sérieux \\
\hline
社会 & shèhuì & nom & société \\
\hline
问题 & wèntí & nom & problème, question \\
\hline
\end{tabularx}
\end{center}

\begin{definition}[犯罪 fànzuì et 罪 zuì]
\zh{犯罪} (\emph{fànzuì}) est à la fois un \textbf{verbe} (« commettre un crime ») et un \textbf{nom} (« crime, délit »). Il est formé de \zh{犯} (\emph{fàn}, « violer, commettre ») et de \zh{罪} (\emph{zuì}, « crime, faute »). Exemple : \zh{他没有犯罪。} (\emph{Tā méiyǒu fànzuì.}) « Il n'a pas commis de crime. » Notez la négation du passé avec \zh{没有} (\emph{méiyǒu}), sans \zh{了}.
\end{definition}

\begin{definition}[青少年 qīngshàonián et 未成年人 wèichéngniánrén]
\zh{青少年} (\emph{qīngshàonián}) désigne les adolescents et les jeunes ; il est formé de \zh{青} (\emph{qīng}, « vert, jeune ») et de \zh{少年} (\emph{shàonián}, « jeune garçon, adolescent »). Le terme juridique, lui, est \zh{未成年人} (\emph{wèichéngniánrén}), littéralement « personne pas encore adulte » : \zh{未} (\emph{wèi}, « pas encore ») + \zh{成年} (\emph{chéngnián}, « devenir adulte ») + \zh{人} (\emph{rén}, « personne »).
\end{definition}

\begin{propriete}[Construire un mot composé : nom + verbe = notion]
Le chinois fabrique de grandes notions en soudant des mots simples : \zh{青少年} (jeunesse) + \zh{犯罪} (commettre un crime) = « délinquance juvénile ». Le français utilise un adjectif savant (\emph{juvénile}) ; le chinois reste concret : littéralement « les jeunes commettent des crimes ». Retenez ce mécanisme : il vous permettra de deviner le sens de nombreux mots nouveaux, par exemple \zh{社会问题} (société + problème = « problème social »).
\end{propriete}

\begin{attention}[Piège pour les francophones]
Ne traduisez pas \zh{犯罪} par « criminel » : le criminel (la personne) se dit \zh{罪犯} (\emph{zuìfàn}) ou \zh{犯人} (\emph{fànrén}). \zh{犯罪} désigne l'\emph{action} de commettre un crime ou le crime lui-même, jamais la personne. Comparez : \zh{他犯罪了。} (\emph{Tā fànzuì le.}) « Il a commis un crime. » et \zh{他是一个罪犯。} (\emph{Tā shì yí ge zuìfàn.}) « C'est un criminel. »
\end{attention}

\begin{savaistu}[L'âge de la responsabilité pénale]
En Chine, l'âge de la responsabilité pénale est fixé en principe à 16 ans, et à 14 ans pour les crimes les plus graves. Au Cameroun, la responsabilité pénale des mineurs est également encadrée par un régime spécial. Dans les deux pays, la justice des mineurs privilégie l'éducation et la réinsertion plutôt que la prison : on dit en chinois que l'on veut \zh{帮助} (\emph{bāngzhù}, « aider ») les jeunes, pas seulement les punir.
\end{savaistu}

\courssub{Carte mentale du thème}

Organisons dès maintenant les trois axes autour desquels tout texte sur ce sujet est construit : les \cle{causes}, les \cle{conséquences} et les \cle{solutions}.

\begin{popfigure}
\begin{tikzpicture}[mindmap, grow cyclic, every node/.style={concept, align=center, font=\small},
root concept/.append style={concept color=popblue, font=\bfseries},
level 1/.append style={level distance=3.2cm, sibling angle=120},
level 2/.append style={level distance=2.4cm, sibling angle=45}]
\node[root concept, align=center]{青少年犯罪\\ \emph{délinquance}\\ \emph{juvénile}}
  child[concept color=poporange] { node[align=center]{原因\\ \emph{causes}}
    child { node[concept color=poporangeL, text=popdark, align=center]{贫穷\\ \emph{pauvreté}} }
    child { node[concept color=poporangeL, text=popdark, align=center]{家庭问题\\ \emph{problèmes}\\ \emph{familiaux}} }
    child { node[concept color=poporangeL, text=popdark, align=center]{毒品\\ \emph{drogue}} }
  }
  child[concept color=poppink] { node[align=center]{后果\\ \emph{conséquences}}
    child { node[concept color=poppinkL, text=popdark, align=center]{失学\\ \emph{déscolarisation}} }
    child { node[concept color=poppinkL, text=popdark, align=center]{进监狱\\ \emph{prison}} }
  }
  child[concept color=popgreen] { node[align=center]{解决办法\\ \emph{solutions}}
    child { node[concept color=popgreenL, text=popdark, align=center]{教育\\ \emph{éducation}} }
    child { node[concept color=popgreenL, text=popdark, align=center]{法律\\ \emph{loi}} }
    child { node[concept color=popgreenL, text=popdark, align=center]{家庭关爱\\ \emph{soutien}\\ \emph{familial}} }
  };
\end{tikzpicture}
\legende{Carte mentale lexicale : les trois axes du thème 青少年犯罪.}
\end{popfigure}

\begin{exemplebox}[Lire les mots de la carte mentale]
\begin{itemize}
\item \zh{原因} (\emph{yuányīn}) : « cause, raison ». \zh{犯罪的原因是什么？} (\emph{Fànzuì de yuányīn shì shénme?}) « Quelle est la cause de la délinquance ? »
\item \zh{后果} (\emph{hòuguǒ}) : « conséquence ». \zh{后果很严重。} (\emph{Hòuguǒ hěn yánzhòng.}) « Les conséquences sont graves. »
\item \zh{解决办法} (\emph{jiějué bànfǎ}) : « solution ». \zh{我们必须找到解决办法。} (\emph{Wǒmen bìxū zhǎodào jiějué bànfǎ.}) « Nous devons trouver des solutions. »
\end{itemize}
\end{exemplebox}

\courssub{Première écoute du texte support}

\begin{methode}[Avant de lire un texte chinois]
\begin{enumerate}
\item Lisez d'abord le \textbf{titre} : il annonce le sujet (ici, \zh{青少年犯罪}).
\item Lisez ensuite la \textbf{première phrase} : elle donne souvent l'idée générale (ici : « un grave problème social »).
\item \textbf{Prédisez} le contenu : le texte va probablement parler des causes, des conséquences et des solutions.
\item À la première écoute (sans support écrit), ne cherchez pas à tout comprendre : repérez seulement les mots connus et l'idée générale.
\end{enumerate}
\end{methode}

L'enseignant lit maintenant le texte support à voix haute, deux fois, \textbf{sans que vous regardiez l'écrit}. Écoutez les sons, les tons, et essayez de reconnaître les mots de la carte mentale : \zh{原因}, \zh{后果}, \zh{解决办法}, \zh{家庭}, \zh{学校}, \zh{教育}…

\courssub{Questions de repérage global}

Après cette première écoute, répondez oralement, en français ou avec les mots chinois que vous connaissez déjà :

\begin{enumerate}
\item \textbf{De qui parle le texte ?} (De jeunes, d'adolescents : \zh{青少年}.)
\item \textbf{Quel problème est décrit ?} (Certains jeunes commettent des délits : \zh{犯罪}.)
\item \textbf{Quelles solutions sont proposées ?} (L'éducation, la famille, l'école, la loi : \zh{教育、家庭、学校、法律}.)
\end{enumerate}

\begin{exoresolu}[Repérage global à l'écoute]
Énoncé : à la première écoute du texte, vous avez relevé les mots \zh{家庭} (\emph{jiātíng}), \zh{毒品} (\emph{dúpǐn}) et \zh{教育} (\emph{jiàoyù}). Classez-les dans la carte mentale : cause, conséquence ou solution ?
\tcblower
Solution détaillée : \zh{家庭} (la famille) et \zh{毒品} (la drogue) appartiennent à la branche \zh{原因} (causes) : les problèmes familiaux et la drogue poussent certains jeunes vers la délinquance. \zh{教育} (l'éducation) appartient à la branche \zh{解决办法} (solutions) : éduquer les jeunes, c'est prévenir la délinquance. Aucun de ces trois mots n'est une conséquence (\zh{后果}) : une conséquence serait par exemple la prison (\zh{进监狱}) ou la déscolarisation (\zh{失学}).
\end{exoresolu}

\begin{exercice}[À vous de jouer]
Classez les mots suivants dans les trois branches de la carte mentale (causes, conséquences, solutions) : \zh{贫穷} (\emph{pínqióng}, pauvreté), \zh{进监狱} (\emph{jìn jiānyù}, entrer en prison), \zh{法律} (\emph{fǎlǜ}, loi), \zh{失学} (\emph{shīxué}, déscolarisation).
\end{exercice}

\begin{corrige}[Exercice : classification lexicale]
Causes (\zh{原因}) : \zh{贫穷} (la pauvreté pousse certains jeunes à voler). Conséquences (\zh{后果}) : \zh{进监狱} (entrer en prison) et \zh{失学} (quitter l'école). Solutions (\zh{解决办法}) : \zh{法律} (la loi protège et encadre les mineurs).
\end{corrige}

\begin{aretenir}[L'essentiel de la mise en route]
\begin{itemize}
\item \zh{青少年犯罪} (\emph{qīngshàonián fànzuì}) = « délinquance juvénile » : jeunesse + commettre un crime.
\item Tout texte sur ce sujet s'organise autour de trois axes : \zh{原因} (causes), \zh{后果} (conséquences), \zh{解决办法} (solutions).
\item Stratégie de lecture : titre, première phrase, prédiction, puis écoute globale.
\item Attention au classificateur : \zh{一个问题}, jamais \zh{一问题} ; et \zh{一} devient \emph{yí} devant un quatrième ton.
\end{itemize}
\end{aretenir}

Nous sommes prêts : passons à la lecture guidée du texte support.

\coursec{Mise en route et présentation du thème}

\begin{experience}[Activation des connaissances : la jeunesse face aux défis]
Le professeur projette ou écrit au tableau trois photos légendées en chinois (sans traduction) : \zh{快乐的家庭} (famille heureuse), \zh{紧张的考试} (examen stressant), \zh{不良少年团伙} (bande de jeunes délinquants). 
\begin{enumerate}
\item \textbf{Observation} : En binômes, les élèves identifient les mots connus (\zh{家庭}, \zh{考试}, \zh{少年}) et devinent le sens global de chaque image.
\item \textbf{Mise en commun} : Le professeur guide vers le thème : \zh{青少年犯罪} (\emph{qīngshàonián fànzuì}) — \cle{délinquance juvénile}. On écrit le titre au tableau avec le pinyin et la traduction.
\item \textbf{Brainstorming lexical} : Chaque binôme note 3 à 5 mots chinois (caractères + pinyin) liés au thème (ex. : \zh{警察}, \zh{监狱}, \zh{教育}, \zh{父母}, \zh{学校}). Partage oral ; le professeur complète au tableau en corrigeant les tons.
\end{enumerate}
\textbf{Transition :} « Aujourd'hui, nous allons lire un texte qui analyse les causes de ce phénomène et propose des solutions. Mais d'abord, rappelons la méthode de lecture. »
\end{experience}

\coursec{Texte support : lecture guidée}

Dans la section précédente, nous avons activé nos connaissances sur la famille (\zh{家庭}), l'école (\zh{学校}) et la société (\zh{社会}) face au phénomène de la \cle{délinquance juvénile} (\zh{青少年犯罪}, \emph{qīngshàonián fànzuì}). Nous allons maintenant aborder le texte support en suivant une démarche rigoureuse en quatre temps : écoute globale, lecture chorale, analyse phrase par phrase, et repérage de la logique argumentative.

\courssub{Première approche : écoute et lecture chorale}

\begin{methode}[Comment aborder un texte chinois inconnu ?]
\begin{enumerate}
\item \textbf{Écoute globale (livres fermés)} : Le professeur lit le texte à vitesse normale. Les élèves ne regardent pas le texte ; ils identifient le ton général (grave, urgent, espoir) et repèrent les mots « phares » entendus (\zh{犯罪}, \zh{严重}, \zh{政府}...).
\item \textbf{Lecture chorale (livres ouverts)} : Toute la classe lit ensemble en suivant le pointeur du professeur. Insistance sur la prosodie : groupes respiratoires, tons, liaison \zh{一个} (\emph{yí ge} / \emph{yì ge} selon le ton suivant).
\item \textbf{Lecture individuelle diagnostique} : Chaque élève lit une phrase. Le professeur corrige immédiatement les tons (surtout les 3\up{e} tons consécutifs, les 4\up{e} tons qui « tombent » trop bas).
\item \textbf{Traduction littérale puis libre} : Phrase par phrase, on donne le sens mot à mot, puis la formulation française naturelle.
\end{enumerate}
\end{methode}

\begin{attention}[Points de vigilance phonétique — Deux mots « pièges »]
\begin{itemize}
\item \zh{犯罪} \emph{fànzuì} : Deux 4\up{e} tons (\emph{fàn-zuì}). La voix doit descendre franchement deux fois. Erreur fréquente : monter sur \emph{zuì} (confusion avec 3\up{e} ton) ou adoucir le deuxième ton. \textbf{Exercice} : Répéter \emph{fàn-zuì, fàn-zuì} en tapant la main sur la cuisse à chaque syllabe pour marquer la chute.
\item \zh{严重} \emph{yánzhòng} : 2\up{e} ton (\emph{yán}, montée nette) + 4\up{e} ton (\emph{zhòng}, chute nette). Erreur : dire \emph{yǎnzhòng} (3\up{e} ton) ou \emph{yánzhǒng} (3\up{e} ton). \textbf{Astuce} : \zh{重} (\emph{zhòng}, « lourd, grave ») se distingue de \zh{种} (\emph{zhǒng}, « espèce, graine ») par le ton seulement.
\end{itemize}
Rappel : En chinois, \textbf{un ton faux = un mot faux}. Le français n'utilise pas les tons pour distinguer le sens lexical ; cet apprentissage demande un entraînement quotidien (5 min / séance).
\end{attention}

\courssub{Le texte : lecture phrase par phrase avec analyse}

\begin{textebox}[青少年犯罪 — La délinquance juvénile]
现在，青少年犯罪是一个严重的社会问题。\\
\emph{Xiànzài, qīngshàonián fànzuì shì yí ge yánzhòng de shèhuì wèntí.}\\
« Aujourd'hui, la délinquance juvénile est un grave problème de société. »\\[4pt]
有些青少年因为家庭问题、学习压力或者坏朋友的影响而违法。\\
\emph{Yǒuxiē qīngshàonián yīnwèi jiātíng wèntí, xuéxí yālì huòzhě huài péngyou de yǐngxiǎng ér wéifǎ.}\\
« Certains adolescents enfreignent la loi à cause de problèmes familiaux, de la pression scolaire ou de l'influence de mauvais amis. »\\[4pt]
他们有的偷东西，有的打架，还有的吸毒。\\
\emph{Tāmen yǒude tōu dōngxi, yǒude dǎjià, hái yǒude xīdú.}\\
« Certains volent des choses, d'autres se battent, d'autres encore se droguent. »\\[4pt]
犯罪的原因很多：父母太忙，没有时间关心孩子；学校只重视成绩，不重视教育；社会上也有不好的影响。\\
\emph{Fànzuì de yuányīn hěn duō : fùmǔ tài máng, méiyǒu shíjiān guānxīn háizi ; xuéxiào zhǐ zhòngshì chéngjì, bù zhòngshì jiàoyù ; shèhuì shang yě yǒu bù hǎo de yǐngxiǎng.}\\
« Les causes de la délinquance sont nombreuses : les parents sont trop occupés et n'ont pas le temps de se soucier de leurs enfants ; l'école ne valorise que les notes et néglige l'éducation ; la société exerce aussi de mauvaises influences. »\\[4pt]
为了解决这个问题，家庭、学校和社会应该一起努力。\\
\emph{Wèile jiějué zhège wèntí, jiātíng, xuéxiào hé shèhuì yīnggāi yìqǐ nǔlì.}\\
« Pour résoudre ce problème, la famille, l'école et la société doivent travailler ensemble. »\\[4pt]
父母要关心孩子，学校要加强教育，政府要保护青少年，帮助他们走上正确的道路。\\
\emph{Fùmǔ yào guānxīn háizi, xuéxiào yào jiāqiáng jiàoyù, zhèngfǔ yào bǎohù qīngshàonián, bāngzhù tāmen zǒushang zhèngquè de dàolù.}\\
« Les parents doivent se soucier de leurs enfants, l'école doit renforcer l'éducation, le gouvernement doit protéger les adolescents et les aider à prendre le bon chemin. »\\[4pt]
只有这样，才能减少青少年犯罪。\\
\emph{Zhǐyǒu zhèyàng, cái néng jiǎnshǎo qīngshàonián fànzuì.}\\
« C'est seulement ainsi que l'on pourra réduire la délinquance juvénile. »
\end{textebox}

\courssub{Observation grammaticale : la structure d'énumération 有的...有的...还有的...}

Reprenons la troisième phrase : \zh{他们有的偷东西，有的打架，还有的吸毒。} 
Que remarquez-vous ? Le mot \zh{有的} (\emph{yǒude}) revient trois fois. Il ne signifie pas « avoir » ici, mais « certains (d'entre eux) ». C'est une structure figée pour énumérer des sous-groupes aux comportements différents.

\begin{propriete}[Structure d'énumération : 有的... 有的... (还有的...)]
\textbf{Schéma :} (Groupe de référence) + 有的 + V₁ / SV₁ , 有的 + V₂ / SV₂ , (还有的 + V₃ / SV₃) 。

\textbf{Emploi :} Présenter une répartition interne à un ensemble. Équivalent français : « Les uns..., les autres..., d'autres encore... » ou « Certains..., d'autres... ».

\textbf{Contraintes syntaxiques :}
\begin{itemize}
\item Le groupe de référence (\zh{他们}, \zh{学生们}, \zh{青少年}...) est énoncé \textbf{une seule fois} en tête (sujet de la phrase).
\item \zh{有的} fonctionne comme un \textbf{sujet partiel} : il se place \textbf{devant le verbe} (ou l'adjectif prédicatif), jamais seul en fin de phrase.
\item \zh{还有的} (\emph{hái yǒude}) introduit le dernier élément de l'énumération (« et d'autres encore ») ; on peut s'arrêter à deux \zh{有的}.
\end{itemize}
\end{propriete}

\begin{exemplebox}[Exemples authentiques — Caractères + Pinyin + Traduction]
\begin{itemize}
\item \zh{班上有的同学喜欢看书，有的喜欢听音乐。} \emph{Bān shang yǒude tóngxué xǐhuan kànshū, yǒude xǐhuan tīng yīnyuè.} « Dans la classe, certains camarades aiment lire, d'autres aiment écouter de la musique. »
\item \zh{市场上，有的水果很贵，有的很便宜。} \emph{Shìchǎng shang, yǒude shuǐguǒ hěn guì, yǒude hěn piányi.} « Au marché, certains fruits sont chers, d'autres sont bon marché. » (Ici \zh{有的} + adjectif).
\item \zh{周末，同学们有的去踢足球，有的去图书馆，还有的在家打游戏。} \emph{Zhōumò, tóngxuémen yǒude qù tī zúqiú, yǒde qù túshūguǎn, hái yǒude zài jiā dǎ yóuxì.} « Le week-end, certains vont jouer au foot, d'autres à la bibliothèque, d'autres encore jouent aux jeux vidéo à la maison. »
\end{itemize}
\end{exemplebox}

\begin{attention}[Faux amis phonétiques : 吸毒 vs 吸烟]
\begin{center}
\begin{tabularx}{\linewidth}{|X|X|X|X|}
\hline
\rowcolor{popblueL} Caractères & Pinyin & Nature & Traduction \\
\hline
\zh{吸毒} & \emph{xīdú} & v. & se droguer (litt. « aspirer du poison ») \\
\hline
\zh{吸烟} & \emph{xīyān} & v. & fumer (litt. « aspirer de la fumée/tabac ») \\
\hline
\end{tabularx}
\end{center}
Les deux commencent par \zh{吸} (\emph{xī}, 1\up{er} ton, « aspirer, inhaler »). Le second caractère change radicalement le sens : \zh{毒} (\emph{dú}, poison/drogue) vs \zh{烟} (\emph{yān}, tabac/fumée). Dans le texte, il s'agit de \zh{吸毒} (acte illégal, délit). Ne pas écrire \zh{吸烟} pour « se droguer » !
\end{attention}

\courssub{Repérage de la logique textuelle : les connecteurs argumentatifs}

\begin{methode}[Décrypter le plan d'un texte argumentatif chinois]
Soulignez dans le texte ces trois connecteurs structurants ; ils dessinent le squelette logique :
\begin{enumerate}
\item \zh{因为} (\emph{yīnwèi}) — « parce que, à cause de » $\rightarrow$ introduit la \textbf{CAUSE} (phrase 2 : \zh{因为家庭问题...而违法}). Souvent corrélé à \zh{所以} (\emph{suǒyǐ}) « donc » (absent ici, remplacé par \zh{而} \emph{ér} « et alors » formel).
\item \zh{为了} (\emph{wèile}) — « pour, afin de » $\rightarrow$ introduit le \textbf{BUT} (phrase 5 : \zh{为了解决这个问题}).
\item \zh{只有...才...} (\emph{zhǐyǒu... cái...}) — « seulement si... alors... » $\rightarrow$ exprime la \textbf{CONDITION NÉCESSAIRE} (phrase 7 : \zh{只有这样，才能减少...}). \zh{只有} = « seulement avoir » ; \zh{才} = « seulement alors » (adverbe de restriction placé devant le verbe).
\end{enumerate}
\textbf{Plan du texte} : 1. Constat (phr. 1) → 2. Causes multiples (phr. 2-4) → 3. Appel à l'action collective (phr. 5-6) → 4. Condition de réussite (phr. 7).
\end{methode}

\coursec{Glossaire du thème : vocabulaire ciblé}

Le tableau ci-dessous regroupe le lexique essentiel du texte et du thème « Délinquance juvénile ». La colonne \textbf{Nature} indique la catégorie grammaticale (n. = nom, v. = verbe, adj. = adjectif, adv. = adverbe, num. = numéral, clas. = classificateur, part. = particule). Les mots marqués \cle{★} sont \textbf{nouveaux} pour cette leçon ; les autres sont \textbf{connus} (réactivation).

\begin{center}
\begin{tabularx}{\linewidth}{|X|X|X|X|}
\hline
\rowcolor{popblueL} \textbf{Caractères} & \textbf{Pinyin} & \textbf{Nature} & \textbf{Traduction} \\
\hline
\zh{现在} & \emph{xiànzài} & adv. / n. & maintenant, actuellement ; le présent \\
\hline
\zh{青少年} & \emph{qīngshàonián} & n. & adolescent(s), jeune(s) \cle{★} \\
\hline
\zh{犯罪} & \emph{fànzuì} & v. / n. & commettre un crime ; délinquance, criminalité \cle{★} \\
\hline
\zh{是} & \emph{shì} & v. & être (verbe copule) \\
\hline
\zh{一个} & \emph{yí ge} / \emph{yì ge} & num.+clas. & un (classificateur \zh{个}) \\
\hline
\zh{严重} & \emph{yánzhòng} & adj. & grave, sérieux, sévère \cle{★} \\
\hline
\zh{的} & \emph{de} & part. & marque de la relation modificatif-nom (possessif, attributif) \\
\hline
\zh{社会} & \emph{shèhuì} & n. & société \\
\hline
\zh{问题} & \emph{wèntí} & n. & problème, question \\
\hline
\zh{有些} & \emph{yǒuxiē} & pron. & certains, quelques-uns \\
\hline
\zh{因为} & \emph{yīnwèi} & conj. & parce que, à cause de \\
\hline
\zh{家庭} & \emph{jiātíng} & n. & famille, foyer \\
\hline
\zh{学习} & \emph{xuéxí} & v. / n. & étudier, apprentissage ; études \\
\hline
\zh{压力} & \emph{yālì} & n. & pression (physique ou psychologique) \cle{★} \\
\hline
\zh{或者} & \emph{huòzhě} & conj. & ou, ou bien (alternative) \\
\hline
\zh{坏} & \emph{huài} & adj. & mauvais, cassé, gâté \\
\hline
\zh{朋友} & \emph{péngyou} & n. & ami \\
\hline
\zh{影响} & \emph{yǐngxiǎng} & v. / n. & influencer ; influence \cle{★} \\
\hline
\zh{而} & \emph{ér} & conj. & et (formel, lie cause $\rightarrow$ conséquence) ; mais (adversatif) \cle{★} \\
\hline
\zh{违法} & \emph{véifǎ} & v. & enfreindre la loi, agir illégalement \cle{★} \\
\hline
\zh{他们} & \emph{tāmen} & pron. & ils, elles \\
\hline
\zh{有的} & \emph{yǒude} & pron. & certains (d'entre eux) \cle{★} \\
\hline
\zh{偷} & \emph{tōu} & v. & voler, dérober \cle{★} \\
\hline
\zh{东西} & \emph{dōngxi} & n. & chose, objet, truc \\
\hline
\zh{打架} & \emph{dǎjià} & v. & se battre, bagarrer \cle{★} \\
\hline
\zh{还有的} & \emph{hái yǒude} & pron. & d'autres encore \cle{★} \\
\hline
\zh{吸毒} & \emph{xīdú} & v. & se droguer, consommer de la drogue \cle{★} \\
\hline
\zh{原因} & \emph{yuányīn} & n. & cause, raison, origine \cle{★} \\
\hline
\zh{很多} & \emph{hěn duō} & adj. & beaucoup (de) \\
\hline
\zh{父母} & \emph{fùmǔ} & n. & parents (père et mère) \\
\hline
\zh{太} & \emph{tài} & adv. & trop, excessivement \\
\hline
\zh{忙} & \emph{máng} & adj. & occupé \\
\hline
\zh{没有} & \emph{méiyǒu} & v. & ne pas avoir ; n'avoir pas (négation de \zh{有}) \\
\hline
\zh{时间} & \emph{shíjiān} & n. & temps \\
\hline
\zh{关心} & \emph{guānxīn} & v. & se soucier de, s'inquiéter pour, prendre soin de \cle{★} \\
\hline
\zh{孩子} & \emph{háizi} & n. & enfant \\
\hline
\zh{学校} & \emph{xuéxiào} & n. & école \\
\hline
\zh{只} & \emph{zhǐ} & adv. & seulement, rien que \\
\hline
\zh{重视} & \emph{zhòngshì} & v. & attacher de l'importance à, valoriser \cle{★} \\
\hline
\zh{成绩} & \emph{chéngjì} & n. & résultats scolaires, notes, performance \\
\hline
\zh{不} & \emph{bù} & adv. & non, ne ... pas (négation standard) \\
\hline
\zh{教育} & \emph{jiàoyù} & v. / n. & éduquer ; éducation \\
\hline
\zh{社会上} & \emph{shèhuì shang} & loc. & dans la société, sur le plan social \\
\hline
\zh{也} & \emph{yě} & adv. & aussi, également \\
\hline
\zh{有} & \emph{yǒu} & v. & avoir, il y a \\
\hline
\zh{不好} & \emph{bù hǎo} & adj. & pas bon, mauvais \\
\hline
\zh{为了} & \emph{wèile} & prep. & pour, afin de (but) \cle{★} \\
\hline
\zh{解决} & \emph{jiějué} & v. & résoudre (un problème, un conflit) \cle{★} \\
\hline
\zh{这个} & \emph{zhège} & pron. & ce, cet, cette \\
\hline
\zh{和} & \emph{hé} & conj. & et (coordonne noms/pronoms) \\
\hline
\zh{应该} & \emph{yīnggāi} & v. aux. & devoir, doit (obligation morale/conseil) \\
\hline
\zh{一起} & \emph{yìqǐ} & adv. & ensemble \\
\hline
\zh{努力} & \emph{nǔlì} & v. / adj. & faire des efforts, travailler dur ; assidu \\
\hline
\zh{要} & \emph{yào} & v. aux. & devoir, vouloir, aller (futur proche), avoir besoin de \\
\hline
\zh{加强} & \emph{jiāqiáng} & v. & renforcer, intensifier \cle{★} \\
\hline
\zh{政府} & \emph{zhèngfǔ} & n. & gouvernement \cle{★} \\
\hline
\zh{保护} & \emph{bǎohù} & v. & protéger \cle{★} \\
\hline
\zh{帮助} & \emph{bāngzhù} & v. & aider, assistance \\
\hline
\zh{走上} & \emph{zǒushang} & v. & s'engager sur (une voie), emprunter (un chemin) \cle{★} \\
\hline
\zh{正确} & \emph{zhèngquè} & adj. & correct, juste, exact \cle{★} \\
\hline
\zh{道路} & \emph{dàolù} & n. & route, chemin, voie (litt. ou fig.) \\
\hline
\zh{只有} & \emph{zhǐyǒu} & conj. & seulement si (corrélé à \zh{才}) \cle{★} \\
\hline
\zh{这样} & \emph{zhèyàng} & pron. & ainsi, de cette manière \\
\hline
\zh{才} & \emph{cái} & adv. & seulement alors (marque la condition nécessaire) \cle{★} \\
\hline
\zh{能} & \emph{néng} & v. aux. & pouvoir, être capable de \\
\hline
\zh{减少} & \emph{jiǎnshǎo} & v. & réduire, diminuer, faire baisser \cle{★} \\
\hline
\end{tabularx}
\end{center}

\coursec{Clés (radicaux) et ordre des traits : focus sur 6 caractères nouveaux}

La maîtrise de l'écriture passe par la connaissance des \cle{clés (radicaux)} \zh{部首} (\emph{bùshǒu}) et le respect de l'\cle{ordre des traits} \zh{笔顺} (\emph{bǐshùn}). Règles de base : haut $\rightarrow$ bas, gauche $\rightarrow$ droite, horizontal avant vertical, trait central avant les ailes, cadre avant contenu, trait de fermeture en dernier.

\begin{definition}[Analyse de 6 caractères clés du thème]
\begin{itemize}
\item \zh{\textbf{犯}} (\emph{fàn}) — \textbf{Clé :} \zh{犭} (chien/quadrupède, 3 traits, « bête féroce ») — \textbf{Traits :} 5 — \textbf{Ordre :} \zh{犭} (point, courbe, point) + \zh{文} (point, horizontale, verticale, horizontale, horizontale). \textit{Sens originel :} bête qui attaque $\rightarrow$ offenser, commettre (un crime).
\item \zh{\textbf{罪}} (\emph{zuì}) — \textbf{Clé :} \zh{网} (filet, 5 traits, « filet ») — \textbf{Traits :} 13 — \textbf{Ordre :} \zh{网} (verticale, horizontale pliée, horizontale, horizontale, grande diagonale, 4 petits traits intérieurs) + \zh{非} (\emph{fēi}, « ne pas », 8 traits). \textit{Sens :} pris dans le filet $\rightarrow$ faute, crime.
\item \zh{\textbf{严}} (\emph{yán}) — \textbf{Clé :} \zh{厂} (rocher/falaise, 2 traits, « usine » à ne pas confondre) — \textbf{Traits :} 7 — \textbf{Ordre :} \zh{厂} (horizontale, verticale pliée) + \zh{工} (horizontale, verticale, horizontale) + \zh{口} (verticale, horizontale pliée, horizontale). \textit{Sens :} haut comme une falaise $\rightarrow$ strict, sévère.
\item \zh{\textbf{违}} (\emph{véi/wéi}) — \textbf{Clé :} \zh{辶} (marche, 3 traits, « avancer ») — \textbf{Traits :} 7 — \textbf{Ordre :} \zh{韦} (haut : \zh{十} + \zh{田} sans trait du bas) + \zh{辶} (point, horizontale pliée, grande diagonale). \textit{Sens :} aller dans le sens opposé $\rightarrow$ violer, enfreindre.
\item \zh{\textbf{吸}} (\emph{xī}) — \textbf{Clé :} \zh{口} (bouche, 3 traits) — \textbf{Traits :} 7 — \textbf{Ordre :} \zh{口} (verticale, horizontale pliée, horizontale) + \zh{及} (horizontale, verticale, horizontale, point). \textit{Sens :} bouche qui aspire $\rightarrow$ inhaler, absorber.
\item \zh{\textbf{护}} (\emph{hù}) — \textbf{Clé :} \zh{扌} (main, 3 traits, « action de la main ») — \textbf{Traits :} 7 — \textbf{Ordre :} \zh{扌} (horizontale, verticale, crochet) + \zh{户} (verticale, horizontale pliée, horizontale, point). \textit{Sens :} main qui protège la porte $\rightarrow$ protéger, garder.
\end{itemize}
\end{definition}

\begin{exercice}[Entraînement à l'écriture — Carnet de caractères]
\begin{enumerate}
\item Écrivez chaque caractère ci-dessus \textbf{10 fois} dans votre carnet en respectant l'ordre des traits et en prononçant la syllabe à chaque trait.
\item Pour chaque caractère, identifiez la clé, comptez le nombre de traits \textbf{hors clé} (utile pour la recherche dictionnaire) : \zh{犯}(2), \zh{罪}(8), \zh{严}(5), \zh{违}(4), \zh{吸}(4), \zh{护}(4).
\item \textbf{Dictée guidée :} Le professeur dicte : \zh{犯罪}, \zh{严重}, \zh{违法}, \zh{吸毒}, \zh{政府}, \zh{保护}. Les élèves écrivent les caractères + pinyin + traduction.
\end{enumerate}
\end{exercice}

\coursec{Compréhension du texte : questions graduées}

\begin{exercice}[Questions de compréhension — Réponses en français ou en chinois selon la consigne]
\textbf{A. Compréhension globale (choix multiple)}
\begin{enumerate}
\item Le texte est principalement : a) une histoire vécue \quad b) un texte argumentatif \quad c) un poème \quad d) une lettre administrative.
\item Le sujet central est : a) l'école \quad b) la famille \quad c) la délinquance juvénile \quad d) le gouvernement.
\item La conclusion du texte exprime : a) le découragement \quad b) la nécessité d'une action conjointe \quad c) la responsabilité unique des parents \quad d) l'inefficacité de la police.
\end{enumerate}

\textbf{B. Compréhension détaillée (réponses courtes)}
\begin{enumerate}
\item Citez les trois causes explicites mentionnées à la phrase 2 pour expliquer pourquoi « certains adolescents enfreignent la loi ».
\item À la phrase 3, trois comportements délinquants sont énumérés. Lesquels ? (Donnez les caractères et le sens).
\item Selon la phrase 4, que reproche le texte à l'école ? (Citez les deux verbes \zh{重视} et \zh{不重视} avec leurs objets).
\item Quels sont les trois acteurs sociaux appelés à « travailler ensemble » (phrase 5) ?
\item Quelle action spécifique est demandée au \zh{政府} (gouvernement) à la phrase 6 ?
\end{enumerate}

\textbf{C. Inférence et vocabulaire en contexte}
\begin{enumerate}
\item Expliquez la différence de sens entre \zh{问题} (\emph{wèntí}) à la phrase 1 et \zh{问题} à la phrase 2 (même mot, contextes différents).
\item Pourquoi l'auteur utilise-t-il \zh{而} (\emph{ér}) et non \zh{和} (\emph{hé}) dans \zh{因为...而违法} ?
\item Traduisez en chinois : « Ce n'est qu'ainsi qu'on pourra réduire la délinquance. » (Réutilisez la structure de la dernière phrase).
\end{enumerate}

\textbf{D. Réaction personnelle (expression orale/écrite)}
\begin{enumerate}
\item Selon vous, dans le contexte camerounais (Yaoundé, Douala, Buea...), quelle cause est la plus importante ? Justifiez en 2 phrases en chinois (vocabulaire du texte autorisé).
\item Proposez une mesure concrète que l'école pourrait prendre (\zh{学校应该...}).
\end{enumerate}
\end{exercice}

\begin{corrige}[Exercice « Compréhension du texte »]
\textbf{A.} 1. b) 2. c) 3. b)
\textbf{B.} 1. Problèmes familiaux (\zh{家庭问题}), pression des études (\zh{学习压力}), influence de mauvais amis (\zh{坏朋友的影响}). 2. \zh{偷东西} (voler), \zh{打架} (se battre), \zh{吸毒} (se droguer). 3. L'école \zh{只重视成绩} (ne valorise que les notes) et \zh{不重视教育} (ne valorise pas l'éducation morale/civique). 4. \zh{家庭} (famille), \zh{学校} (école), \zh{社会} (société). 5. \zh{保护青少年，帮助他们走上正确的道路} (protéger les adolescents, les aider à prendre le bon chemin).
\textbf{C.} 1. Phr. 1 : « problème de société » (phénomène global) ; Phr. 2 : « problèmes familiaux » (difficultés concrètes au foyer). 2. \zh{而} lie une cause (\zh{因为}) à sa conséquence directe de façon formelle ; \zh{和} coordonne des noms, pas des propositions. 3. \zh{只有这样，才能减少青少年犯罪。}
\textbf{D.} Réponses libres attendues. Exemple : \zh{我认为学习压力最严重。学生太累了。} / \zh{学校应该加强心理健康教育。}
\end{corrige}

\coursec{Production guidée et réinvestissement}

\courssub{Expression écrite : paragraphe argumentatif court}

\begin{methode}[Rédiger un micro-texte argumentatif (80-100 caractères) sur le modèle du texte support]
\begin{enumerate}
\item \textbf{Choisir un angle} : Causes familiales / Causes scolaires / Solutions éducatives / Rôle du gouvernement.
\item \textbf{Structurer} : 1. Phrase d'accroche (\zh{现在...是一个问题}) → 2. Cause(s) avec \zh{因为} / \zh{原因} → 3. Solution(s) avec \zh{为了} / \zh{应该} → 4. Conclusion avec \zh{只有...才...}.
\item \textbf{Utiliser le vocabulaire cible} : Minimum 6 mots nouveaux du glossaire (surlignez-les).
\item \textbf{Relire} : Vérifier l'ordre des mots (Sujet + Temps + Lieu + Manière + Verbe + Complément), les particules (\zh{的}, \zh{得}, \zh{地}), la ponctuation chinoise.
\end{enumerate}
\end{methode}

\begin{exemplebox}[Modèle de production d'élève (niveau HSK 3)]
\zh{现在，青少年吸毒是一个严重的问题。因为有些父母太忙，没有时间关心孩子，所以孩子交坏朋友。为了解决这个问题，父母要关心孩子，学校要加强教育。只有家庭和学校一起努力，才能减少青少年吸毒。}\\
\emph{Xiànzài, qīngshàonián xīdú shì yí ge yánzhòng de wèntí. Yīnwèi yǒuxiē fùmǔ tài máng, méiyǒu shíjiān guānxīn háizi, suǒyǐ háizi jiāo huài péngyou. Wèile jiějué zhège wèntí, fùmǔ yào guānxīn háizi, xuéxiào yào jiāqiáng jiàoyù. Zhǐyǒu jiātíng hé xuéxiào yìqǐ nǔlì, cái néng jiǎnshǎo qīngshàonián xīdú.}\\
« Aujourd'hui, la toxicomanie des jeunes est un grave problème. Parce que certains parents sont trop occupés pour s'occuper des enfants, ceux-ci fréquentent de mauvais amis. Pour résoudre ce problème, les parents doivent s'occuper des enfants, l'école doit renforcer l'éducation. Ce n'est que si famille et école travaillent ensemble qu'on pourra réduire la toxicomanie juvénile. »
\end{exemplebox}

\begin{exercice}[À vous d'écrire — Devoir / Travail en classe]
Rédigez un paragraphe de 6 à 8 phrases en caractères chinois (avec pinyin au brouillon) sur l'un des deux sujets :
\begin{enumerate}
\item \textbf{Sujet A :} \zh{谈谈“学习压力”与青少年犯罪} (Parlez de la « pression scolaire » et la délinquance juvénile). Utilisez : \zh{因为}, \zh{所以}, \zh{重视}, \zh{加强}, \zh{只有...才...}.
\item \textbf{Sujet B :} \zh{家庭、学校、社会：谁的责任最大？} (Famille, école, société : qui a la plus grande responsabilité ?). Utilisez : \zh{有的...有的...}, \zh{应该}, \zh{保护}, \zh{正确的道路}.
\end{enumerate}
\textbf{Critères d'évaluation (20 pts) :} Respect du plan (4) / Vocabulaire cible utilisé correctement (6) / Grammaire \& ordre des mots (6) / Caractères \& ponctuation (4).
\end{exercice}

\courssub{Expression orale : débat structuré}

\begin{experience}[Débat : « La délinquance juvénile : faute des parents ou de la société ? »]
\textbf{Organisation :} 2 groupes (Affirmatif / Négatif) + 1 jury (3 élèves). Durée : 15 min préparation + 10 min débat.
\begin{enumerate}
\item \textbf{Préparation (groupes)} : Chaque groupe prépare 3 arguments en chinois sur fiches. Vocabulaire imposé : \zh{因为}, \zh{所以}, \zh{但是}, \zh{例如}, \zh{关心}, \zh{教育}, \zh{影响}, \zh{责任} (\emph{zérèn}, responsabilité — mot nouveau à donner).
\item \textbf{Débat} : Tour de parole 1 min / groupe. Le jury note : prononciation (tons), structures correctes, richesse lexicale, réactivité.
\item \textbf{Synthèse} : Le jury donne son verdict en chinois simple : \zh{肯定组胜利，因为...} / \zh{否定组胜利，因为...}.
\end{enumerate}
\textbf{Prolongement} : Écrire un compte-rendu de 50 caractères : \zh{辩论结果：... 我的看法：...}
\end{experience}

\begin{savaistu}[Regards croisés : Prévention de la délinquance juvénile — Cameroun vs Chine]
\begin{itemize}
\item \textbf{Cameroun :} Âge de la majorité pénale : 18 ans (Code pénal). Mesures : Centres d'éducation spéciale (ex. Centre de Bétamba), tribunaux pour enfants, rôle des chefs traditionnels et associations (ex. RECAJ). Défis : surpopulation carcérale, manque d'éducateurs spécialisés, stigmatisation.
\item \textbf{Chine :} Âge de la responsabilité pénale : 16 ans (14 ans pour crimes graves — homicide, vol à main armée — depuis réforme 2021). Système : \zh{未成年人保护法} (Loi sur la protection des mineurs), \zh{预防未成年人犯罪法} (Loi sur la prévention de la délinquance juvénile). Mesures : \zh{训诫} (avertissement), \zh{责令具结悔过} (engagement de repentir), \zh{专门矫治教育} (éducation correctionnelle spécialisée dans écoles fermées), scellage des casiers judiciaires (\zh{犯罪记录封存}) à 18 ans pour favoriser réinsertion.
\item \textbf{Point commun :} Les deux pays privilégient de plus en plus l'\textbf{éducation et la réinsertion} (\zh{教育挽救}, \zh{矫治}) à la simple répression, et impliquent famille, école, communauté (\zh{社区}) et justice (\zh{检察院} / Parquet).
\end{itemize}
\textbf{Question débat :} L'abaissement de l'âge de responsabilité pénale (Chine : 14 ans pour crimes graves) est-il une solution efficace ou un risque d'injustice ? Argumentez en français.
\end{savaistu}

\courssub{Bilan de la leçon : auto-évaluation}

\begin{aretenir}[Check-list des compétences — Leçon 5 : Délinquance juvénile]
\begin{itemize}
\item[\square] Je sais lire le texte à voix haute avec les tons corrects (focus \zh{犯罪}, \zh{严重}, \zh{吸毒}).
\item[\square] Je reconnais, écris (ordre des traits) et traduis les 30 mots du glossaire (nature grammaticale identifiée).
\item[\square] Je sais segmenter une phrase chinoise sans espaces en mots de 1 à 3 caractères.
\item[\square] J'utilise correctement la structure \zh{有的...有的...还有的...} pour énumérer.
\item[\square] J'identifie et j'emploie les connecteurs logiques : \zh{因为}, \zh{为了}, \zh{只有...才...}, \zh{而}.
\item[\square] Je réponds aux questions de compréhension (globale, détaillée, inférence).
\item[\square] J'écris un paragraphe argumentatif structuré (problème $\rightarrow$ causes $\rightarrow$ solutions $\rightarrow$ conclusion).
\item[\square] Je participe à un débat oral en réinvestissant le vocabulaire du thème.
\item[\square] Je connais les radicaux et l'ordre des traits de \zh{犯}, \zh{罪}, \zh{严}, \zh{违}, \zh{吸}, \zh{护}.
\item[\square] Je peux comparer les approches camerounaise et chinoise de la justice des mineurs.
\end{itemize}
\end{aretenir}

\begin{exoresolu}[Modèle de rédaction guidée — Sujet A : « Pression scolaire et délinquance »]
Énoncé : Rédigez 6 phrases en chinois en suivant le plan : 1. Constat, 2. Cause (pression), 3. Conséquence (comportement), 4. Solution famille, 5. Solution école, 6. Conclusion conditionnelle.
\tcblower
\textbf{Proposition corrigée :}
\begin{enumerate}
\item \zh{现在，青少年犯罪是一个严重的社会问题。} (Constat)
\item \zh{因为学习压力太大，有些学生不想上学。} (Cause : \zh{因为} + \zh{压力} + \zh{太大})
\item \zh{他们有的逃学，有的离家出走，还有的违法犯罪。} (Conséquence : structure \zh{有的...有的...还有的...})
\item \zh{父母要关心孩子的心理健康，不只看成绩。} (Solution famille : \zh{关心}, \zh{心理健康} \emph{xīnlǐ jiànkāng} — santé mentale, \zh{不只} « pas seulement »)
\item \zh{学校应该减少作业，加强体育和艺术教育。} (Solution école : \zh{减少}, \zh{作业} \emph{zuòyè}, \zh{体育} \emph{tǐyù}, \zh{艺术} \emph{yìshù})
\item \zh{只有家庭和学校一起努力，才能减少青少年犯罪。} (Conclusion : \zh{只有...才...})
\end{enumerate}
\textbf{Vérification :} Tous les mots nouveaux du glossaire sont réinvestis (sauf \zh{政府}, \zh{保护}, \zh{吸毒}, \zh{走上}, \zh{正确}, \zh{道路} — à réinvestir dans Sujet B). Ponctuation chinoise respectée.
\end{exoresolu}

\begin{corrige}[Exercice « Entraînement à l'écriture — Carnet de caractères » (Q2)]
Nombre de traits hors clé (résiduel) :
\begin{itemize}
\item \zh{犯} : clé \zh{犭} (3 traits) → résiduel \zh{文} = 4 traits (souvent compté 4 en dict. moderne).
\item \zh{罪} : clé \zh{网} (5 traits) → résiduel \zh{非} = 8 traits.
\item \zh{严} : clé \zh{厂} (2 traits) → résiduel = 5 traits.
\item \zh{违} : clé \zh{辶} (3 traits) → résiduel \zh{韦} = 4 traits (simplifié).
\item \zh{吸} : clé \zh{口} (3 traits) → résiduel \zh{及} = 3 traits (mais \zh{及} est 3 traits, total 6? Vérif : \zh{吸} total 7, clé 3, résiduel 4). \textit{Correction :} \zh{及} = 3 traits (horiz, vert, horiz) + point final = 4 traits. Oui.
\item \zh{护} : clé \zh{扌} (3 traits) → résiduel \zh{户} = 4 traits.
\end{itemize}
\end{corrige}

\coursec{Mise en route : la délinquance juvénile, un défi de société}

\begin{experience}[Discussion déclencheur]
En classe entière, le professeur projette ou écrit au tableau les deux questions suivantes. Les élèves réfléchissent 2 minutes, puis échangent en binôme (3 min) avant une mise en commun orale (5 min).
\begin{enumerate}
\item \textbf{Au Cameroun comme en Chine,} quels comportements considère-t-on comme de la « délinquance juvénile » ? (Vol, bagarre, consommation de drogue, absentéisme scolaire, cyber-harcèlement…)
\item Selon vous, quelles sont les \emph{causes principales} qui poussent un adolescent à basculer dans l'illégalité ? (Pression scolaire, problèmes familiaux, influence des pairs, manque d'encadrement, pauvreté…)
\end{enumerate}
Le professeur note les idées-clés au tableau en français et en chinois (ex. : \zh{偷东西}, \zh{打架}, \zh{吸毒}, \zh{压力}, \zh{家庭问题}) pour introduire le lexique de la leçon.
\end{experience}

\coursec{Texte support : lecture guidée}

\begin{textebox}[La délinquance juvénile : causes et solutions]
青少年犯罪是一个严重的社会问题。在喀麦隆和中国，越来越多的青少年因为压力大、家庭缺乏关心、受坏人影响，开始违法犯罪：偷东西、打架、甚至吸毒。这些行为不仅伤害别人，也毁了自己的前途。\\
\emph{Qīngshàonián fànzuì shì yí ge yánzhòng de shèhuì wèntí. Zài Kāmàilóng hé Zhōngguó, yuè lái yuè duō de qīngshàonián yīnwèi yālì dà, jiātíng quēfá guānxīn, shòu huàirén yǐngxiǎng, kāishǐ wéifǎ fànzuì: tōu dōngxi, dǎjià, shènzhì xīdú. Zhèxiē xíngwéi bùjǐn shānghài biérén, yě huǐ le zìjǐ de qiántú.}\\
« La délinquance juvénile est un grave problème de société. Au Cameroun et en Chine, de plus en plus d'adolescents, à cause d'une forte pression, d'un manque d'attention familiale, sous l'influence de mauvaises fréquentations, commencent à enfreindre la loi et commettre des crimes : voler, se battre, voire se droguer. Ces actes blessent non seulement les autres, mais ruinent aussi leur propre avenir. »\\[6pt]
为了解决这个问题，社会、学校和家庭必须共同努力。政府要完善法律，加强教育；学校要关心学生，预防犯罪；父母要保护孩子，帮助他们走上正确的道路。对于已经犯错的青少年，不能只惩罚，更要教育和保护，让他们明白：违法不能解决问题，只有遵守法律，才能拥有美好的未来。\\
\emph{Wèile jiějué zhège wèntí, shèhuì, xuéxiào hé jiātíng bìxū gòngtóng nǔlì. Zhèngfǔ yào wánshàn fǎlǜ, jiāqiáng jiàoyù; xuéxiào yào guānxīn xuéshēng, yùfáng fànzuì; fùmǔ yào bǎohù háizi, bāngzhù tāmen zǒushàng zhèngquè de dàolù. Duìyú yǐjīng cuòwù de qīngshàonián, bùnéng zhǐ chéngfá, gèng yào jiàoyù hé bǎohù, ràng tāmen míngbai: wéifǎ bù néng jiějué wèntí, zhǐyǒu zūnshǒu fǎlǜ, cáinéng yǒngyǒu měihǎo de wèilái.}\\
« Pour résoudre ce problème, la société, l'école et la famille doivent unir leurs efforts. Le gouvernement doit perfectionner la législation, renforcer l'éducation ; l'école doit se soucier des élèves, prévenir la délinquance ; les parents doivent protéger leurs enfants, les aider à s'engager sur le bon chemin. Pour les adolescents qui ont déjà fauté, on ne peut pas seulement punir, il faut encore plus les éduquer et les protéger, pour qu'ils comprennent : enfreindre la loi ne résout pas les problèmes, seul le respect de la loi permet d'avoir un bel avenir. »
\end{textebox}

\begin{methode}[Lecture guidée en 3 étapes]
\begin{enumerate}
\item \textbf{Lecture globale (écoute + lecture silencieuse) :} Identifiez le sujet, le ton (alarmiste / constructif) et la structure (constat $\rightarrow$ causes $\rightarrow$ solutions).
\item \textbf{Lecture détaillée :} Surlignez les mots connus, entourez les mots nouveaux. Repérez les connecteurs logiques : \zh{因为…所以…} (cause-conséquence), \zh{不仅…也…} (progression), \zh{对于…} (changement de focus), \zh{只有…才…} (condition nécessaire).
\item \textbf{Lecture à voix haute :} Travail des tons (notamment \zh{法律 fǎlǜ}, \zh{原因 yuányīn}, \zh{保护 bǎohù}) et de la prosodie des phrases longues.
\end{enumerate}
\end{methode}

\coursec{Glossaire du thème}

\courssub{Tableau de vocabulaire}

Voici les vingt mots essentiels du thème « délinquance juvénile ». Pour chaque mot, observez les caractères, lisez le pinyin à voix haute en respectant les tons, vérifiez la nature grammaticale, puis mémorisez la traduction.

\begin{center}
\begin{tabularx}{\linewidth}{|X|X|X|X|}
\hline
\rowcolor{popblueL} Caractères & Pinyin & Nature & Traduction \\
\hline
青少年 & qīngshàonián & n. & adolescent, jeune \\
\hline
犯罪 & fànzuì & v./n. & commettre un crime ; crime, délit \\
\hline
罪 & zuì & n. & crime, faute, péché \\
\hline
严重 & yánzhòng & adj. & grave, sérieux \\
\hline
社会 & shèhuì & n. & société \\
\hline
问题 & wèntí & n. & problème ; question \\
\hline
违法 & wéifǎ & v. & violer la loi, enfreindre la loi (acte illégal) \\
\hline
法律 & fǎlǜ & n. & loi, législation \\
\hline
偷 & tōu & v. & voler, dérober \\
\hline
打架 & dǎjià & v. & se battre, se disputer (venez aux mains) \\
\hline
吸毒 & xīdú & v. & se droguer, consommer de la drogue \\
\hline
原因 & yuányīn & n. & cause, raison \\
\hline
压力 & yālì & n. & pression, stress \\
\hline
影响 & yǐngxiǎng & n./v. & influence ; influencer \\
\hline
关心 & guānxīn & v. & se soucier de, s'intéresser à, prendre soin de \\
\hline
教育 & jiàoyù & n./v. & éducation ; éduquer \\
\hline
保护 & bǎohù & v. & protéger \\
\hline
帮助 & bāngzhù & v./n. & aider ; aide \\
\hline
惩罚 & chéngfá & v./n. & punir ; punition, sanction \\
\hline
解决 & jiějué & v. & résoudre, régler (un problème) \\
\hline
\end{tabularx}
\end{center}

\begin{definition}[Nuance : \zh{违法} vs \zh{犯罪}]
\zh{违法} (\emph{wéifǎ}) = « enfreindre la loi » (sens large). Tout acte contraire à la loi, même mineur (excès de vitesse, stationnement interdit, vol à l'étalage).\\
\zh{犯罪} (\emph{fànzuì}) = « commettre un crime/délit » (sens pénal strict). Acte grave prévu et puni par le Code pénal (vol à main armée, trafic, homicide).\\
\emph{Règle :} Tout crime (\zh{犯罪}) est un acte illégal (\zh{违法}), mais tout acte illégal n'est pas un crime. En français : \emph{illégal} $\neq$ \emph{criminel}.
\end{definition}

\begin{exemplebox}[Le vocabulaire en phrases]
\zh{他因为偷东西受到了惩罚。}\\
\emph{Tā yīnwèi tōu dōngxi shòudào le chéngfá.}\\
« Il a été puni parce qu'il a volé des choses. »\\[4pt]
\zh{父母应该关心孩子的成长。}\\
\emph{Fùmǔ yīnggāi guānxīn háizi de chéngzhǎng.}\\
« Les parents doivent se soucier de la croissance de leurs enfants. »\\[4pt]
\zh{青少年犯罪是一个严重的社会问题。}\\
\emph{Qīngshàonián fànzuì shì yí ge yánzhòng de shèhuì wèntí.}\\
« La délinquance juvénile est un grave problème de société. »\\[4pt]
\zh{我们要遵守法律，不违法。}\\
\emph{Wǒmen yào zūnshǒu fǎlǜ, bù wéifǎ.}\\
« Nous devons respecter la loi, ne pas l'enfreindre. »
\end{exemplebox}

\courssub{Classer les mots par familles sémantiques}

\begin{methode}[Apprendre par familles de sens]
N'apprenez jamais une liste dans l'ordre alphabétique ou aléatoire. Procédez ainsi :
\begin{enumerate}
\item Lisez chaque mot et demandez : « De quoi parle-t-il ? » (Personne ? Action ? Cause ? Solution ? Cadre ?)
\item Regroupez logiquement (voir tableau ci-dessous).
\item Fabriquez une phrase \emph{personnelle} avec au moins un mot de chaque famille.
\item Le lendemain : cachez la colonne « Traduction », testez-vous famille par famille.
\end{enumerate}
\end{methode}

\begin{center}
\begin{tabularx}{\linewidth}{|X|X|X|}
\hline
\rowcolor{popblueL} Famille & Mots chinois + pinyin & Traduction \\
\hline
Personnes & 青少年 qīngshàonián & adolescent, jeune \\
\hline
Actes délictueux & 犯罪 fànzuì ; 违法 wéifǎ ; 偷 tōu ; 打架 dǎjià ; 吸毒 xīdú & crime ; illégalité ; voler ; se battre ; se droguer \\
\hline
Causes & 原因 yuányīn ; 压力 yālì ; 影响 yǐngxiǎng & cause ; pression ; influence \\
\hline
Solutions & 关心 guānxīn ; 教育 jiàoyù ; 保护 bǎohù ; 帮助 bāngzhù ; 惩罚 chéngfá ; 解决 jiějué & se soucier ; éduquer ; protéger ; aider ; punir ; résoudre \\
\hline
Cadre général & 社会 shèhuì ; 问题 wèntí ; 法律 fǎlǜ ; 严重 yánzhòng ; 罪 zuì & société ; problème ; loi ; grave ; crime \\
\hline
\end{tabularx}
\end{center}

\begin{popfigure}
\begin{tikzpicture}[mindmap, grow cyclic, every node/.style={concept, align=center, font=\small}, root concept/.append style={concept color=popblue, font=\bfseries\small}, level 1/.append style={level distance=3.2cm, sibling angle=90}, level 1 concept/.append style={concept color=poporangeL}]
\node[root concept, align=center]{青少年犯罪\\ \emph{délinquance}\\ \emph{juvénile}}
  child { node[align=center]{Personnes\\ 青少年\\ \emph{qīngshàonián}} }
  child { node[align=center]{Actes\\ 偷 打架 吸毒\\ \emph{tōu dǎjià xīdú}} }
  child { node[align=center]{Causes\\ 压力 影响\\ \emph{yālì yǐngxiǎng}} }
  child { node[align=center]{Solutions\\ 关心 教育 保护\\ \emph{guānxīn jiàoyù bǎohù}} };
\end{tikzpicture}
\legende{Carte mentale lexicale : organiser le vocabulaire autour du thème central.}
\end{popfigure}

\courssub{Entraînement à la prononciation : les paires de tons difficiles}

Trois mots du glossaire posent des difficultés de tons aux francophones. Répétez-les lentement, syllabe par syllabe, puis à vitesse normale.

\begin{itemize}
\item \zh{法律} \emph{fǎlǜ} (3$^{\text{ème}}$-4$^{\text{ème}}$ ton) : le 3$^{\text{ème}}$ ton descend puis remonte (ton creusé), le 4$^{\text{ème}}$ tombe net. Attention : \zh{律} se prononce \emph{lǜ} avec le son \emph{ü} (lèvres arrondies comme pour « u » français en prononçant « i »).
\item \zh{原因} \emph{yuányīn} (2$^{\text{ème}}$-1$^{\text{er}}$ ton) : le 2$^{\text{ème}}$ ton monte comme une question, le 1$^{\text{er}}$ reste haut et plat. Ne confondez pas avec \emph{yuànyìn} (inexistant).
\item \zh{保护} \emph{bǎohù} (3$^{\text{ème}}$-4$^{\text{ème}}$ ton) : même profil que \zh{法律}. \textbf{Règle de sandhi :} Devant un 4$^{\text{ème}}$ ton, le 3$^{\text{ème}}$ ton devient « demi-troisième » : il descend seulement, sans remonter. Dites \emph{bǎo} (bas) + \emph{hù} (tombe).
\end{itemize}

\begin{attention}[Le piège des tons plats]
Les francophones ont tendance à prononcer toutes les syllabes sur un ton moyen, comme en français. Or en chinois, un mauvais ton change le sens : \zh{tōu} (voler, 1$^{\text{er}}$ ton) $\neq$ \zh{tóu} (tête, 2$^{\text{ème}}$ ton) ; \zh{mǎi} (acheter, 3$^{\text{ème}}$) $\neq$ \zh{mài} (vendre, 4$^{\text{ème}}$). \textbf{Entraînez-vous toujours à voix haute, jamais mentalement.}
\end{attention}

\courssub{Collocations à mémoriser (搭配)}

En chinois, les mots s'associent de manière conventionnelle (collocations). Les maîtriser évite le « chinois traduit mot à mot ».

\begin{center}
\begin{tabularx}{\linewidth}{|X|X|X|}
\hline
\rowcolor{popblueL} Collocation & Pinyin & Traduction \\
\hline
预防犯罪 & yùfáng fànzuì & prévenir la délinquance / le crime \\
\hline
受到惩罚 & shòudào chéngfá & être puni, recevoir une sanction (passif) \\
\hline
解决问题 & jiějué wèntí & résoudre le problème \\
\hline
关心孩子 & guānxīn háizi & se soucier des enfants / prendre soin d'eux \\
\hline
走上正确的道路 & zǒushàng zhèngquè de dàolù & s'engager sur le bon chemin (réinsertion) \\
\hline
\end{tabularx}
\end{center}

\begin{propriete}[Structure passive avec \zh{受到}]
\zh{受到} (\emph{shòudào}, « recevoir, subir ») + [Nom d'action négative/neutre] = Passif formel (écrit, journalistique, administratif).
\zh{受到惩罚} (être puni), \zh{受到保护} (être protégé), \zh{受到教育} (être éduqué / recevoir une éducation), \zh{受到影响} (être influencé).
\emph{Ne pas dire :} \zh{被惩罚} (oral, plus direct) — les deux existent, \zh{受到} est plus formel.
\end{propriete}

\begin{exemplebox}[Collocations en contexte]
\zh{学校和家庭应帮助青少年走上正确的道路。}\\
\emph{Xuéxiào hé jiātíng yīng bāngzhù qīngshàonián zǒushàng zhèngquè de dàolù.}\\
« L'école et la famille doivent aider les jeunes à s'engager sur le bon chemin. »\\[4pt]
\zh{政府采取措施预防犯罪，解决社会问题。}\\
\emph{Zhèngfǔ cǎiqǔ cuòshī yùfáng fànzuì, jiějué shèhuì wèntí.}\\
« Le gouvernement prend des mesures pour prévenir la délinquance et résoudre les problèmes sociaux. »
\end{exemplebox}

\begin{attention}[\zh{问题} : « question » ou « problème » ?]
Le sens dépend du verbe collocatif :
\begin{itemize}
\item \zh{回答问题} (\emph{huídá wèntí}) = répondre à une \textbf{question} (examen, oral).
\item \zh{解决问题} (\emph{jiějué wèntí}) = résoudre un \textbf{problème} (difficulté, conflit).
\item \zh{没有问题} (\emph{méiyǒu wèntí}) = « pas de problème / pas de souci / OK ».
\end{itemize}
Regardez toujours le verbe qui précède ou suit !
\end{attention}

\coursec{Clés (radicaux) et ordre des traits}

\begin{definition}[Rappel : Clé (radical) et ordre des traits]
La \cle{clé} (部首, \emph{bùshǒu}) classe le caractère dans le dictionnaire et donne souvent un indice de sens. L'\cle{ordre des traits} (笔顺, \emph{bǐshùn}) suit des règles fixes : haut $\rightarrow$ bas, gauche $\rightarrow$ droite, horizontal avant vertical, trait central avant les ailes, cadre avant contenu, trait de fermeture en dernier. Respecter l'ordre assure l'équilibre, la vitesse et la reconnaissance par les OCR.
\end{definition}

Étudions 6 caractères clés du thème. Pour chacun : identifiez la clé (surlignée en \textbf{gras} dans le tableau**), comptez les traits, mémorisez l'ordre.

\begin{center}
\begin{tabularx}{\linewidth}{|c|X|X|X|X|}
\hline
\rowcolor{popblueL} Car. \& Clé (Radical) \& Nb traits \& Ordre des traits (description) \& Sens de la clé \\
\hline
\zh{法} \& \zh{氵} (trois points d'eau) \& 8 \& 1.点 2.点 3.提 (eau) ; 4.横 5.竖 6.横折 7.横 8.竖弯钩 (去) \& La loi « plane comme l'eau » : équité, justice. \\
\hline
\zh{罪} \& \zh{罒} (filet/réseau) \& 13 \& 1-5. Filet (罒) : 竖, 横折, 横, 横, 竖 ; 6-13. \zh{非} (non) en bas : 横, 竖, 提, 横, 竖, 横折, 横, 横 \& Être pris dans le filet = faute, crime. \\
\hline
\zh{压} \& \zh{厂} (rocher/falaise) \& 5 \& 1.横 2.竖 3.横 (厂) ; 4.竖弯钩 5.横 (厶) \& Un rocher qui écrase = pression, peser. \\
\hline
\zh{护} \& \zh{扌} (main) \& 7 \& 1.横 2.竖钩 3.提 (main) ; 4.横 5.竖 6.横折 7.横 (戶) \& La main qui protège une porte/maison. \\
\hline
\zh{犯} \& \zh{犭} (chien/quadrupède) \& 5 \& 1.撇 2.弯钩 3.提 (chien) ; 4.横 5.竖弯钩 (然 simplifié) \& Le chien qui mord/attaque = offenser, violer. \\
\hline
\zh{解} \& \zh{角} (corne) \& 13 \& 1-7. \zh{角} (corne) : 横, 竖, 横折, 横, 竖, 横折, 横 ; 8-13. \zh{刀} (couteau) + \zh{牛} (bœuf) : 刀 (撇, 竖) + 牛 (横, 竖, 横折, 横) \& Séparer les cornes avec un couteau = délier, résoudre, comprendre. \\
\hline
\end{tabularx}
\end{center}

\begin{experience}[Atelier écriture : « Course aux traits »]
Par groupes de 4. Le professeur dicte un caractère du tableau (ex. \zh{护}). Le 1$^{\text{er}}$ élève écrit le 1$^{\text{er}}$ trait au tableau, passe la craie au 2$^{\text{e}}$ pour le 2$^{\text{e}}$ trait, etc. Le groupe qui finit le premier \textbf{avec l'ordre exact} gagne. Répéter pour \zh{法}, \zh{罪}, \zh{解}.
\end{experience}

\coursec{Compréhension du texte support}

Répondez en chinois (phrase complète si possible) ou en français aux questions ci-dessous, en vous basant \textbf{uniquement} sur le texte support (section 2).

\begin{exercice}[Questions de compréhension]
\begin{enumerate}
\item 根据课文，青少年犯罪的三个主要原因是什么？(D'après le texte, quelles sont les 3 causes principales ?)
\item 课文提到了哪三类违法犯罪行为？(Quels trois types d'actes illégaux sont cités ?)
\item 谁应该共同努力来解决这个问题？(Qui doit unir ses efforts pour résoudre le problème ?)
\item 对于已经犯错的青少年，只能惩罚吗？还应该怎么做？(Pour les jeunes fautifs, la punition suffit-elle ? Que faut-il faire d'autre ?)
\item 翻译这句话：\zh{只有遵守法律，才能拥有美好的未来。}
\end{enumerate}
\end{exercice}

\begin{corrige}[Compréhension du texte support]
\begin{enumerate}
\item \zh{压力大、家庭缺乏关心、受坏人影响。} (Forte pression, manque d'attention familiale, influence de mauvaises fréquentations.)
\item \zh{偷东西、打架、吸毒。} (Voler, se battre, se droguer.)
\item \zh{社会、学校和家庭。} (La société, l'école et la famille.)
\item \zh{不能只惩罚，更要教育和保护。} (On ne peut pas seulement punir, il faut surtout éduquer et protéger.)
\item « Ce n'est qu'en respectant la loi qu'on peut avoir un bel avenir. » (Structure \zh{只有…才…} = condition nécessaire).
\end{enumerate}
\end{corrige}

\coursec{Production guidée et réinvestissement}

\courssub{Mini-texte de réinvestissement (Contexte camerounais)}

Lisons un court texte réemployant le vocabulaire dans une situation locale.

\begin{textebox}[雅温得的一个帮助中心]
在雅温得，有一个帮助青少年的中心。一些年轻人因为压力大、家庭问题多，开始偷东西或者打架。中心的工作人员关心他们，听他们说话，也教他们学习法律。\\
\emph{Zài Yǎwēndé, yǒu yí ge bāngzhù qīngshàonián de zhōngxīn. Yìxiē niánqīngrén yīnwèi yālì dà, jiātíng wèntí duō, kāishǐ tōu dōngxi huòzhě dǎjià. Zhōngxīn de gōngzuò rényuán guānxīn tāmen, tīng tāmen shuōhuà, yě jiāo tāmen xuéxí fǎlǜ.}\\
« À Yaoundé, il existe un centre d'aide aux adolescents. Certains jeunes, à cause d'une forte pression et de nombreux problèmes familiaux, se mettent à voler ou à se battre. Les employés du centre se soucient d'eux, les écoutent et leur apprennent aussi la loi. »\\[6pt]
慢慢地，这些孩子明白了：违法不能解决问题。他们受到了教育和保护，走上了正确的道路。\\
\emph{Mànmàn de, zhèxiē háizi míngbai le: wéifǎ bù néng jiějué wèntí. Tāmen shòudào le jiàoyù hé bǎohù, zǒushàng le zhèngquè de dàolù.}\\
« Petit à petit, ces enfants ont compris : enfreindre la loi ne résout pas les problèmes. Ils ont reçu éducation et protection, et se sont engagés sur le bon chemin. »
\end{textebox}

\begin{exercice}[Questions rapides sur le mini-texte]
\begin{enumerate}
\item 中心在什么地方？(Où se trouve le centre ?)
\item 年轻人为什么偷东西或打架？(Pourquoi les jeunes volent/se battent-ils ?)
\item 工作人员教他们什么？(Qu'est-ce que le personnel leur enseigne ?)
\item 最后这些孩子怎么样了？(Que sont devenus ces enfants à la fin ?)
\end{enumerate}
\end{exercice}

\courssub{Exercice résolu : Complétion lexicale}

\begin{exoresolu}[Choisir le bon mot : \zh{惩罚}，\zh{压力}，\zh{解决}，\zh{违法}]
Complétez les phrases. Attention à la nature grammaticale (verbe, nom, adjectif).
\begin{enumerate}
\item 偷东西是 \blank 的行为。 (Voler est un acte \ldots)
\item 他学习 \blank 很大，所以心情不好。 (Sa pression d'études est forte, il est de mauvaise humeur.)
\item 警察惩罚了那个犯罪的人，他受到了 \blank{}。 (La police a puni le criminel, il a \textbf{reçu} \ldots)
\item 我们一起想办法 \blank 这个问题。 (Cherchons ensemble à \ldots ce problème.)
\end{enumerate}
\tcblower
\textbf{Solutions et justifications :}
\begin{enumerate}
\item \zh{违法} (\emph{wéifǎ}). Contexte : qualification d'un acte. \zh{违法} (adj./v.) = illégal. \zh{犯罪} serait trop fort (crime pénal).
\item \zh{压力} (\emph{yālì}). Structure : \zh{压力} + \zh{大} (adjectif « grand/fort ») = « forte pression ». \zh{压力大} est un syntagme figé.
\item \zh{惩罚} (\emph{chéngfá}). Collocation passive apprise : \zh{受到惩罚} (recevoir une punition). Ici \zh{惩罚} fonctionne comme nom.
\item \zh{解决} (\emph{jiějué}). Collocation : \zh{解决问题} (résoudre le problème). Verbe transitif direct.
\end{enumerate}
\end{exoresolu}

\courssub{Exercices d'entraînement (À vous de jouer)}

\begin{exercice}[Production écrite et orale]
\begin{enumerate}
\item \textbf{Traduction (Thème) :} « Les parents doivent protéger leurs enfants contre les mauvaises influences. » (Indice : \zh{不受…影响} ou \zh{不让…受…影响}).
\item \textbf{Classement :} Rangez ces 5 mots en deux colonnes \textbf{Causes} / \textbf{Solutions} : \zh{压力}，\zh{教育}，\zh{影响}，\zh{关心}，\zh{帮助}。
\item \textbf{Phonétique :} Lisez à voix haute 3 fois devant un camarade : \zh{法律}，\zh{原因}，\zh{保护}. Cochez la case quand les tons sont stables : \fbox{\phantom{ok}} \zh{法律} \quad \fbox{\phantom{ok}} \zh{原因} \quad \fbox{\phantom{ok}} \zh{保护}.
\item \textbf{Rédaction courte (40-60 caractères) :} Écrivez 3 phrases en chinois pour décrire une action concrète que l'école ou la famille peut faire pour \zh{预防犯罪}. Utilisez au moins 3 mots du glossaire.
\end{enumerate}
\end{exercice}

\begin{corrige}[Exercice « À vous de jouer »]
\begin{enumerate}
\item \zh{父母应该保护孩子，不让他们受坏的影响。} (\emph{Fùmǔ yīnggāi bǎohù háizi, bù ràng tāmen shòu huài de yǐngxiǎng.})\\
  Variante acceptée : \zh{父母应该保护孩子不受坏的影响。}
\item \textbf{Causes :} \zh{压力}，\zh{影响}. \quad \textbf{Solutions :} \zh{教育}，\zh{关心}，\zh{帮助}.
\item Auto-évaluation orale : \emph{fǎ-lǜ} (3-4, demi-3$^{\text{ème}}$), \emph{yuán-yīn} (2-1), \emph{bǎo-hù} (3-4, demi-3$^{\text{ème}}$). Le camarade valide avec le professeur.
\item \textbf{Exemple de production :}\\
  \zh{学校应该加强法律教育，关心学生的压力，帮助他们解决问题，预防犯罪。}\\
  \emph{Xuéxiào yīnggāi jiāqiáng fǎlǜ jiàoyù, guānxīn xuéshēng de yālì, bāngzhù tāmen jiějué wèntí, yùfáng fànzuì.}\\
  « L'école doit renforcer l'éducation légale, se soucier de la pression des élèves, les aider à résoudre leurs problèmes, prévenir la délinquance. » (Mots utilisés : \zh{教育}, \zh{压力}, \zh{帮助}, \zh{解决}, \zh{预防犯罪}).
\end{enumerate}
\end{corrige}

\begin{savaistu}[Étymologie : \zh{法} , la loi « plane comme l'eau »]
Le caractère \zh{法} (\emph{fǎ}) combine \zh{氵} (eau) et \zh{去} (partir/ôter). Dans l'Antiquité, \zh{法} évoquait un animal mythique (\zh{獬豸} \emph{xièzhì}) capable de distinguer le bien du mal d'un coup de corne. L'eau, par nature plane et horizontale, symbolise l'équité : la loi doit être « plane comme la surface de l'eau » (\zh{水平如一}), égale pour tous, sans favoritisme. Retenir cette image (eau + justice) ancre durablement le sens de \zh{法} et \zh{法律}.
\end{savaistu}

\begin{aretenir}[L'essentiel de la leçon 5 — Vocabulaire \& Compréhension]
\begin{itemize}
\item \textbf{20 mots clés} classés en 5 familles : Personnes (\zh{青少年}), Actes (\zh{偷/打架/吸毒/违法/犯罪}), Causes (\zh{压力/影响/原因}), Solutions (\zh{关心/教育/保护/帮助/惩罚/解决}), Cadre (\zh{社会/法律/问题/严重/罪}).
\item \textbf{Distinction cruciale :} \zh{违法} (illégal, général) $\neq$ \zh{犯罪} (crime, pénal, grave).
\item \textbf{Polysémie de \zh{问题} :} \zh{回答问题} = question ; \zh{解决问题} = problème.
\item \textbf{5 collocations figées :} \zh{预防犯罪}, \zh{受到惩罚}, \zh{解决问题}, \zh{关心孩子}, \zh{走上正确的道路}.
\item \textbf{Tons à maîtriser :} \emph{fǎ-lǜ} (3-4), \emph{yuán-yīn} (2-1), \emph{bǎo-hù} (3-4 avec sandhi).
\item \textbf{6 caractères à écrire :} \zh{法} (氵, 8 tr.), \zh{罪} (罒, 13 tr.), \zh{压} (厂, 5 tr.), \zh{护} (扌, 7 tr.), \zh{犯} (犭, 5 tr.), \zh{解} (角, 13 tr.).
\item \textbf{Structure texte :} Constat (gravité) $\rightarrow$ Causes ($\zh{因为}$) $\rightarrow$ Solutions ($\zh{必须}$, \zh{对于}... \zh{不能只…更要…}) $\rightarrow$ Message final ($\zh{只有…才…}$).
\end{itemize}
\end{aretenir}

\courssub{Devoirs / Travail personnel}
\begin{enumerate}
\item Recopier 2 fois chaque caractère du tableau « Clés et ordre des traits » en respectant l'ordre (utiliser une feuille à carreaux chinois \emph{田字格}).
\item Apprendre par cœur les 5 collocations + les 20 mots du glossaire (test écrit la prochaine séance).
\item Rédiger un court paragraphe (60-80 caractères) : « Comment aider un ami qui a des ennuis ? » en utilisant \zh{关心}, \zh{帮助}, \zh{解决}, \zh{走上正确的道路}.
\item (Optionnel) Regarder une vidéo courte (2-3 min) en chinois sur \zh{青少年法治教育} (éducation légale des jeunes) sur Bilibili ou YouTube, noter 3 mots reconnus.
\end{enumerate}

\coursec{Clés (radicaux) et ordre des traits}

Dans la section précédente, nous avons dressé le glossaire du thème de la délinquance juvénile : \zh{犯罪} (fànzuì, commettre un crime), \zh{法律} (fǎlǜ, la loi), \zh{保护} (bǎohù, protéger), \zh{教育} (jiàoyù, éducation), \zh{惩罚} (chéngfá, punir). Vous savez désormais lire et prononcer ces mots. Mais pour les \emph{écrire} correctement — et surtout pour les \emph{mémoriser} durablement —, il faut comprendre comment les caractères sont construits. C'est l'objet de cette section : observer la structure interne des caractères, identifier leur \cle{clé} (radical) et appliquer les règles d'\cle{ordre des traits}.

\courssub{Observer avant d'écrire : pourquoi les radicaux ?}

Un caractère chinois n'est presque jamais un dessin arbitraire. La plupart des caractères sont des \emph{caractères phono-idéographiques} (形声字) composés de deux éléments : une partie qui indique le \emph{champ de sens} (la clé, ou radical) et une partie qui indique souvent la \emph{prononciation} (la phonétique). Observer d'abord, écrire ensuite : telle est la bonne démarche.

\begin{propriete}[Le radical donne le champ sémantique]
La clé (radical) d'un caractère renseigne très souvent sur sa famille de sens :
\begin{itemize}
\item \zh{犭} (quǎn, variante de \zh{犬}, chien) : famille des \textbf{animaux} — \zh{犯} (enfreindre, à l'origine : l'animal qui attaque), \zh{狗} (gǒu, chien), \zh{猫} (māo, chat).
\item \zh{扌} (shǒu, variante de \zh{手}, main) : famille des \textbf{actions de la main} — \zh{护} (hù, protéger), \zh{打} (dǎ, frapper), \zh{拉} (lā, tirer).
\item \zh{氵} (shuǐ, variante de \zh{水}, eau) : famille de l'\textbf{eau et des liquides} — \zh{法} (fǎ, la loi ; sens ancien : l'eau qui coule à niveau, symbole d'équité), \zh{河} (hé, rivière), \zh{洗} (xǐ, laver).
\item \zh{罒} (wǎng, variante de \zh{网}, filet) : famille du \textbf{filet}, par extension ce qui \textbf{prend au piège} — \zh{罪} (zuì, crime : celui qui est pris dans les mailles de la justice), \zh{罚} (fá, punir), \zh{罗} (luó, filet, attraper).
\end{itemize}
\end{propriete}

\begin{methode}[Deviner le sens d'un caractère inconnu]
Face à un caractère que vous ne connaissez pas, procédez en trois étapes :
\begin{enumerate}
\item \textbf{Repérer le radical} : il se trouve le plus souvent à gauche (\zh{扌}, \zh{氵}, \zh{犭}) ou en haut (\zh{罒}, \zh{艹}).
\item \textbf{Dégager le champ de sens} : main ? eau ? filet ? Cela donne une hypothèse de signification.
\item \textbf{Observer la partie restante} : c'est souvent la phonétique, qui suggère la prononciation. Exemple : dans \zh{护} (hù), la partie \zh{户} se prononce aussi \emph{hù} — c'est elle qui donne le son.
\end{enumerate}
Vérifiez ensuite votre hypothèse dans le dictionnaire ou le glossaire.
\end{methode}

\courssub{Les règles générales de l'ordre des traits}

Avant d'étudier nos cinq caractères, rappelons les quatre règles d'or de la calligraphie chinoise, valables pour tous les caractères :

\begin{aretenir}[Les quatre règles d'or du tracé]
\begin{enumerate}
\item \textbf{De haut en bas} (从上到下 cóng shàng dào xià) : on trace d'abord la partie supérieure. Exemple : \zh{三} (sān, trois), trois traits du haut vers le bas.
\item \textbf{De gauche à droite} (从左到右 cóng zuǒ dào yòu) : on trace d'abord la partie gauche. Exemple : \zh{法} : d'abord \zh{氵}, ensuite \zh{去}.
\item \textbf{L'extérieur avant l'intérieur} (先外后内 xiān wài hòu nèi) : pour les caractères encadrés, on trace le cadre, puis le contenu, et on « ferme la porte » en dernier. Exemple : \zh{国} (guó, pays).
\item \textbf{L'horizontal avant le vertical} (先横后竖 xiān héng hòu shù) : quand deux traits se croisent, l'horizontal passe d'abord. Exemple : \zh{十} (shí, dix) : d'abord \zh{一}, ensuite \zh{丨}.
\end{enumerate}
\end{aretenir}

\begin{attention}[Erreur fréquente des francophones]
En français, nous écrivons « comme ça vient », en reliant les lettres. En chinois, \textbf{chaque trait est séparé et suit un ordre fixe}. Ne « dessinez » pas les caractères : un tracé dans le désordre produit des caractères illisibles et empêche de compter les traits, ce qui rend impossible la recherche dans un dictionnaire classé par nombre de traits.
\end{attention}

\courssub{Étude détaillée des cinq caractères du thème}

\textbf{1. 罪 (zuì) — crime, faute}\par
Structure : \zh{罒} (wǎng, filet) au-dessus + \zh{非} (fēi, « n'est pas », le mal) en dessous. Image : celui qui fait le mal est pris dans le filet de la justice. Total : \textbf{13 traits}. On applique la règle « de haut en bas » : on trace d'abord le filet \zh{罒} (5 traits : vertical, horizontal-crochet, deux verticaux intérieurs, horizontal de fermeture), puis \zh{非} (8 traits : trois horizontaux + vertical à gauche, trois horizontaux + vertical à droite).

\begin{popfigure}
\begin{tikzpicture}[scale=1.0]
\node[draw=popblue, thick, rounded corners, fill=popblueL, minimum width=2.2cm, minimum height=1.1cm, align=center] (top) at (0,2.2) {\zh{罒}\\ \footnotesize filet (5 traits)};
\node[draw=poporange, thick, rounded corners, fill=poporangeL, minimum width=2.2cm, minimum height=1.1cm, align=center] (bottom) at (0,0) {\zh{非}\\ \footnotesize le mal (8 traits)};
\node[draw=popgreen, very thick, rounded corners, fill=popgreenL, minimum width=1.4cm, minimum height=1.4cm, align=center] (result) at (4.2,1.1) {\Huge \zh{罪}};
\draw[-{Stealth}, very thick, popblue] (top.east) -- node[above, font=\footnotesize, align=center]{1. d'abord\\ le haut} (result.north west);
\draw[-{Stealth}, very thick, poporange] (bottom.east) -- node[below, font=\footnotesize, align=center]{2. ensuite\\ le bas} (result.south west);
\node[font=\footnotesize, align=center, text=popdark] at (4.2,-1.1) {de haut en bas\\ 13 traits au total};
\end{tikzpicture}
\legende{Décomposition de 罪 : le filet 罒 au-dessus, 非 en dessous ; on trace de haut en bas.}
\end{popfigure}

\textbf{2. 法 (fǎ) — loi, droit}\par
Structure : \zh{氵} (shuǐ, eau, 3 traits à gauche) + \zh{去} (qù, aller) à droite. Total : \textbf{8 traits}. Règle « de gauche à droite » : on trace d'abord les trois gouttes d'eau \zh{氵} (point, point, trait relevé), puis \zh{去} (5 traits : horizontal, vertical, horizontal, trait plié, point). L'eau évoque l'équité : la loi doit être plane comme la surface de l'eau.

\textbf{3. 犯 (fàn) — enfreindre, commettre}\par
Structure : \zh{犭} (quǎn, chien/animal, 3 traits à gauche) + \zh{㔾} (jié) à droite. Total : \textbf{5 traits}. Ordre : d'abord les traits du radical \zh{犭} (trait oblique, crochet courbe, trait oblique), ensuite la partie droite \zh{㔾} (2 traits). Sens ancien : l'animal qui bondit, qui « franchit » la limite — d'où « transgresser ». C'est le caractère de \zh{犯罪} (fànzuì, commettre un crime) et de \zh{犯人} (fànrén, délinquant).

\textbf{4. 护 (hù) — protéger}\par
Structure : \zh{扌} (shǒu, main, 3 traits à gauche) + \zh{户} (hù, porte, famille) à droite. Total : \textbf{7 traits}. La main \zh{扌} évoque l'action concrète de protéger ; \zh{户} donne le son \emph{hù}. Image : la main placée devant la porte de la maison pour la défendre. On le trouve dans \zh{保护} (bǎohù, protéger) : \zh{法律保护青少年。} (Fǎlǜ bǎohù qīngshàonián.) « La loi protège les adolescents. »

\textbf{5. 教 (jiào) — enseigner, éduquer}\par
Structure : \zh{孝} (xiào, piété filiale, l'enfant) à gauche + \zh{攵} (pū, taper doucement avec un bâton) à droite. Total : \textbf{11 traits}. Le radical est \zh{攵} : le sens ancien est « guider l'enfant, le corriger pour qu'il apprenne ». Aujourd'hui, \zh{教} signifie simplement « enseigner » : \zh{教育} (jiàoyù, éducation), \zh{教师} (jiàoshī, enseignant). Exemple : \zh{老师教育学生。} (Lǎoshī jiàoyù xuésheng.) « Le professeur éduque les élèves. »

\begin{exemplebox}[Les cinq caractères en contexte]
\zh{法律保护青少年。} \emph{Fǎlǜ bǎohù qīngshàonián.} « La loi protège les adolescents. »\\
\zh{老师教育学生。} \emph{Lǎoshī jiàoyù xuésheng.} « Le professeur éduque les élèves. »\\
\zh{他犯了罪。} \emph{Tā fàn le zuì.} « Il a commis un crime. »\\
\zh{学校和家庭应该保护孩子。} \emph{Xuéxiào hé jiātíng yīnggāi bǎohù háizi.} « L'école et la famille doivent protéger les enfants. »
\end{exemplebox}

\courssub{Tableau récapitulatif}

\begin{center}
\begin{tabularx}{\linewidth}{|X|X|X|X|}
\hline
\rowcolor{popblueL} Caractère (pinyin) & Clé (radical) & Traits & Structure et image \\
\hline
\zh{罪} zuì, crime & \zh{罒} filet & 13 & haut + bas : pris dans le filet de la justice \\
\hline
\zh{法} fǎ, loi & \zh{氵} eau & 8 & gauche + droite : l'eau plane, symbole d'équité \\
\hline
\zh{犯} fàn, enfreindre & \zh{犭} animal & 5 & gauche + droite : franchir la limite \\
\hline
\zh{护} hù, protéger & \zh{扌} main & 7 & gauche + droite : la main devant la porte \\
\hline
\zh{教} jiào, enseigner & \zh{攵} guider & 11 & gauche + droite : guider l'enfant \\
\hline
\end{tabularx}
\end{center}

\courssub{Attention aux caractères proches}

\begin{attention}[罚 ou 罪 ?]
Ne confondez pas \zh{罚} (fá, punir — clé \zh{罒} + \zh{讠} parole + \zh{刂} couteau) et \zh{罪} (zuì, crime — clé \zh{罒} + \zh{非}). Ils partagent le filet \zh{罒}, mais :
\begin{itemize}
\item \zh{犯罪} (fànzuì) = \textbf{commettre un crime} : \zh{青少年不应该犯罪。} (Qīngshàonián bù yīnggāi fànzuì.) « Les adolescents ne devraient pas commettre de crimes. »
\item \zh{惩罚} (chéngfá) = \textbf{punir} : \zh{法律惩罚犯人。} (Fǎlǜ chéngfá fànrén.) « La loi punit les délinquants. »
\end{itemize}
Astuce : dans \zh{罚}, il y a le couteau \zh{刂} à droite — la punition « tranche » ; dans \zh{罪}, il y a \zh{非} (le mal) — le crime est le mal commis.
\end{attention}

\courssub{Entraînement à l'écriture}

\begin{exoresolu}[Identifier le radical et compter les traits]
Énoncé : pour chaque caractère, indiquez le radical et le nombre total de traits : \zh{犯}, \zh{护}, \zh{法}.
\tcblower
Solution détaillée :
\begin{itemize}
\item \zh{犯} : radical \zh{犭} (animal) à gauche ; 3 traits + 2 traits = \textbf{5 traits}.
\item \zh{护} : radical \zh{扌} (main) à gauche ; 3 traits + 4 traits (\zh{户}) = \textbf{7 traits}.
\item \zh{法} : radical \zh{氵} (eau) à gauche ; 3 traits + 5 traits (\zh{去}) = \textbf{8 traits}.
\end{itemize}
Méthode : repérer d'abord la partie gauche (radical fréquent), la tracer mentalement, puis compter la partie droite.
\end{exoresolu}

\begin{exercice}[À votre tour : tracer les caractères]
Sur votre cahier, tracez chaque caractère \textbf{5 fois} en respectant l'ordre des traits et en prononçant le pinyin à voix haute à chaque tracé : \zh{罪} (13 traits, de haut en bas), \zh{法} (8 traits, de gauche à droite), \zh{犯} (5 traits), \zh{护} (7 traits), \zh{教} (11 traits). Vérifiez ensuite que vous savez écrire les mots \zh{犯罪}, \zh{法律}, \zh{保护} et \zh{教育} sans modèle.
\end{exercice}

\begin{corrige}[Exercice « À votre tour »]
Critères de réussite : (1) pour \zh{罪}, le filet \zh{罒} est tracé avant \zh{非} ; (2) pour \zh{法}, \zh{犯} et \zh{护}, le radical de gauche est tracé en premier ; (3) pour \zh{教}, la partie gauche \zh{孝} précède \zh{攵} ; (4) le nombre de traits est respecté : 13, 8, 5, 7, 11. Auto-correction : comparez votre tracé au tableau récapitulatif et refaites tout caractère dont l'ordre était faux.
\end{corrige}

\begin{savaistu}[Le radical, clé du dictionnaire]
Les dictionnaires chinois traditionnels classent les caractères par leurs 214 radicaux historiques (système du dictionnaire Kangxi). Pour chercher un caractère inconnu, on identifie son radical, puis on compte les traits \emph{restants}. Savoir décomposer \zh{罪} en \zh{罒} + \zh{非}, c'est donc déjà savoir utiliser un dictionnaire chinois — une compétence précieuse pour la suite de vos études, au Cameroun comme en Chine.
\end{savaistu}

\coursec{Compréhension du texte}

Après avoir décodé les caractères du thème — leurs clés, leurs traits, leur prononciation — nous sommes maintenant prêts à vérifier que nous \emph{comprenons} réellement le texte sur la délinquance juvénile. Savoir lire les caractères ne suffit pas : il faut pouvoir répondre à des questions, retrouver une information précise, dire si une phrase est vraie ou fausse, et reformuler avec ses propres mots. C'est exactement ce que l'on vous demandera aux évaluations de fin d'année. Cette section vous donne une méthode rigoureuse, des réponses modèles et des exercices résolus.

\courssub{Le texte de référence (rappel)}

Avant de répondre aux questions, relisons le texte support étudié dans cette leçon. Toutes les réponses se trouvent \emph{dans} le texte : un bon élève ne répond jamais « de tête », il retourne chercher la phrase exacte.

\begin{textebox}[青少年犯罪 — Qīngshàonián fànzuì]
现在，青少年犯罪是一个严重的问题。\\
Xiànzài, qīngshàonián fànzuì shì yí ge yánzhòng de wèntí.\\
« Aujourd'hui, la délinquance juvénile est un problème grave. »\\[4pt]
青少年犯罪的原因有很多：有的家庭不关心孩子，有的学校只重视成绩，不重视教育，还有一些年轻人交了坏朋友。\\
Qīngshàonián fànzuì de yuányīn yǒu hěn duō : yǒu de jiātíng bù guānxīn háizi, yǒu de xuéxiào zhǐ zhòngshì chéngjì, bù zhòngshì jiàoyù, hái yǒu yìxiē niánqīngrén jiāo le huài péngyou.\\
« Les causes de la délinquance juvénile sont nombreuses : certaines familles ne s'occupent pas de leurs enfants, certaines écoles ne valorisent que les notes et non l'éducation, et certains jeunes se font de mauvais amis. »\\[4pt]
他们有的偷东西，有的打架，还有的喝酒、吸毒。\\
Tāmen yǒu de tōu dōngxi, yǒu de dǎjià, hái yǒu de hē jiǔ, xīdú.\\
« Certains volent des objets, d'autres se battent, d'autres encore boivent de l'alcool et se droguent. »\\[4pt]
谁应该帮助这些青少年呢？家庭、学校和社会应该一起帮助他们。\\
Shéi yīnggāi bāngzhù zhèxiē qīngshàonián ne ? Jiātíng, xuéxiào hé shèhuì yīnggāi yìqǐ bāngzhù tāmen.\\
« Qui doit aider ces adolescents ? La famille, l'école et la société doivent les aider ensemble. »\\[4pt]
父母应该多和孩子说话，了解他们的想法。\\
Fùmǔ yīnggāi duō hé háizi shuōhuà, liǎojiě tāmen de xiǎngfǎ.\\
« Les parents doivent parler davantage avec leurs enfants et comprendre leurs idées. »\\[4pt]
只有家庭、学校和社会一起努力，才能减少青少年犯罪。\\
Zhǐyǒu jiātíng, xuéxiào hé shèhuì yìqǐ nǔlì, cái néng jiǎnshǎo qīngshàonián fànzuì.\\
« Ce n'est que si la famille, l'école et la société travaillent ensemble que l'on pourra réduire la délinquance juvénile. »
\end{textebox}

\courssub{Méthode : 问什么，答什么}

En chinois comme en français, la première règle d'une bonne réponse de compréhension est de \cle{reprendre les mots de la question} dans la réponse. Les Chinois résument cela par une formule à retenir : \zh{问什么，答什么} (Wèn shénme, dá shénme) — « Ce qu'on demande, on le répond ».

\begin{methode}[问什么，答什么 — « Ce qu'on demande, on le répond »]
\begin{enumerate}
\item \textbf{Repérer le mot interrogatif} : \zh{谁} (shéi, qui), \zh{什么} (shénme, quoi), \zh{怎么} (zěnme, comment), \zh{为什么} (wèishénme, pourquoi). Il indique le type d'information attendu.
\item \textbf{Retrouver dans le texte la phrase clé} qui contient la réponse : cherchez les mots de la question (\zh{原因}, \zh{帮助}, \zh{父母}...) dans le texte.
\item \textbf{Recopier le squelette de la question} et remplacer le mot interrogatif par l'information trouvée : \zh{谁应该帮助青少年？} → \zh{家庭、学校和社会应该帮助青少年。}
\item \textbf{Vérifier} : la réponse est une phrase complète, avec sujet + verbe, et la ponctuation chinoise 。 à la fin.
\end{enumerate}
\end{methode}

\begin{exemplebox}[Question → réponse modèle]
Q : \zh{谁应该帮助青少年？} (Shéi yīnggāi bāngzhù qīngshàonián ?) — « Qui doit aider les adolescents ? »\\
R : \zh{家庭、学校和社会应该帮助青少年。} (Jiātíng, xuéxiào hé shèhuì yīnggāi bāngzhù qīngshàonián.) — « La famille, l'école et la société doivent aider les adolescents. »\\[4pt]
Observez : la réponse reprend \emph{tous} les mots de la question (\zh{应该帮助青少年}) ; seul le mot interrogatif \zh{谁} est remplacé par la réponse \zh{家庭、学校和社会}. C'est la technique du « squelette ».
\end{exemplebox}

\begin{popfigure}
\begin{tikzpicture}[node distance=0.55cm, every node/.style={font=\small}]
\node[draw=popblue, fill=popblueL, rounded corners, align=center, minimum width=4.2cm] (q) {1. Question\\\zh{谁应该帮助青少年？}};
\node[draw=poporange, fill=poporangeL, rounded corners, align=center, minimum width=4.2cm, right=1.1cm of q] (rep) {2. Repérage de la phrase clé\\\zh{家庭、学校和社会}\\\zh{应该一起帮助他们。}};
\node[draw=poppurple, fill=poppurpleL, rounded corners, align=center, minimum width=4.2cm, below=0.9cm of q] (ref) {3. Reformulation\\remplacer \zh{谁}\\par l'information trouvée};
\node[draw=popgreen, fill=popgreenL, rounded corners, align=center, minimum width=4.2cm, right=1.1cm of ref] (ans) {4. Réponse complète\\\zh{家庭、学校和社会}\\\zh{应该帮助青少年。}};
\draw[-{Stealth}, thick, popdark] (q) -- (rep);
\draw[-{Stealth}, thick, popdark] (rep) |- (ref);
\draw[-{Stealth}, thick, popdark] (ref) -- (ans);
\end{tikzpicture}
\legende{Organigramme de réponse : de la question à la réponse complète en quatre étapes.}
\end{popfigure}

\courssub{Questions de compréhension et réponses attendues}

Voici les cinq questions de la leçon, avec pour chacune la phrase clé du texte et la réponse complète attendue. Apprenez à produire ces réponses \emph{par écrit}, avec les caractères.

\begin{center}
\begin{tabularx}{\linewidth}{|X|X|X|}
\hline
\rowcolor{popblueL} Question & Phrase clé du texte & Réponse complète attendue \\
\hline
\zh{青少年犯罪的原因是什么？} & \zh{原因有很多：家庭不关心孩子…} & \zh{青少年犯罪的原因有很多：有的家庭不关心孩子，有的学校只重视成绩，还有一些年轻人交了坏朋友。} \\
\hline
\zh{青少年犯什么罪？} & \zh{有的偷东西，有的打架…} & \zh{他们有的偷东西，有的打架，还有的喝酒、吸毒。} \\
\hline
\zh{谁应该帮助青少年？} & \zh{家庭、学校和社会应该一起帮助他们。} & \zh{家庭、学校和社会应该帮助青少年。} \\
\hline
\zh{父母应该怎么做？} & \zh{父母应该多和孩子说话…} & \zh{父母应该多和孩子说话，了解他们的想法。} \\
\hline
\zh{怎么才能减少青少年犯罪？} & \zh{只有…一起努力，才能减少…} & \zh{只有家庭、学校和社会一起努力，才能减少青少年犯罪。} \\
\hline
\end{tabularx}
\end{center}

Traductions des questions et des réponses :
\begin{enumerate}
\item \zh{青少年犯罪的原因是什么？} (Qīngshàonián fànzuì de yuányīn shì shénme ?) — « Quelles sont les causes de la délinquance juvénile ? » Réponse : « Les causes sont nombreuses : certaines familles ne s'occupent pas des enfants, certaines écoles ne valorisent que les notes, et certains jeunes se font de mauvais amis. »
\item \zh{青少年犯什么罪？} (Qīngshàonián fàn shénme zuì ?) — « Quels actes commettent les adolescents ? » Réponse : « Certains volent, d'autres se battent, d'autres encore boivent et se droguent. » Remarquez : \zh{犯罪} (fànzuì) est un verbe séparable « commettre + crime » ; à la question, on peut donc dire \zh{犯什么罪} « commettre quels actes ».
\item \zh{谁应该帮助青少年？} (Shéi yīnggāi bāngzhù qīngshàonián ?) — « Qui doit aider les adolescents ? » Réponse : « La famille, l'école et la société. »
\item \zh{父母应该怎么做？} (Fùmǔ yīnggāi zěnme zuò ?) — « Que doivent faire les parents ? » Réponse : « Les parents doivent parler davantage avec leurs enfants et comprendre leurs idées. »
\item \zh{怎么才能减少青少年犯罪？} (Zěnme cái néng jiǎnshǎo qīngshàonián fànzuì ?) — « Comment réduire la délinquance juvénile ? » Réponse : « Ce n'est qu'en travaillant ensemble — famille, école, société — qu'on peut la réduire. »
\end{enumerate}

Remarquez la structure \zh{有的...有的...还有的...} (yǒu de... yǒu de... hái yǒu de...) : « certains... d'autres... d'autres encore... ». Elle est idéale pour énumérer des cas dans une réponse de compréhension ou une production écrite.

\courssub{Point de grammaire : 只有...才...}

La dernière phrase du texte contient une structure essentielle du niveau HSK 3, très souvent exigée dans les réponses de type « comment ? ».

\begin{propriete}[La structure 只有...才... (zhǐyǒu...cái...)]
\textbf{Forme} : 只有 + condition，才 + résultat.\\
\textbf{Sens} : « seulement si... alors... », « ce n'est que... que... ». Elle exprime une \cle{condition nécessaire} : sans la condition, le résultat est impossible.\\[4pt]
Exemple du texte : \zh{只有家庭、学校和社会一起努力，才能减少青少年犯罪。} (Zhǐyǒu jiātíng, xuéxiào hé shèhuì yìqǐ nǔlì, cái néng jiǎnshǎo qīngshàonián fànzuì.) — « Ce n'est que si la famille, l'école et la société travaillent ensemble que l'on peut réduire la délinquance juvénile. »\\[4pt]
Autres exemples :
\begin{itemize}
\item \zh{只有好好学习，才能考上大学。} (Zhǐyǒu hǎohao xuéxí, cái néng kǎo shàng dàxué.) — « Ce n'est qu'en travaillant sérieusement qu'on peut être reçu au concours d'entrée à l'université. »
\item \zh{只有多练习，才能写好汉字。} (Zhǐyǒu duō liànxí, cái néng xiě hǎo Hànzì.) — « Ce n'est qu'en s'entraînant souvent qu'on peut bien écrire les caractères chinois. »
\end{itemize}
\textbf{Comparaison avec le français} : le français dit « seulement si... » ou « il faut... pour... » ; le chinois encadre la phrase avec deux mots obligatoires : \zh{只有} devant la condition, \zh{才} devant le résultat. On ne peut pas en omettre un. Attention à ne pas confondre avec \zh{只要...就...} (zhǐyào...jiù..., « du moment que... alors... »), qui exprime une condition \emph{suffisante} : \zh{只要努力，就会进步。} (Zhǐyào nǔlì, jiù huì jìnbù.) — « Du moment qu'on fait des efforts, on progresse. »
\end{propriete}

\begin{attention}[Piège : 才 ne veut pas dire « seulement » ici]
Les francophones connaissent \zh{才} dans le sens de « seulement » (quantité) : \zh{他才十五岁。} (Tā cái shíwǔ suì.) — « Il n'a que quinze ans. » Mais dans la structure \zh{只有...才...}, \zh{才} ne signifie \emph{pas} « seulement » : il marque que le résultat devient possible \emph{grâce à} la condition. Ne traduisez donc jamais mot à mot \zh{只有...才...} par « seulement... seulement... ». Traduction correcte : « ce n'est que... que... » ou « il faut... pour... ».
\end{attention}

\courssub{Vrai ou faux : vérifier une affirmation}

Le vrai/faux est un exercice de précision : une seule différence de mot peut rendre une phrase fausse. Il faut citer le texte pour justifier. En chinois, on répond \zh{对} (duì, vrai) ou \zh{错} (cuò, faux).

\begin{exoresolu}[Vrai ou faux ?]
Énoncé : \zh{学校只重视教育，不重视成绩。} (Xuéxiào zhǐ zhòngshì jiàoyù, bù zhòngshì chéngjì.) — « L'école ne valorise que l'éducation, pas les notes. » Vrai ou faux ? Justifiez avec le texte.
\tcblower
\textbf{Solution} : FAUX (\zh{错} / cuò). Le texte dit exactement l'inverse : \zh{有的学校只重视成绩，不重视教育。} (Yǒu de xuéxiào zhǐ zhòngshì chéngjì, bù zhòngshì jiàoyù.) — « Certaines écoles ne valorisent que les notes, pas l'éducation. » L'affirmation a inversé \zh{成绩} (les notes) et \zh{教育} (l'éducation). Méthode : comparer mot à mot l'affirmation et la phrase du texte ; ici, l'ordre des deux noms est interverti, donc le sens est contraire.
\end{exoresolu}

\courssub{Recherche dans le texte et reformulation}

\textbf{Recherche.} Consigne : retrouvez la phrase qui exprime la \emph{solution principale} du problème. La solution principale n'est pas une action isolée (parler aux enfants), mais l'idée de la \emph{coopération}. C'est la phrase : \zh{家庭、学校和社会应该一起努力。} (Jiātíng, xuéxiào hé shèhuì yīnggāi yìqǐ nǔlì.) — « La famille, l'école et la société doivent travailler ensemble. » Le mot-clé qui signale la solution est \zh{一起} (yìqǐ, ensemble).

\textbf{Reformulation.} Expliquer « avec ses propres mots » signifie : garder le sens, changer les mots et la structure. Voici trois reformulations possibles de la conclusion du texte, en chinois simple :
\begin{itemize}
\item \zh{大家一起努力，青少年犯罪就会减少。} (Dàjiā yìqǐ nǔlì, qīngshàonián fànzuì jiù huì jiǎnshǎo.) — « Si tout le monde travaille ensemble, la délinquance juvénile diminuera. »
\item \zh{要减少青少年犯罪，家庭、学校和社会都要帮忙。} (Yào jiǎnshǎo qīngshàonián fànzuì, jiātíng, xuéxiào hé shèhuì dōu yào bāngmáng.) — « Pour réduire la délinquance, la famille, l'école et la société doivent tous aider. »
\item \zh{一个人不能解决这个问题，大家要一起努力。} (Yí ge rén bù néng jiějué zhège wèntí, dàjiā yào yìqǐ nǔlì.) — « Une personne seule ne peut pas résoudre ce problème, tout le monde doit coopérer. »
\end{itemize}

\begin{exercice}[À vous de jouer]
\begin{enumerate}
\item Répondez en phrases complètes : \zh{青少年犯罪的原因是什么？} / \zh{父母应该怎么做？}
\item Vrai ou faux (justifiez) : \zh{只有学校应该帮助青少年。}
\item Reformulez avec vos mots : \zh{只有一起努力，才能减少犯罪。}
\item Traduisez en chinois : « Ce n'est qu'en parlant avec les enfants que les parents peuvent les comprendre. »
\end{enumerate}
\end{exercice}

\begin{corrige}[Exercice « À vous de jouer »]
1. \zh{青少年犯罪的原因有很多：有的家庭不关心孩子，有的学校只重视成绩，还有一些年轻人交了坏朋友。} / \zh{父母应该多和孩子说话，了解他们的想法。}\\
2. FAUX (\zh{错}). Le texte dit \zh{家庭、学校和社会应该一起帮助他们} : ce n'est pas seulement l'école, mais les trois ensemble.\\
3. Exemple : \zh{大家一起努力，犯罪就会减少。} (Dàjiā yìqǐ nǔlì, fànzuì jiù huì jiǎnshǎo.)\\
4. \zh{只有多和孩子说话，父母才能了解他们。} (Zhǐyǒu duō hé háizi shuōhuà, fùmǔ cái néng liǎojiě tāmen.)
\end{corrige}

\courssub{Discussion guidée : et au Cameroun ?}

\begin{experience}[Débat : la délinquance juvénile au Cameroun]
Par groupes de quatre, discutez en chinois simple (phrases courtes avec \zh{应该}, \zh{只有...才...}, \zh{有的...有的...}) :
\begin{itemize}
\item Au Cameroun, certains enfants ne vont plus à l'école : \zh{有的孩子不上学。} (Yǒu de háizi bù shàngxué.) — « Certains enfants ne vont pas à l'école. » Pourquoi ?
\item Dans les grandes villes comme Douala et Yaoundé, il y a des enfants de la rue. Qui peut les aider ? \zh{家庭、政府和社会组织应该帮助他们。} (Jiātíng, zhèngfǔ hé shèhuì zǔzhī yīnggāi bāngzhù tāmen.) — « Les familles, le gouvernement et les organisations sociales (ONG) doivent les aider. »
\item Conclusion du débat, à formuler avec la structure du jour : \zh{只有大家一起努力，才能帮助这些孩子。} (Zhǐyǒu dàjiā yìqǐ nǔlì, cái néng bāngzhù zhèxiē háizi.) — « Ce n'est qu'en travaillant tous ensemble qu'on peut aider ces enfants. »
\end{itemize}
\end{experience}

\begin{savaistu}[Cameroun et Chine : deux réponses, une même idée]
En Chine comme au Cameroun, la prévention de la délinquance juvénile repose sur le trio \zh{家庭、学校、社会} (famille, école, société). En Chine, la loi sur la protection des mineurs insiste sur la coresponsabilité des parents et de l'école ; au Cameroun, de nombreuses ONG et associations de quartier jouent un rôle central auprès des enfants déscolarisés ou des enfants de la rue. Dans les deux pays, la conclusion est la même : aucun acteur ne peut réussir seul — \zh{只有一起努力，才能成功。} (Zhǐyǒu yìqǐ nǔlì, cái néng chénggōng.) — « Ce n'est qu'en travaillant ensemble qu'on peut réussir. »
\end{savaistu}

\begin{aretenir}[L'essentiel de la section]
\begin{itemize}
\item Méthode de réponse : \zh{问什么，答什么} — reprendre les mots de la question et remplacer le mot interrogatif par l'information du texte.
\item Structure de condition nécessaire : \zh{只有 + condition，才 + résultat} = « ce n'est que... que... » ; \zh{才} n'est pas « seulement » ici ; ne pas confondre avec \zh{只要...就...} (condition suffisante).
\item Vrai/faux : toujours justifier en citant la phrase exacte du texte ; répondre \zh{对} (vrai) ou \zh{错} (faux).
\item Reformuler = garder le sens, changer les mots ; la solution principale du texte : \zh{家庭、学校和社会应该一起努力。}
\end{itemize}
\end{aretenir}

\coursec{Production guidée et réinvestissement}

Après avoir lu, compris et questionné le texte sur la délinquance juvénile, il est temps de passer à l'étape décisive de toute leçon de langue : \cle{produire} à votre tour. Comprendre un texte, c'est bien ; savoir réemployer son vocabulaire et ses structures dans une production personnelle, écrite puis orale, c'est le vrai objectif du programme de Première. Dans cette section, vous allez d'abord rédiger un court texte guidé de cinq à six phrases, puis mener une interview orale en binôme, corriger collectivement une production au tableau, et enfin découvrir le format de la mini-dissertation type baccalauréat.

\courssub{Observer avant de rédiger : un modèle de production}

Avant d'écrire, observons comment un élève peut organiser ses idées. Lisez ce modèle de production rédigé par une élève de Première A de Yaoundé. Repérez l'ordre des idées : problème, causes, solutions, conclusion.

\begin{textebox}[我们城市的青少年犯罪问题 — Modèle de production]
在我们城市，有些青少年犯罪，这是一个严重的问题。\\
Zài wǒmen chéngshì, yǒuxiē qīngshàonián fànzuì, zhè shì yí ge yánzhòng de wèntí.\\
« Dans notre ville, certains adolescents commettent des délits, c'est un problème grave. »\\[4pt]
犯罪的原因有很多：第一，有些家庭不关心孩子；第二，一些青少年不上学，也没有工作。\\
Fànzuì de yuányīn yǒu hěn duō : dì-yī, yǒuxiē jiātíng bù guānxīn háizi; dì-èr, yìxiē qīngshàonián bú shàngxué, yě méiyǒu gōngzuò.\\
« Les causes de la délinquance sont nombreuses : premièrement, certaines familles ne s'occupent pas de leurs enfants ; deuxièmement, certains adolescents ne vont pas à l'école et n'ont pas de travail. »\\[4pt]
为了解决这个问题，政府应该给青少年更多的工作机会，学校也应该教育学生遵守法律。\\
Wèile jiějué zhège wèntí, zhèngfǔ yīnggāi gěi qīngshàonián gèng duō de gōngzuò jīhuì, xuéxiào yě yīnggāi jiàoyù xuésheng zūnshǒu fǎlǜ.\\
« Pour résoudre ce problème, le gouvernement devrait donner plus d'opportunités de travail aux jeunes, et l'école devrait aussi éduquer les élèves à respecter la loi. »\\[4pt]
只有家庭、学校和社会一起努力，才能帮助青少年走上正确的道路。\\
Zhǐyǒu jiātíng, xuéxiào hé shèhuì yìqǐ nǔlì, cái néng bāngzhù qīngshàonián zǒushàng zhèngquè de dàolù.\\
« Ce n'est que si la famille, l'école et la société travaillent ensemble qu'on pourra aider les adolescents à suivre le droit chemin. »
\end{textebox}

Observez la logique du texte : chaque phrase a une fonction précise. La première présente le \cle{problème}, la deuxième donne deux \cle{causes} (marquées par \zh{第一} et \zh{第二}), la troisième propose deux \cle{solutions} (introduites par \zh{为了}), la dernière conclut avec la structure \zh{只有……才能……}. C'est exactement la trame que vous allez suivre.

\courssub{La trame de production : problème, causes, solutions, conclusion}

\begin{definition}[La trame argumentative en quatre étapes]
Un court texte argumentatif chinois de niveau Première suit un plan fixe : \cle{问题} (le problème) $\rightarrow$ \cle{原因} (les causes, au moins deux) $\rightarrow$ \cle{办法} (les solutions, au moins deux) $\rightarrow$ \cle{结论} (la conclusion). Chaque étape correspond à une ou deux phrases.
\end{definition}

\begin{popfigure}
\begin{tikzpicture}[
  node distance=0.55cm and 1.1cm,
  box/.style={draw=popblue, fill=popblueL, rounded corners=3pt, align=center, minimum width=2.6cm, minimum height=1cm, font=\small},
  arr/.style={-{Stealth[length=3mm]}, thick, poporange}
]
\node[box, align=center] (p){问题\\ \emph{le problème}};
\node[box, right=of p, align=center] (c){原因1 + 原因2\\ \emph{deux causes}};
\node[box, right=of c, align=center] (s){办法1 + 办法2\\ \emph{deux solutions}};
\node[box, right=of s, align=center] (k){结论\\ \emph{conclusion}};
\draw[arr] (p) -- (c);
\draw[arr] (c) -- (s);
\draw[arr] (s) -- (k);
\end{tikzpicture}
\legende{La trame de production en quatre étapes : du problème à la conclusion.}
\end{popfigure}

\courssub{Banque de phrases-modèles}

Pour chaque étape de la trame, voici des phrases-modèles à trous. Vous les recopiez et vous les complétez avec le glossaire de la leçon. Ce sont des \cle{structures passe-partout} : elles fonctionnent pour n'importe quel sujet de société (la délinquance, mais aussi la drogue, le chômage, la violence).

\begin{center}
\begin{tabularx}{\linewidth}{|X|X|X|}
\hline
\rowcolor{popblueL} Étape & Phrase-modèle (caractères + pinyin) & Traduction \\
\hline
Problème & 在我们城市，有些青少年…… Zài wǒmen chéngshì, yǒuxiē qīngshàonián... & Dans notre ville, certains adolescents... \\
\hline
Causes & 犯罪的原因有…… Fànzuì de yuányīn yǒu... & Les causes de la délinquance sont... \\
\hline
Solutions & 为了解决这个问题，……应该…… Wèile jiějué zhège wèntí, ... yīnggāi... & Pour résoudre ce problème, ... devrait... \\
\hline
Conclusion & 只有这样，才能…… Zhǐyǒu zhèyàng, cái néng... & C'est seulement ainsi qu'on pourra... \\
\hline
\end{tabularx}
\end{center}

\begin{propriete}[La structure 只有……才能……]
Structure : \cle{只有} + condition, \cle{才能} + résultat. Elle signifie « ce n'est que si... que... ». Le mot \zh{才} (cái) souligne que le résultat n'est possible \emph{qu'à cette condition}. Exemple : \zh{只有努力学习，才能通过考试。} (Zhǐyǒu nǔlì xuéxí, cái néng tōngguò kǎoshì.) « Ce n'est qu'en travaillant dur qu'on peut réussir l'examen. » Attention : en français on dit « seulement si », en chinois \zh{只有} se place \emph{avant} la condition et \zh{才} \emph{avant le verbe} du résultat.
\end{propriete}

\begin{propriete}[La structure 为了……，……应该……]
\cle{为了} + but, sujet + \cle{应该} + verbe : « pour + but, ... devrait... ». \zh{为了} introduit le but et se place en tête de phrase ; \zh{应该} exprime le devoir, la recommandation. Exemple : \zh{为了保护青少年，家长应该多和孩子交流。} (Wèile bǎohù qīngshàonián, jiāzhǎng yīnggāi duō hé háizi jiāoliú.) « Pour protéger les adolescents, les parents devraient communiquer davantage avec leurs enfants. »
\end{propriete}

\begin{attention}[La cause avant la conséquence]
En chinois, la \cle{cause précède toujours la conséquence} dans la phrase : \zh{因为家庭问题，有些青少年违法。} (Yīnwèi jiātíng wèntí, yǒuxiē qīngshàonián wéifǎ.) « À cause de problèmes familiaux, certains adolescents enfreignent la loi. » Le français permet l'ordre inverse (« Certains adolescents enfreignent la loi parce que... »), mais en chinois on ne dit jamais \zh{有些青少年违法，因为家庭问题}. De même, avec \zh{所以} : \zh{因为他们没有工作，所以他们犯罪。} (Yīnwèi tāmen méiyǒu gōngzuò, suǒyǐ tāmen fànzuì.) « Parce qu'ils n'ont pas de travail, ils délinquent. » Gardez toujours l'ordre : cause $\rightarrow$ conséquence.
\end{attention}

\begin{methode}[Rédiger en trois temps : plan, traduction, relecture]
\begin{enumerate}
\item \textbf{Plan en français.} Écrivez d'abord votre plan en français, en quatre lignes : le problème, deux causes, deux solutions, une conclusion. Ne cherchez pas encore le chinois.
\item \textbf{Traduction phrase par phrase.} Traduisez chaque ligne en chinois en vous appuyant sur le glossaire de la leçon et sur les phrases-modèles. Une phrase = une idée. Restez simple : Sujet + Verbe + Complément.
\item \textbf{Relecture.} Vérifiez trois choses : les \cle{tons} du pinyin (un mauvais ton change le sens), l'\cle{ordre des mots} (la cause avant la conséquence, le complément de temps ou de lieu avant le verbe), et les \cle{caractères} (traits, clés, homophones comme \zh{犯} et \zh{饭}).
\end{enumerate}
\end{methode}

\begin{exoresolu}[Rédiger les deux premières phrases de la production]
Énoncé : en suivant la trame, rédigez la phrase du problème et la phrase des causes sur le sujet « 我们城市的青少年犯罪问题 ». Utilisez les mots \zh{有些}, \zh{犯罪}, \zh{原因}, \zh{家庭}, \zh{关心}.
\tcblower
Solution détaillée :\\
Étape 1 — Plan en français : « Problème : dans notre ville, certains adolescents commettent des délits. Causes : certaines familles ne s'occupent pas de leurs enfants, et certains jeunes ne vont pas à l'école. »\\
Étape 2 — Traduction avec les modèles :\\
\zh{在我们城市，有些青少年犯罪。} (Zài wǒmen chéngshì, yǒuxiē qīngshàonián fànzuì.) « Dans notre ville, certains adolescents commettent des délits. »\\
\zh{犯罪的原因有很多：有些家庭不关心孩子，有些青少年不上学。} (Fànzuì de yuányīn yǒu hěn duō : yǒuxiē jiātíng bù guānxīn háizi, yǒuxiē qīngshàonián bú shàngxué.) « Les causes sont nombreuses : certaines familles ne s'occupent pas de leurs enfants, certains adolescents ne vont pas à l'école. »\\
Étape 3 — Relecture : \zh{犯罪} (fànzuì) porte bien le 4\textsuperscript{e} ton puis le 4\textsuperscript{e} ton ; \zh{不} devant \zh{关心} (1\textsuperscript{er} ton) reste \zh{bù}, mais devant \zh{上学} (shàng, 4\textsuperscript{e} ton) il devient \zh{bú} (règle du \emph{sandhi} tonal) ; la cause (\zh{家庭不关心孩子}) est bien placée dans la phrase des causes, après le deux-points.
\end{exoresolu}

\begin{exercice}[À vous d'écrire]
Rédigez votre production complète « 我们城市的青少年犯罪问题 » en cinq à six phrases : une phrase de problème, une phrase de deux causes, une phrase de deux solutions avec \zh{为了……应该……}, et une conclusion avec \zh{只有……才能……}. Employez au moins cinq mots du glossaire : \zh{青少年}, \zh{犯罪}, \zh{原因}, \zh{解决}, \zh{法律}, \zh{社会}, \zh{努力}.
\end{exercice}

\courssub{Expression orale en binômes : l'interview}

\begin{experience}[Journaliste et expert : interview sur la délinquance juvénile au Cameroun]
Par groupes de deux : l'élève A joue le \cle{journaliste} d'une chaîne de télévision de Douala, l'élève B joue un \cle{expert} (éducateur, sociologue ou policier). Le journaliste pose au moins quatre questions ; l'expert répond en phrases complètes.\\
Règle du jeu : l'expert doit employer \textbf{au moins cinq mots du glossaire} de la leçon, par exemple \zh{青少年}, \zh{犯罪}, \zh{原因}, \zh{家庭}, \zh{法律}, \zh{解决}, \zh{社会}.\\
Questions possibles pour le journaliste :\\
\zh{为什么有些青少年犯罪？} (Wèishénme yǒuxiē qīngshàonián fànzuì ?) « Pourquoi certains adolescents délinquent-ils ? »\\
\zh{家庭有什么问题？} (Jiātíng yǒu shénme wèntí ?) « Quels problèmes y a-t-il dans les familles ? »\\
\zh{政府应该怎么做？} (Zhèngfǔ yīnggāi zěnme zuò ?) « Que devrait faire le gouvernement ? »\\
\zh{学校能帮助青少年吗？} (Xuéxiào néng bāngzhù qīngshàonián ma ?) « L'école peut-elle aider les adolescents ? »\\
Ensuite, échangez les rôles. Le professeur écoute et note les erreurs de tons pour la correction collective.
\end{experience}

\begin{corrige}[Exercice « À vous d'écrire » — éléments de production attendus]
Une production réussie contient : (1) la phrase d'introduction du problème avec \zh{在我们城市，有些青少年……} ; (2) deux causes reliées par \zh{第一……第二……} ou par un deux-points ; (3) deux solutions introduites par \zh{为了解决这个问题，……应该……} ; (4) une conclusion en \zh{只有……才能……}. Exemple de conclusion attendue : \zh{只有家庭、学校和社会一起努力，才能帮助青少年走上正确的道路。} Vérifiez : cause avant conséquence, \zh{应该} devant le verbe, tons de \zh{犯罪} (fànzuì) et de \zh{解决} (jiějué) correctement marqués.
\end{corrige}

\courssub{Correction collective au tableau}

Le professeur recopie au tableau une production d'élève (anonyme) contenant des erreurs typiques. La classe cherche et corrige ensemble. Voici les trois familles d'erreurs à repérer, avec des exemples fautifs et leur correction.

\begin{center}
\begin{tabularx}{\linewidth}{|X|X|X|}
\hline
\rowcolor{popblueL} Type d'erreur & Phrase fautive & Correction (caractères + pinyin + traduction) \\
\hline
Ordre des mots (conséquence avant cause) & 有些青少年违法，因为家庭问题。 & 因为家庭问题，有些青少年违法。 Yīnwèi jiātíng wèntí, yǒuxiē qīngshàonián wéifǎ. « À cause de problèmes familiaux, certains adolescents enfreignent la loi. » \\
\hline
Erreur de ton & \zh{解决} écrit \emph{jiéjué} & \zh{解决} se lit \emph{jiějué} (3\textsuperscript{e} + 2\textsuperscript{e} ton) : « résoudre ». Un mauvais ton peut créer un mot inexistant ou un autre mot. \\
\hline
Caractère confondu & \zh{饭罪} (fànzuì avec \zh{饭} « riz ») & \zh{犯罪} : \zh{犯} (clé du chien \zh{犭}) « commettre », non \zh{饭} (clé de la nourriture \zh{饣}) « riz, repas ». \\
\hline
Place de \zh{应该} & 政府给应该青少年工作。 & 政府应该给青少年工作。 Zhèngfǔ yīnggāi gěi qīngshàonián gōngzuò. « Le gouvernement devrait donner du travail aux jeunes. » \zh{应该} se place avant le verbe. \\
\hline
\end{tabularx}
\end{center}

\begin{attention}[Les trois contrôles avant de rendre votre copie]
Avant de rendre toute production écrite, relisez en vous posant trois questions : (1) \textbf{Tons} : ai-je marqué le bon ton sur chaque syllabe du pinyin, en particulier \zh{犯罪} (fànzuì), \zh{原因} (yuányīn), \zh{法律} (fǎlǜ) ? (2) \textbf{Ordre des mots} : la cause est-elle avant la conséquence ? Le complément de lieu (\zh{在我们城市}) est-il avant le verbe ? \zh{应该} est-il devant le verbe ? (3) \textbf{Caractères} : ai-je écrit \zh{犯} et non \zh{饭}, \zh{法} et non \zh{去}, \zh{青} et non \zh{清} ?
\end{attention}

\courssub{Complément utile à l'examen : la mini-dissertation type Bac}

\begin{aretenir}[Mini-dissertation type baccalauréat]
Au baccalauréat, on peut vous demander un texte argumentatif de huit à douze phrases sur un sujet de société. La structure attendue est :\\
\textbf{Introduction} (1-2 phrases) : présenter le sujet avec \zh{在我们城市 / 现在，……是一个严重的问题。}\\
\textbf{Développement en deux parties} : partie 1, les causes (\zh{犯罪的原因有……：第一……第二……}) ; partie 2, les solutions (\zh{为了解决这个问题，……应该……，……也应该……}).\\
\textbf{Conclusion} (1-2 phrases) : ouvrir avec \zh{只有……才能……} ou \zh{只有这样，才能……}.\\
C'est exactement la trame de cette leçon, simplement allongée : qui maîtrise la production de cinq phrases sait déjà écrire la dissertation du Bac.
\end{aretenir}

\begin{savaistu}[La loi chinoise sur la prévention de la délinquance juvénile]
En 2021, la Chine a promulgué une version révisée de sa loi sur la prévention de la délinquance juvénile (\zh{预防未成年人犯罪法}, Yùfáng wèichéngniánrén fànzuì fǎ). Cette loi insiste sur l'\cle{éducation} plutôt que sur la punition : elle demande aux familles, aux écoles et à la société de travailler ensemble pour aider les jeunes en difficulté avant qu'ils ne commettent des délits. C'est exactement l'idée de la phrase-modèle de notre leçon : \zh{只有家庭、学校和社会一起努力，才能帮助青少年走上正确的道路。} Au Cameroun aussi, les éducateurs répètent que prévenir vaut mieux que punir : les deux pays partagent la même conviction.
\end{savaistu}

\newpage

\coursec{Activité d'intégration}

\coursec{Leçon 5 — Vocabulaire et compréhension de texte : délinquance juvénile}

\courssub{Mise en route et présentation du thème}

\begin{experience}[Brainstorming : « Jeunes en danger ? »]
\begin{enumerate}
\item \textbf{Observation d'images.} Le professeur projette ou affiche trois photos : (a) des lycéens en uniforme devant le lycée Leclerc à Yaoundé ; (b) un groupe d'adolescents fumant dans un quartier de Douala (Bessengue) ; (c) une famille réunie autour du repas (ndolé, plantains) qui discute.
\item \textbf{Mots-clés.} En binôme, listez 5 mots chinois (caractères + pinyin) qui vous viennent à l'esprit pour décrire la situation (b) et 5 mots pour la situation (c). Exemple : \zh{吸烟} (xīyān), \zh{逃学} (táoxué), \zh{关心} (guānxīn), \zh{教育} (jiàoyù).
\item \textbf{Mise en commun.} Au tableau, construction d'une carte mentale : \zh{青少年问题} au centre, branches \zh{原因} (famille, école, société, internet) et \zh{解决办法} (éducation, loi, psychologie, sport).
\end{enumerate}
\end{experience}

\begin{popfigure}
\begin{tikzpicture}[mindmap, grow cyclic, every node/.style=concept, concept color=popblue!60,
    level 1/.append style={level distance=4.5cm, sibling angle=90},
    level 2/.append style={level distance=3cm, sibling angle=45}]
\node[concept color=popdark] {青少年问题}
    child[concept color=poporange] { node {原因}
        child { node {家庭} }
        child { node {学校} }
        child { node {社会} }
        child { node {网络} }
    }
    child[concept color=popgreen] { node {表现}
        child { node {逃学} }
        child { node {打架} }
        child { node {吸毒} }
        child { node {偷窃} }
    }
    child[concept color=poppurple] { node {对策}
        child { node {家庭教育} }
        child { node {学校管理} }
        child { node {法律保护} }
        child { node {自律} }
    };
\end{tikzpicture}
\legende{Carte mentale : causes, manifestations et solutions de la délinquance juvénile}
\end{popfigure}

\begin{definition}[Définition du thème]
\zh{青少年犯罪} (qīngshàonián fànzuì) désigne les actes illégaux commis par des mineurs (généralement 14--18 ans en Chine, 13--18 ans au Cameroun) : vol, violence, trafic de drogue, vandalisme, etc. Le thème de cette leçon porte sur le \textbf{vocabulaire} nécessaire pour en parler, en comprendre les \cle{causes} (\zh{原因}) et proposer des \cle{solutions} (\zh{办法}), tout en réinvestissant les structures \zh{因为……所以……}, \zh{有的……有的……} et \zh{只有……才……}.
\end{definition}

\courssub{Texte support : lecture guidée}

\begin{textebox}[Texte original : 青少年犯罪的原因与对策 — Causes et solutions de la délinquance juvénile]
近年来，青少年犯罪在喀麦隆和中国都引起了社会的关注。\\
Jìnnián lái, qīngshàonián fànzuì zài Kāmàilóng hé Zhōngguó dōu yǐnqǐle shèhuì de guānzhù.\\
« Ces dernières années, la délinquance juvénile a suscité l'attention de la société au Cameroun comme en Chine. »\\[4pt]
\zh{有的青少年因为家庭缺乏关爱，所以走上了歧路。}\\
Yǒude qīngshàonián yīnwèi jiātíng quēfá guān'ài, suǒyǐ zǒushàngle qílù.\\
« Certains jeunes, parce que leur famille manque d'affection, ont emprunté la mauvaise voie. »\\[4pt]
\zh{有的学生在学校受到坏同学的影响，逃课、打架、甚至偷窃。}\\
Yǒude xuésheng zài xuéxiào shòudào huài tóngxué de yǐngxiǎng, tàokè, dǎjià, shènzhì tōuqiè.\\
« Certains élèves, sous l'influence de mauvais camarades à l'école, sèchent les cours, se battent, voire volent. »\\[4pt]
\zh{网络游戏和短视频也让一些孩子沉迷，不想学习。}\\
Wǎngluò yóuxì hé duǎn shìpín yě ràng yìxiē háizi chénmí, bù xiǎng xuéxí.\\
« Les jeux en ligne et les vidéos courtes rendent aussi certains enfants accros, ne voulant plus étudier. »\\[4pt]
\zh{只有家庭、学校和社会一起努力，才能保护青少年健康成长。}\\
Zhǐyǒu jiātíng, xuéxiào hé shèhuì yìqǐ nǔlì, cái néng bǎohù qīngshàonián jiànkāng chéngzhǎng.\\
« Seule une action conjointe de la famille, de l'école et de la société peut protéger la croissance saine des adolescents. »\\[4pt]
\zh{我们要加强法治教育，帮助他们分辨是非，远离犯罪。}\\
Wǒmen yào jiāqiáng fǎzhì jiàoyù, bāngzhù tāmen fēnbiàn shìfēi, yuǎnlí fànzuì.\\
« Nous devons renforcer l'éducation juridique, les aider à distinguer le bien du mal et à s'éloigner du crime. »
\end{textebox}

\begin{methode}[Lecture guidée en 4 étapes]
\begin{enumerate}
\item \textbf{Lecture globale (1 min) :} Lisez le texte sans dictionnaire. Identifiez le sujet principal (titre implicite) et le ton (informatif, préventif).
\item \textbf{Repérage des structures (3 min) :} Surlignez dans le texte : \cle{une phrase avec \zh{因为……所以……}}, \cle{une phrase avec \zh{有的……有的……}}, \cle{une phrase avec \zh{只有……才……}}.
\item \textbf{Lecture détaillée (5 min) :} Avec le glossaire (section 3), traduisez mentalement chaque paragraphe. Notez les mots inconnus.
\item \textbf{Analyse socioculturelle (2 min) :} Comparez : le texte cite le Cameroun et la Chine. Quels points communs ? (Famille, école, internet). Différences ? (Législation, rôle du quartier/village au Cameroun).
\end{enumerate}
\end{methode}

\courssub{Glossaire du thème (Vocabulaire de base HSK 3--4)}

\begin{center}
\begin{tabularx}{\linewidth}{|X|X|X|X|}
\hline
\rowcolor{popblueL} Caractères & Pinyin & Nature & Traduction \\
\hline
青少年犯罪 & qīngshàonián fànzuì & n. & délinquance juvénile \\
\hline
原因 & yuányīn & n. & cause, raison \\
\hline
保护 & bǎohù & v. & protéger \\
\hline
帮助 & bāngzhù & v. & aider, assistance \\
\hline
教育 & jiàoyù & n./v. & éducation ; éduquer \\
\hline
办法 & bànfǎ & n. & solution, moyen, méthode \\
\hline
家长 & jiāzhǎng & n. & parent d'élève, tuteur \\
\hline
警察 & jǐngchá & n. & policier \\
\hline
努力 & nǔlì & v./adj. & faire des efforts, s'efforcer ; assidu \\
\hline
一起 & yìqǐ & adv. & ensemble \\
\hline
关爱 & guān'ài & v./n. & affection, soin, bienveillance \\
\hline
缺乏 & quēfá & v. & manquer de, faire défaut \\
\hline
歧路 & qílù & n. & mauvaise voie (fig.) \\
\hline
影响 & yǐngxiǎng & n./v. & influence ; influencer \\
\hline
逃课 & tàokè & v. & sécher les cours \\
\hline
偷窃 & tōuqiè & v. & voler (à la tire, vol simple) \\
\hline
沉迷 & chénmí & v. & être accro, s'adonner excessivement \\
\hline
法治 & fǎzhì & n. & État de droit, règle de droit \\
\hline
分辨 & fēnbiàn & v. & distinguer, discerner \\
\hline
是非 & shìfēi & n. & le bien et le mal, le vrai et le faux \\
\hline
远离 & yuǎnlí & v. & s'éloigner, éviter \\
\hline
责任 & zérèn & n. & responsabilité \\
\hline
\end{tabularx}
\end{center}

\begin{propriete}[Points de vocabulaire clés]
\begin{itemize}
\item \zh{青少年} (qīngshàonián) = \zh{青年} (jeune adulte) + \zh{少年} (adolescent). Désigne les 12--18 ans environ.
\item \zh{犯罪} (fànzuì) : \zh{犯} (enfreindre) + \zh{罪} (crime/péché). Ne pas confondre avec \zh{犯错} (fàncuò, faire une erreur).
\item \zh{关爱} (guān'ài) : registre plus formel/affectif que \zh{关心} (guānxīn, se soucier de).
\item \zh{歧路} (qílù) : litt. « chemin de traverse » $\rightarrow$ fig. « mauvaise voie ». Expression figée : \zh{走上歧路}.
\item \zh{法治教育} (fǎzhì jiàoyù) : éducation à la citoyenneté et au droit (programme commun Cameroun/Chine).
\end{itemize}
\end{propriete}

\courssub{Clés (radicaux) et ordre des traits}

\begin{definition}[Analyse de 6 caractères du thème]
Le tableau ci-dessous présente la décomposition graphique pour faciliter la mémorisation et l'écriture correcte (ordre des traits : horizontal avant vertical, gauche avant droite, haut avant bas, extérieur avant intérieur).
\end{definition}

\begin{center}
\begin{tabularx}{\linewidth}{|c|X|X|X|X|X|}
\hline
\rowcolor{popblueL} Car. & Radical (Clé) & Sens du radical & Structure / Décomposition & Nb traits & Ordre des traits (résumé) \\
\hline
\zh{犯} & \zh{犭} (quǎn, chien) & Animal / comportement & \zh{犭} + \zh{然} (rán) & 5 + 12 = 17 & 1. \zh{犭} (point, courbe, trait, crochet, point) 2. \zh{然} (haut $\to$ bas) \\
\hline
\zh{罪} & \zh{网} (wǎng, filet) & Filet, piège & \zh{网} + \zh{非} (fēi) & 6 + 8 = 14 & 1. \zh{网} (cadre ext. puis int.) 2. \zh{非} (haut $\to$ bas) \\
\hline
\zh{原} & \zh{厶} (sī, privé) & Petit, détail & \zh{厶} + \zh{泉} (quán) & 4 + 9 = 10 & 1. \zh{厶} 2. \zh{泉} (eau \zh{水} + blanc \zh{白}) \\
\hline
\zh{护} & \zh{扌} (shǒu, main) & Action de la main & \zh{扌} + \zh{护} (phonétique \zh{户} + \zh{言}) & 3 + 14 = 17 & 1. \zh{扌} (horizontal, crochet, point) 2. Partie sup. \zh{互} puis \zh{言} \\
\hline
\zh{教} & \zh{攵} (bū, taper) & Action d'enseigner & \zh{孝} (xiào) + \zh{攵} & 7 + 4 = 11 & 1. \zh{孝} (haut $\to$ bas) 2. \zh{攵} (point, horiz., crochet, point) \\
\hline
\zh{警} & \zh{言} (yán, parole) & Parole, avertissement & \zh{言} + \zh{敬} (jìng) & 7 + 12 = 19 & 1. \zh{言} (point, horiz., horiz., horiz., vertical, horiz., point) 2. \zh{敬} (complexe, haut $\to$ bas) \\
\hline
\end{tabularx}
\end{center}

\begin{experience}[Atelier écriture : « Course aux radicaux »]
En groupes de 4. Chaque groupe reçoit 6 cartes « radical » (\zh{犭}, \zh{网}, \zh{厶}, \zh{扌}, \zh{攵}, \zh{言}) et 6 cartes « caractère complet ». Relier le radical au caractère en moins de 2 minutes. Puis, un élève écrit le caractère au tableau en dictant l'ordre des traits à voix haute en chinois : \zh{第一笔：横} (di yī bǐ: héng), \zh{第二笔：竖} (dì èr bǐ: shù)...
\end{experience}

\courssub{Compréhension du texte}

\begin{exercice}[Questions sur le texte « Causes et solutions »]
Répondez en chinois (phrase complète) ou en français selon la consigne.
\begin{enumerate}
\item \textbf{Vrai ou Faux ?} Justifiez par une phrase du texte.
\begin{itemize}
\item[a)] Le texte dit que la délinquance ne concerne que la Chine. (Faux : \zh{在喀麦隆和中国都引起了社会的关注})
\item[b)] L'école n'a aucune responsabilité. (Faux : \zh{在学校受到坏同学的影响})
\item[c)] La solution repose uniquement sur la police. (Faux : \zh{只有家庭、学校和社会一起努力})
\end{itemize}
\item \textbf{Questions ouvertes (réponse en chinois) :}
\begin{itemize}
\item[a)] 根据课文，青少年走上歧路的原因有哪些？(D'après le texte, quelles sont les causes ?) \\
\emph{Réponse attendue :} \zh{因为家庭缺乏关爱，受坏同学影响，沉迷网络游戏和短视频。}
\item[b)] 课文提出了什么办法来保护青少年？(Quelles solutions le texte propose-t-il ?) \\
\emph{Réponse attendue :} \zh{家庭、学校和社会一起努力，加强法治教育，帮助他们分辨是非，远离犯罪。}
\end{itemize}
\item \textbf{Travail lexical :} Trouvez dans le texte les synonymes ou paraphrases de :
\begin{itemize}
\item[a)] « mauvais chemin » $\rightarrow$ \zh{歧路}
\item[b)] « ne pas avoir assez » $\rightarrow$ \zh{缺乏}
\item[c)] « distinguer le bien du mal » $\rightarrow$ \zh{分辨是非}
\end{itemize}
\item \textbf{Traduction (Version) :} Traduisez en français le 4\textsuperscript{e} paragraphe : \zh{只有家庭、学校和社会一起努力，才能保护青少年健康成长。}
\end{enumerate}
\end{exercice}

\begin{corrige}[Exercice Compréhension]
\begin{enumerate}
\item[a] Faux. \zh{近年来，青少年犯罪在喀麦隆和中国都引起了社会的关注。} \\
\item[b] Faux. \zh{有的学生在学校受到坏同学的影响，逃课、打架、甚至偷窃。} \\
\item[c] Faux. \zh{只有家庭、学校和社会一起努力，才能保护青少年健康成长。} \\
\item[a] \zh{根据课文，原因是：家庭缺乏关爱、受坏同学影响、沉迷网络。} \\
\item[b] \zh{办法是：家庭、学校、社会一起努力，加强法治教育，帮助分辨是非，远离犯罪。} \\
\item[a] \zh{歧路} ; [b] \zh{缺乏} ; [c] \zh{分辨是非}. \\
\item « Seule une action conjointe de la famille, de l'école et de la société peut protéger la croissance saine des adolescents. »
\end{enumerate}
\end{corrige}

\courssub{Production guidée et réinvestissement (Activité d'intégration)}

\courssub{Objectifs de l'activité}

À l'issue de cette activité d'intégration, l'élève sera capable de :
\begin{itemize}
\item réinvestir à l'oral le vocabulaire du thème de la délinquance juvénile (\zh{青少年犯罪}, \zh{原因}, \zh{家庭}, \zh{学校}, \zh{帮助}, \zh{保护}, \zh{教育}, \zh{关爱}, \zh{法治}…) dans une situation de communication authentique ;
\item employer correctement les structures étudiées : \zh{因为……所以……}, \zh{有的……有的……}, \zh{只有……才……} ;
\item comprendre un échange oral en chinois et en prendre des notes (stratégies : repérer les mots-clés, les causes, les solutions) ;
\item produire par écrit un court article de presse en chinois (6 phrases) résumant un débat ;
\item travailler en équipe et respecter un rôle social dans un jeu de simulation.
\end{itemize}

\courssub{Situation}

La CRTV, la télévision nationale du Cameroun, prépare une émission spéciale en chinois pour les lycéens. Le thème du jour : « Comment protéger nos jeunes ? ». Sur le plateau, un animateur reçoit quatre invités : un parent d'élève de Yaoundé, un professeur de lycée, un policier de Douala et un adolescent de seize ans. Chacun donne son avis sur les causes de la délinquance juvénile et propose des solutions. Ta classe joue l'émission en direct !

\begin{textebox}[Annonce de l'émission — CRTV]
大家好！欢迎收看今天的节目。今天的主题是：怎么保护我们的年轻人？\\
Dàjiā hǎo! Huānyíng shōukàn jīntiān de jiémù. Jīntiān de zhǔtí shì: zěnme bǎohù wǒmen de niánqīngrén?\\
« Bonjour à tous ! Bienvenue dans l'émission d'aujourd'hui. Le thème du jour : comment protéger nos jeunes ? »\\[4pt]
为什么有些青少年会犯罪？是因为家庭的问题，还是因为学校的问题？\\
Wèishénme yǒuxiē qīngshàonián huì fànzuì? Shì yīnwèi jiātíng de wèntí, háishì yīnwèi xuéxiào de wèntí?\\
« Pourquoi certains adolescents commettent-ils des délits ? Est-ce à cause de problèmes familiaux ou de problèmes scolaires ? »\\[4pt]
今天我们请了四位客人：一位家长、一位老师、一位警察和一个十六岁的孩子。请听他们说！\\
Jīntiān wǒmen qǐngle sì wèi kèrén: yí wèi jiāzhǎng, yí wèi lǎoshī, yí wèi jǐngchá hé yí ge shíliù suì de háizi. Qǐng tīng tāmen shuō!\\
« Aujourd'hui, nous avons invité quatre personnes : un parent, un professeur, un policier et un jeune de seize ans. Écoutons-les ! »
\end{textebox}

\courssub{Consignes}

\begin{enumerate}
\item \textbf{Former les équipes.} La classe se divise en groupes. Un groupe de cinq élèves joue le plateau : un animateur et quatre invités (le parent, le professeur, le policier, l'adolescent). Les autres élèves sont le public et les journalistes.
\item \textbf{Préparer son rôle (10 minutes).} Chaque invité prépare un court discours de 3 à 4 phrases en chinois. Il doit obligatoirement utiliser :
\begin{itemize}
\item au moins \cle{5 mots du glossaire} de la leçon (par exemple \zh{青少年犯罪}, \zh{原因}, \zh{家庭}, \zh{学校}, \zh{教育}, \zh{帮助}, \zh{保护}, \zh{关爱}, \zh{法治}, \zh{责任}) ;
\item au moins \cle{une structure de la leçon} : \zh{因为……所以……} ou \zh{有的……有的……} ou \zh{只有……才……}.
\end{itemize}
L'animateur prépare ses questions : \zh{你觉得青少年犯罪的原因是什么？} (Nǐ juéde qīngshàonián fànzuì de yuányīn shì shénme ?) et \zh{我们应该怎么帮助他们？} (Wǒmen yīnggāi zěnme bāngzhù tāmen ?).
\item \textbf{Jouer l'émission (10 minutes).} L'animateur pose ses questions, chaque invité répond à son tour. Le public écoute en silence.
\item \textbf{Prendre des notes.} Pendant l'émission, chaque élève du public complète une fiche d'écoute : il note pour chaque invité \cle{une cause} citée et \cle{une solution} proposée, en chinois ou en français.
\item \textbf{Répondre aux questions de compréhension.} Après l'émission, réponds par écrit, en chinois, à trois questions :
\begin{itemize}
\item[(a)] \zh{家长认为问题的原因是什么？}
\item[(b)] \zh{警察提出了什么办法？}
\item[(c)] \zh{那个十六岁的孩子说了什么？}
\end{itemize}
\item \textbf{Rédiger l'article de presse.} En binôme, écris un court article de 6 phrases en chinois pour le journal du lycée : titre, présentation de l'émission, deux opinions d'invités, une solution, une phrase de conclusion. Utilise au moins une fois \zh{因为……所以……} et une fois \zh{只有……才……}.
\end{enumerate}

\begin{attention}[Pièges à éviter pendant la production]
\begin{itemize}
\item Ne dis pas \zh{因为……所以……} dans toutes les phrases : une ou deux fois suffisent. En chinois, on peut aussi utiliser \zh{所以} seul en début de phrase (ex. \zh{家庭缺乏关爱，所以孩子走上了歧路。}).
\item Avec \zh{只有……才……}, n'oublie jamais le second élément \zh{才} : \zh{只有大家一起努力，才能保护年轻人。} $\neq$ \zh{只有大家一起努力，能保护年轻人。} (Sans \zh{才}, la nuance de « condition nécessaire et suffisante » disparaît).
\item Attention à l'ordre des mots : le complément de cause avec \zh{因为} se place \textbf{avant} le sujet principal de la conséquence, jamais après. \textbf{Correct :} \zh{因为下雨，所以我不去。} \textbf{Incorrect :} \zh{我不去，因为下雨。} (sauf structure \zh{之所以……是因为……} niveau supérieur).
\item Soigne les tons : \zh{fànzuì} (犯罪, 4\textsuperscript{e} + 4\textsuperscript{e} ton) $\neq$ \zh{fǎnzuò} (反坐, 3\textsuperscript{e} + 4\textsuperscript{e} ton, rare). \zh{quēfá} (缺乏, 1\textsuperscript{e} + 2\textsuperscript{e} ton) $\neq$ \zh{quēfà} (incorrect).
\item Classificateur : \zh{一位} (yí wèi) pour les personnes respectées (prof, policier, parent), \zh{一个} (yí ge) ou \zh{一名} (yí míng) pour l'adolescent.
\end{itemize}
\end{attention}

\courssub{Ressources à mobiliser}

\begin{propriete}[Structures de la leçon à réemployer]
\begin{itemize}
\item \textbf{Cause et conséquence} : \zh{因为 + cause，所以 + conséquence。}\\
\zh{因为家庭缺乏关爱，所以一些孩子走上了歧路。} (Yīnwèi jiātíng quēfá guān'ài, suǒyǐ yìxiē háizi zǒushàngle qílù.)
\item \textbf{Énumération / Contraste} : \zh{有的……有的……} (Sujet \zh{有的} + Prédicat 1，Sujet \zh{有的} + Prédicat 2)。\\
\zh{有的青少年逃课，有的青少年沉迷手机。} (Yǒude qīngshàonián tàokè, yǒude qīngshàonián chénmí shǒujī.)
\item \textbf{Condition nécessaire (seule condition pour le résultat)} : \zh{只有 + condition，才 + résultat。}\\
\zh{只有加强法治教育，才能让青少年远离犯罪。} (Zhǐyǒu jiāqiáng fǎzhì jiàoyù, cái néng ràng qīngshàonián yuǎnlí fànzuì.)
\end{itemize}
\end{propriete}

\begin{center}
\begin{tabularx}{\linewidth}{|X|X|X|X|}
\hline
\rowcolor{popblueL} Caractères & Pinyin & Nature & Traduction \\
\hline
青少年犯罪 & qīngshàonián fànzuì & n. & délinquance juvénile \\
\hline
原因 & yuányīn & n. & cause, raison \\
\hline
保护 & bǎohù & v. & protéger \\
\hline
帮助 & bāngzhù & v. & aider \\
\hline
教育 & jiàoyù & n./v. & éducation ; éduquer \\
\hline
办法 & bànfǎ & n. & solution, moyen \\
\hline
家长 & jiāzhǎng & n. & parent d'élève \\
\hline
警察 & jǐngchá & n. & policier \\
\hline
努力 & nǔlì & v./adj. & faire des efforts \\
\hline
一起 & yìqǐ & adv. & ensemble \\
\hline
关爱 & guān'ài & v./n. & affection, bienveillance \\
\hline
法治 & fǎzhì & n. & État de droit \\
\hline
责任 & zérèn & n. & responsabilité \\
\hline
\end{tabularx}
\end{center}

\courssub{Proposition de production et corrigé commenté}

\begin{exoresolu}[Article de presse modèle — Journal du lycée]
Rédige un article de 6 phrases résumant le débat télévisé « Comment protéger nos jeunes ? ».
\tcblower
\textbf{Production modèle :}\\[4pt]
\zh{怎么保护我们的年轻人？}\\
Zěnme bǎohù wǒmen de niánqīngrén?\\
« Comment protéger nos jeunes ? »\\[3pt]
\zh{昨天，我们在学校看了CRTV的一个节目，主题是青少年犯罪。}\\
Zuótiān, wǒmen zài xuéxiào kànle CRTV de yí ge jiémù, zhǔtí shì qīngshàonián fànzuì.\\
« Hier, nous avons regardé à l'école une émission de la CRTV dont le thème était la délinquance juvénile. »\\[3pt]
\zh{家长说，因为有的家庭缺乏关爱，所以孩子走上了歧路。}\\
Jiāzhǎng shuō, yīnwèi yǒude jiātíng quēfá guān'ài, suǒyǐ háizi zǒushàngle qílù.\\
« Le parent a dit que, parce que certaines familles manquent d'affection, les enfants ont emprunté la mauvaise voie. »\\[3pt]
\zh{老师说，有的学生逃课，有的学生打架，学校应该加强教育。}\\
Lǎoshī shuō, yǒude xuésheng tàokè, yǒude xuésheng dǎjià, xuéxiào yīnggāi jiāqiáng jiàoyù.\\
« Le professeur a dit que certains élèves sèchent les cours, d'autres se battent, et que l'école devrait renforcer l'éducation. »\\[3pt]
\zh{警察说，只有家庭、学校和警察一起努力，才能保护年轻人。}\\
Jǐngchá shuō, zhǐyǒu jiātíng, xuéxiào hé jǐngchá yìqǐ nǔlì, cái néng bǎohù niánqīngrén.\\
« Le policier a dit que seuls la famille, l'école et la police, en travaillant ensemble, peuvent protéger les jeunes. »\\[3pt]
\zh{我们都觉得，保护青少年是大家的责任。}\\
Wǒmen dōu juéde, bǎohù qīngshàonián shì dàjiā de zérèn.\\
« Nous pensons tous que protéger les adolescents est la responsabilité de tous. »\\[6pt]
\textbf{Commentaire des choix de langue :}
\begin{itemize}
\item \textbf{Phrase 1 (Titre) :} Question accroche, vocabulaire du thème (\zh{保护}, \zh{年轻人}).
\item \textbf{Phrase 2 (Contexte) :} \zh{在学校看了} (regardé à l'école) structure lieu + verbe. \zh{主题是} (thème est).
\item \textbf{Phrase 3 (Opinion 1 - Parent) :} Discours rapporté \zh{家长说，...}. Structure \zh{因为……所以……} correcte. Vocabulaire précis : \zh{缺乏关爱}, \zh{走上歧路}.
\item \textbf{Phrase 4 (Opinion 2 - Professeur) :} Structure \zh{有的……有的……} pour énumérer deux comportements distincts. \zh{应该} (devoir moral) + \zh{加强教育}.
\item \textbf{Phrase 5 (Solution - Policier) :} Structure \zh{只有……才……} maîtrisée. Coordination \zh{家庭、学校和警察} + \zh{一起努力}.
\item \textbf{Phrase 6 (Conclusion) :} \zh{都觉得} (tous pensent que). \zh{责任} (responsabilité) mot-clé citoyenneté.
\end{itemize}
\end{exoresolu}

\courssub{Bilan de l'activité}

\begin{aretenir}[Ce que cette activité t'a appris]
\begin{itemize}
\item \textbf{Compréhension orale} : tu as écouté un échange en chinois, repéré les mots-clés (\zh{原因}, \zh{办法}) et distingué les causes des solutions.
\item \textbf{Lexique} : tu as prononcé et employé dans un contexte réel au moins cinq mots du thème de la délinquance juvénile (\zh{青少年犯罪}, \zh{关爱}, \zh{法治}, \zh{歧路}...).
\item \textbf{Grammaire} : tu as utilisé \zh{因为……所以……} (cause), \zh{有的……有的……} (énumération), \zh{只有……才……} (condition nécessaire) dans tes propres phrases.
\item \textbf{Expression écrite} : tu as rédigé un article de presse structuré (titre, faits, opinions, conclusion) en six phrases correctes.
\item \textbf{Compétence sociale} : tu as tenu un rôle (animateur, parent, prof, policier, ado), respecté la parole des autres et travaillé en équipe, comme sur un vrai plateau de télévision camerounais.
\end{itemize}
\end{aretenir}

\begin{experience}[Pour aller plus loin : Extension et évaluation]
\begin{enumerate}
\item \textbf{Enregistrement vidéo.} En équipe, enregistrez votre émission avec un téléphone (mode « reportage »). Diffusez-la en classe ou sur le groupe WhatsApp de la classe. Auto-évaluation : grille « Prononciation / Tons », « Vocabulaire utilisé », « Structures correctes », « Fluidité ».
\item \textbf{Nouveau personnage.} Invitez un \zh{区长} (qūzhǎng, chef de quartier) ou une \zh{社区工作者} (shèqū gōngzuòzhě, travailleur social communautaire) de Buea ou Yaoundé. Il/Elle utilisera la structure \zh{我觉得……} (wǒ juéde…, « je pense que… ») et le vocabulaire \zh{社区} (communauté), \zh{预防} (prévention).
\item \textbf{Débat écrit (Devoir) :} « \zh{预防青少年犯罪，家庭比学校更重要。你同意吗？} » (Prévenir la délinquance, la famille est plus importante que l'école. Êtes-vous d'accord ?). Rédigez 80--100 caractères en utilisant \zh{因为……所以……} et \zh{但是} (dànshì, mais).
\end{enumerate}
\end{experience}

\begin{savaistu}[Regard croisé Cameroun -- Chine : Protection des mineurs]
\begin{itemize}
\item \textbf{Âge de la responsabilité pénale :} En Chine, depuis 2021 (amendement Code pénal), les 12--14 ans peuvent être poursuivis pour crimes graves (homicide, blessure grave) sur décision du Parquet suprême. Au Cameroun, la majorité pénale est à 13 ans (Code pénal, art. 80), mais la responsabilité civile commence plus tôt.
\item \textbf{Écoles spécialisées :} La Chine dispose d'\zh{专门学校} (zhuānmén xuéxiào, écoles spécialisées de correction) pour mineurs délinquants (éducation + formation professionnelle). Le Cameroun a des \emph{Centres de Rééducation des Mineurs} (ex. Centre de Bétamba) mais les effectifs et moyens sont limités.
\item \textbf{Rôle de la communauté :} Au Cameroun, le \zh{chef de quartier} / \zh{chef de village} et les \zh{associations de parents d'élèves} (APEE) jouent un rôle majeur de médiation avant la police. En Chine, les \zh{居委会} (jūwěihuì, comités de quartier) et \zh{村委会} (cūnwěihuì, comités de village) ont un pouvoir administratif fort de surveillance et d'aide aux familles en difficulté.
\item \textbf{Vocabulaire comparé :} \zh{未成年人保护法} (Wèichéngniánrén bǎohù fǎ, Loi protection mineurs - Chine, révisée 2020) vs \textbf{Loi n°2010/012} (Protection de l'enfant - Cameroun).
\end{itemize}
\end{savaistu}

\newpage

\coursec{Méthodes et exercices résolus}

\coursec{Leçon 5 — Vocabulaire et compréhension de texte : délinquance juvénile}

\courssub{Mise en route et présentation du thème}

\begin{experience}[Entrée en matière : un phénomène de société]
\begin{enumerate}
\item \textbf{Observation d'images.} Le professeur projette deux photos contrastées : des jeunes en uniforme scolaire studieux (Cameroun / Chine) vs des jeunes oisifs, fumant ou en conflit avec la police. Vocabulaire émergent : \zh{青少年} (adolescents), \zh{犯罪} (crime/délinquant), \zh{社会问题} (problème social).
\item \textbf{Brainstorming guidé.} En groupes de 4, les élèves listent en français 3 causes possibles de la délinquance juvénile (famille, école, pairs, médias, pauvreté). Mise en commun au tableau ; le professeur fournit les termes chinois clés : \zh{家庭教育}, \zh{不良朋友}, \zh{缺乏关爱}.
\item \textbf{Annonce de l'objectif.} « Aujourd'hui, nous allons lire un texte chinois sur ce thème, en maîtriser le vocabulaire (HSK 3), analyser la structure des phrases (causes \textrightarrow{} solutions) et produire un court paragraphe personnel. »
\end{enumerate}
\end{experience}

\courssub{Texte support : lecture guidée}

\begin{textebox}[Texte original : La délinquance juvénile, un défi commun]
近年来，青少年犯罪是一个严重的社会问题。一些年轻人犯罪的原因有两个：第一，家庭教育不够；第二，他们交了不好的朋友。为了解决这个问题，家庭、学校和社会应该一起努力，帮助这些年轻人。\\
\emph{Jìnnián lái, qīngshàonián fànzuì shì yí ge yánzhòng de shèhuì wèntí. Yìxiē niánqīngrén fànzuì de yuányīn yǒu liǎng ge: dì-yī, jiātíng jiàoyù bù gòu; dì-èr, tāmen jiāo le bù hǎo de péngyou. Wèile jiějué zhège wèntí, jiātíng, xuéxiào hé shèhuì yīnggāi yìqǐ nǔlì, bāngzhù zhèxiē niánqīngrén.}\\
« Ces dernières années, la délinquance juvénile est un grave problème social. Les causes de la criminalité chez certains jeunes sont au nombre de deux : premièrement, l'éducation familiale est insuffisante ; deuxièmement, ils se sont fait de mauvais amis. Pour résoudre ce problème, la famille, l'école et la société doivent faire des efforts ensemble pour aider ces jeunes. »
\end{textebox}

\begin{methode}[Lire un texte sur la délinquance juvénile en 5 étapes]
\begin{enumerate}
\item \textbf{Lecture globale silencieuse.} Lis le texte une première fois sans t'arrêter sur les caractères inconnus. Repère le sujet général : un phénomène social (\zh{青少年犯罪}, la délinquance juvénile), ses causes et ses solutions.
\item \textbf{Repérage des mots-clés du thème.} Souligne les mots du champ lexical : \zh{犯罪} (commettre un crime), \zh{原因} (cause), \zh{家庭} (famille), \zh{学校} (école), \zh{社会} (société), \zh{帮助} (aider), \zh{教育} (éducation).
\item \textbf{Découpage des phrases.} En chinois, la virgule ， sépare des propositions autour d'un même sujet. Identifie le sujet de chaque proposition : souvent il n'est dit qu'une fois (sujet zéro / ellipse du sujet).
\item \textbf{Lecture à voix haute en respectant les tons.} Lis phrase par phrase en prononçant le pinyin avec les tons ; cela fixe le sens et la mémorisation.
\item \textbf{Reformulation en français.} Résume chaque paragraphe (ici chaque phrase complexe) en une phrase française simple avant de répondre aux questions.
\end{enumerate}
\end{methode}

\courssub{Glossaire du thème}

\begin{aretenir}[Vocabulaire clé du thème « Délinquance juvénile » (HSK 2-3)]
\begin{center}
\begin{tabularx}{\linewidth}{|X|X|X|X|}
\hline
\rowcolor{popblueL} \textbf{Caractères} & \textbf{Pinyin} & \textbf{Nature} & \textbf{Traduction} \\
\hline
青少年 & qīngshàonián & nom & adolescent(s), jeune(s) \\
\hline
犯罪 & fànzuì & verbe / nom & commettre un crime / délinquance, crime \\
\hline
严重 & yánzhòng & adjectif & grave, sérieux \\
\hline
社会问题 & shèhuì wèntí & nom & problème social \\
\hline
年轻人 & niánqīngrén & nom & jeune, jeune personne \\
\hline
原因 & yuányīn & nom & cause, raison \\
\hline
家庭教育 & jiātíng jiàoyù & nom & éducation familiale \\
\hline
不够 & bù gòu & verbe / adj. & insuffisant, ne pas suffire \\
\hline
交朋友 & jiāo péngyou & loc. verbale & se faire des amis, lier amitié \\
\hline
不好 & bù hǎo & adjectif & mauvais, pas bien \\
\hline
解决 & jiějué & verbe & résoudre (un problème) \\
\hline
学校 & xuéxiào & nom & école \\
\hline
一起 & yìqǐ & adverbe & ensemble \\
\hline
努力 & nǔlì & verbe / adv. & faire des efforts / avec effort \\
\hline
帮助 & bāngzhù & verbe / nom & aider / aide \\
\hline
法律 & fǎlǜ & nom & loi \\
\hline
违法 & wéifǎ & verbe & violer la loi, illégal \\
\hline
预防 & yùfáng & verbe & prévenir, prévenir (un risque) \\
\hline
关爱 & guān'ài & verbe / nom & chérir, prendre soin de / affection \\
\hline
\end{tabularx}
\end{center}
\end{aretenir}

\courssub{Clés (radicaux) et ordre des traits}

\begin{propriete}[Analyse de caractères clés du texte]
\begin{center}
\begin{tabularx}{\linewidth}{|X|X|X|X|}
\hline
\rowcolor{popblueL} \textbf{Caractère} & \textbf{Clé (Radical)} & \textbf{Sens de la clé} & \textbf{Nombre de traits / Structure} \\
\hline
犯 & \zh{犭} (quǎn, chien) & Animal / comportement & 5 traits ; structure gauche-droite : \zh{犭} + \zh{范} \\
\hline
罪 & \zh{罒} (wǎng, filet) & Filet, capture & 13 traits ; structure haut-bas : \zh{罒} + \zh{非} \\
\hline
家 & \zh{宀} (mián, toit) & Toit, maison & 10 traits ; structure haut-bas : \zh{宀} + \zh{豕} (cochon = richesse) \\
\hline
原 & \zh{厂} (hǎn, falaise) & Falaise, rocher & 10 traits ; structure semi-encadrement : \zh{厂} + \zh{泉} \\
\hline
交 & \zh{亠} (tóu, couvercle) & Couvercle, haut & 6 traits ; structure haut-bas : \zh{亠} + \zh{父} (simplifié) \\
\hline
帮 & \zh{巾} (jīn, tissu) & Tissu, serviette & 12 traits ; structure haut-bas : \zh{邦} + \zh{巾} \\
\hline
\end{tabularx}
\end{center}

\textbf{Ordre des traits détaillé (règles : haut \textrightarrow{} bas, gauche \textrightarrow{} droite, horizontal \textrightarrow{} vertical, extérieur \textrightarrow{} intérieur, fermeture finale).}
\begin{itemize}
\item \zh{家} (10) : 1. point haut \zh{丶}, 2. trait horizontal court \zh{一}, 3. crochet vertical \zh{乛} (clé \zh{宀} finie). 4. horizontal long \zh{一}, 5. vertical \zh{丨}, 6. horizontal \zh{一}, 7. crochet coudé \zh{乛}, 8. horizontal \zh{一}, 9. trait oblique gauche \zh{丿}, 10. trait oblique droit \zh{乀} (partie \zh{豕}).
\item \zh{罪} (13) : Écrire \zh{罒} (5 traits : 2 verticaux, 3 horizontaux), puis \zh{非} (8 traits : haut \zh{飛} simplifié).
\item \zh{帮} (12) : Écrire \zh{邦} en haut (trois horizontaux, vertical, oreille droite \zh{阝}), puis \zh{巾} en bas (vertical, deux horizontaux, crochet final).
\end{itemize}
\end{propriete}

\courssub{Compréhension du texte}

\begin{exoresolu}[Compréhension globale du texte support]
\textbf{Énoncé.} Lis le texte de la boîte \emph{Texte support}, puis réponds en français : (a) De quel problème parle le texte ? (b) Quelles sont les deux causes mentionnées ? (c) Que propose l'auteur pour résoudre ce problème ?

\tcblower
\textbf{Solution détaillée.}
\begin{enumerate}
\item \textbf{Sujet général.} La première phrase pose le thème : \zh{青少年犯罪是一个严重的社会问题} « la délinquance juvénile est un grave problème social ». \textbf{Réponse (a) :} Le texte parle de la délinquance juvénile, présentée comme un problème social grave ces dernières années (\zh{近年来}).
\item \textbf{Causes.} Le marqueur \zh{原因有两个} « les causes sont au nombre de deux » annonce une énumération introduite par \zh{第一}\ldots\zh{第二}\ldots « premièrement\ldots deuxièmement\ldots ». \textbf{Réponse (b) :} (1) L'éducation familiale insuffisante (\zh{家庭教育不够}) ; (2) Les mauvaises fréquentations (\zh{交了不好的朋友}, litt. « ont fait de mauvais amis »).
\item \textbf{Solution.} La structure \zh{为了}\ldots\zh{应该}\ldots « afin de\ldots il faut\ldots » introduit la proposition. \textbf{Réponse (c) :} Pour résoudre ce problème, la famille (\zh{家庭}), l'école (\zh{学校}) et la société (\zh{社会}) doivent faire des efforts ensemble (\zh{一起努力}) pour aider (\zh{帮助}) ces jeunes.
\item \textbf{Points grammaticaux à justifier.} \zh{近年来} « ces dernières années » est un complément de temps placé en tête de phrase ; \zh{为了} + verbe « afin de » se place avant le sujet ; \zh{一起} « ensemble » se place entre le sujet et le verbe ; \zh{两个} (et non \zh{二个}) s'utilise devant un classificateur pour compter.
\end{enumerate}
\end{exoresolu}

\begin{methode}[Répondre aux questions de compréhension en chinois]
\begin{enumerate}
\item \textbf{Repère le mot interrogatif.} \zh{什么} « quoi », \zh{谁} « qui », \zh{为什么} « pourquoi », \zh{怎么样} « comment », \zh{几} « combien (petit nombre) ». La réponse remplace exactement ce mot à la même place dans la phrase.
\item \textbf{Localise la phrase du texte.} Cherche dans le texte les mots de la question ; la réponse se trouve dans la même phrase ou la suivante.
\item \textbf{Réponds par une phrase complète.} Recopie la structure de la question en remplaçant le mot interrogatif par l'information du texte. Ne réponds jamais par un seul mot à l'écrit.
\item \textbf{Pour \zh{为什么}, utilise \zh{因为}\ldots(\zh{所以})\ldots.} « Parce que\ldots (donc)\ldots ». \zh{因为} introduit la cause, \zh{所以} (facultatif) la conséquence.
\item \textbf{Relis ta réponse.} Vérifie l'ordre Sujet + Verbe + Complément, les tons du pinyin et les caractères.
\end{enumerate}
\end{methode}

\begin{exoresolu}[Questions de compréhension détaillée sur le texte]
\textbf{Énoncé.} Réponds en chinois par des phrases complètes, d'après le texte support :\\
1. \zh{青少年犯罪是一个什么样的问题？} \\
2. \zh{年轻人犯罪的原因有几个？} \\
3. \zh{为什么一些年轻人犯罪？} \\
4. \zh{谁应该帮助这些年轻人？}

\tcblower
\textbf{Solution détaillée.}
\begin{enumerate}
\item La question contient \zh{什么样的} « quel genre de ». Le texte dit \zh{严重的社会问题}. \textbf{Réponse :} \zh{青少年犯罪是一个严重的社会问题。} \emph{Qīngshàonián fànzuì shì yí ge yánzhòng de shèhuì wèntí.} « C'est un grave problème social. »
\item Le texte dit \zh{原因有两个}. \textbf{Réponse :} \zh{年轻人犯罪的原因有两个。} \emph{Niánqīngrén fànzuì de yuányīn yǒu liǎng ge.} « Il y a deux causes. » Devant un chiffre + classificateur, on utilise \zh{两} (et non \zh{二}) : \zh{两个}.
\item Question avec \zh{为什么} \textrightarrow{} réponse avec \zh{因为}. Le texte donne deux causes liées par \zh{第一}\ldots\zh{第二}. \textbf{Réponse :} \zh{因为家庭教育不够，而且他们交了不好的朋友。} \emph{Yīnwèi jiātíng jiàoyù bù gòu, érqiě tāmen jiāo le bù hǎo de péngyou.} « Parce que l'éducation familiale est insuffisante et qu'en plus ils se sont fait de mauvais amis. » \zh{而且} « de plus / qui plus est » relie les deux causes.
\item La question porte sur \zh{谁}. Le texte dit \zh{家庭、学校和社会}. \textbf{Réponse :} \zh{家庭、学校和社会应该帮助这些年轻人。} \emph{Jiātíng, xuéxiào hé shèhuì yīnggāi bāngzhù zhèxiē niánqīngrén.} « La famille, l'école et la société doivent aider ces jeunes. » Énumération avec la virgule chinoise \zh{、} et \zh{和} avant le dernier élément.
\end{enumerate}
\end{exoresolu}

\courssub{Production guidée et réinvestissement}

\begin{methode}[Mémoriser et réemployer le lexique de la délinquance juvénile]
\begin{enumerate}
\item \textbf{Apprends le mot avec son ton et sa nature.} Ne mémorise jamais un caractère seul : dis \cle{fànzuì} (\zh{犯罪}, verbe, « commettre un crime ») en marquant le 4\up{e} ton descendant sur les deux syllabes.
\item \textbf{Apprends les collocations (associations fixes).} Un mot chinois s'apprend avec ses partenaires : \zh{交朋友} « se faire des amis », \zh{解决问题} « résoudre un problème », \zh{犯罪的原因} « la cause du crime », \zh{帮助年轻人} « aider les jeunes », \zh{预防犯罪} « prévenir la délinquance ».
\item \textbf{Construis des familles de mots.} Autour d'un caractère-clé : \zh{青} (jeune) \textrightarrow{} \zh{青少年} « adolescent », \zh{年轻人} « jeune » ; \zh{法} (loi) \textrightarrow{} \zh{法律} « loi », \zh{违法} « violer la loi » ; \zh{教} (enseigner) \textrightarrow{} \zh{教育} « éducation », \zh{教师} « professeur ».
\item \textbf{Réemploie dans une phrase personnelle contextualisée.} Pour chaque mot, écris une phrase simple liée au Cameroun : \zh{喀麦隆的学校应该帮助学生。} (\emph{Kāmàilóng de xuéxiào yīnggāi bāngzhù xuéshēng.}) « Les écoles du Cameroun doivent aider les élèves. »
\item \textbf{Vérifie la prononciation à voix haute.} Lis ta phrase en respectant les tons ; un ton faux change le sens.
\end{enumerate}
\end{methode}

\begin{exoresolu}[Entraînement vocabulaire : compléter et employer]
\textbf{Énoncé.} Complète les phrases avec le mot qui convient parmi : \zh{犯罪 / 原因 / 解决 / 帮助 / 严重 / 违法 / 预防 / 关爱}.\\
1. 青少年\zh{……}是一个社会问题。 \\
2. 他\zh{……}的原因是什么？ \\
3. 我们应该\zh{……}这个问题。 \\
4. 老师\zh{……}了那个年轻人。 \\
5. 这个问题很\zh{……}。 \\
6. 社会应该\zh{……}青少年犯罪。 \\
7. 这些孩子需要更多的\zh{……}。

\tcblower
\textbf{Solution détaillée.}
\begin{enumerate}
\item 青少年\textbf{犯罪}是一个社会问题。 \emph{Qīngshàonián fànzuì shì yí ge shèhuì wèntí.} « La délinquance juvénile est un problème social. » \zh{犯罪} est ici le sujet (verbe substantivé), placé avant \zh{是}.
\item 他\textbf{违法}的\textbf{原因}是什么？ \emph{Tā wéifǎ de yuányīn shì shénme ?} « Quelle est la cause de sa violation de la loi ? » \zh{原因} « cause » est précédé du complément du nom marqué par \zh{的}.
\item 我们应该\textbf{解决}这个问题。 \emph{Wǒmen yīnggāi jiějué zhège wèntí.} « Nous devons résoudre ce problème. » \zh{解决} se place après le modal \zh{应该} et avant l'objet \zh{问题}.
\item 老师\textbf{帮助}了那个年轻人。 \emph{Lǎoshī bāngzhù le nà ge niánqīngrén.} « Le professeur a aidé ce jeune. » \zh{了} marque l'accomplissement ; classificateur \zh{个} devant \zh{年轻人}.
\item 这个问题很\textbf{严重}。 \emph{Zhège wèntí hěn yánzhòng.} « Ce problème est très grave. » Adjectif attribut précédé de \zh{很}, sans verbe « être ».
\item 社会应该\textbf{预防}青少年犯罪。 \emph{Shèhuì yīnggāi yùfáng qīngshàonián fànzuì.} « La société doit prévenir la délinquance juvénile. » \zh{预防} prend un objet direct (le risque).
\item 这些孩子需要更多的\textbf{关爱}。 \emph{Zhèxiē háizi xūyào gèng duō de guān'ài.} « Ces enfants ont besoin de plus d'affection/soins. » \zh{关爱} nom « affection, soins ».
\end{enumerate}
\textbf{Conclusion.} Vérifie toujours la nature du mot (verbe, nom, adjectif) pour le placer correctement dans la phrase.
\end{exoresolu}

\begin{methode}[Rédiger 4 à 6 phrases sur un problème social (Modèle « Problème - Causes - Solutions - Espoir »)]
\begin{enumerate}
\item \textbf{Phrase 1 : Présenter le problème.} Modèle : \zh{……是一个严重的问题。} « \ldots est un grave problème. » (Utiliser \zh{也} « aussi » pour comparer).
\item \textbf{Phrase 2 : Donner les causes.} Modèle : \zh{原因有两个：第一，……；第二，……。} ou \zh{因为……，所以……。} « Parce que\ldots, donc\ldots »
\item \textbf{Phrase 3 : Proposer des solutions.} Modèle : \zh{为了解决这个问题，……应该……。} « Pour résoudre ce problème, \ldots doit\ldots » (Sujet : \zh{家庭、学校、社会、政府}).
\item \textbf{Phrase 4 : Action concrète.} Modèle : \zh{他们应该给年轻人更多的教育和关爱。} « Ils doivent donner plus d'éducation et d'affection aux jeunes. » (Structure \zh{给} A B).
\item \textbf{Phrase 5 : Conclure avec un espoir.} Modèle : \zh{我相信……会越来越好。} « Je crois que\ldots ira de mieux en mieux. » (\zh{越来越} + adj.).
\item \textbf{Vérifie.} Une seule idée par phrase, ordre S + V + C, tons corrects, caractères simplifiés, ponctuation chinoise.
\end{enumerate}
\end{methode}

\begin{exoresolu}[Rédaction guidée : La délinquance juvénile au Cameroun]
\textbf{Énoncé.} En t'aidant du texte support, du glossaire et de la méthode, rédige un court paragraphe (5 phrases) en chinois sur la délinquance juvénile dans ton pays (Cameroun), avec pinyin et traduction.

\tcblower
\textbf{Solution détaillée (Modèle de production attendue).}

在喀麦隆，青少年犯罪也是一个严重的社会问题。原因有两个：第一，一些家庭教育不够，父母没时间管孩子；第二，一些年轻人不上学，交了不好的朋友。为了解决这个问题，家庭、学校和政府应该一起努力。他们应该给这些年轻人更多的教育和关爱，帮助他们回归社会。我相信喀麦隆的青少年会越来越好的。\\
\emph{Zài Kāmàilóng, qīngshàonián fànzuì yě shì yí ge yánzhòng de shèhuì wèntí. Yuányīn yǒu liǎng ge: dì-yī, yìxiē jiātíng jiàoyù bù gòu, fùmǔ méi shíjiān guǎn háizi; dì-èr, yìxiē niánqīngrén bú shàngxué, jiāo le bù hǎo de péngyou. Wèile jiějué zhège wèntí, jiātíng, xuéxiào hé zhèngfǔ yīnggāi yìqǐ nǔlì. Tāmen yīnggāi gěi zhèxiē niánqīngrén gèng duō de jiàoyù hé guān'ài, bāngzhù tāmen huíguī shèhuì. Wǒ xiāngxìn Kāmàilóng de qīngshàonián huì yuè lái yuè hǎo de.}\\
« Au Cameroun, la délinquance juvénile est aussi un grave problème social. Il y a deux causes : premièrement, l'éducation familiale de certaines familles est insuffisante, les parents n'ont pas le temps de s'occuper des enfants ; deuxièmement, certains jeunes ne vont pas à l'école et se font de mauvais amis. Pour résoudre ce problème, la famille, l'école et le gouvernement doivent faire des efforts ensemble. Ils doivent donner plus d'éducation et d'affection à ces jeunes, les aider à se réinsérer dans la société. Je crois que les jeunes du Cameroun iront de mieux en mieux. »

\textbf{Justifications grammaticales.}
\begin{enumerate}
\item \zh{也} « aussi » se place avant le verbe \zh{是}.
\item \zh{没时间} (méi shíjián) « ne pas avoir le temps » : négation de \zh{有} par \zh{没}.
\item \zh{管} (guǎn) « s'occuper de, gérer, contrôler » + objet \zh{孩子}.
\item \zh{回归社会} (huíguī shèhuì) « réintégrer la société / retourner dans la société ».
\item \zh{越来越好} (yuè lái yuè hǎo) « de mieux en mieux » : structure comparative de progression.
\end{enumerate}
\end{exoresolu}

\begin{attention}[Les 5 erreurs qui coûtent des points au Bac / Évaluation]
\begin{enumerate}
\item \textbf{Tons négligés.} Écrire \emph{fanzui} sans tons ou avec de mauvais tons : \zh{犯罪} est \emph{fànzuì} (4\up{e} + 4\up{e} ton). Un ton faux change le sens (ex: \zh{犯罪} vs \zh{饭最} « le riz le plus... ») et fait perdre des points en expression écrite comme orale.
\item \textbf{Caractères confondus (homophones / graphies proches).} \zh{帮} (aider, clé \zh{巾}) vs \zh{邦} (pays, clé \zh{阝}) ; \zh{原} (origine, clé \zh{厂}) vs \zh{愿} (souhaiter, clé \zh{心}) ; \zh{罪} (crime, clé \zh{罒}) vs \zh{非} (ne pas être, clé \zh{非}) ; \zh{交} (échanger) vs \zh{较} (comparer). Vérifie la clé et le nombre de traits avant d'écrire.
\item \textbf{Ordre des mots calqué du français.} On n'écrit pas \zh{问题严重是一个} : l'ordre est Sujet + \zh{是} + Complément : \zh{青少年犯罪是一个严重的问题。} Les compléments de temps (\zh{近年来}) et de lieu (\zh{在喀麦隆}) se placent \textbf{avant} le verbe, jamais à la fin comme en français.
\item \textbf{Oubli ou mauvais usage de \zh{了} et \zh{的}.} \zh{了} marque l'action accomplie / aspect parfaitif : \zh{他交了不好的朋友} (il s'est fait de mauvais amis). \zh{的} relie le nom à son complément : \zh{犯罪的原因} « la cause du crime ». Ne mets \textbf{pas} \zh{了} après un état permanent (\zh{他很高了} \textrightarrow{} FAUX, dire \zh{他很高}).
\item \textbf{Classificateur oublié ou faux.} On dit \zh{一个问题}, \zh{两个原因}, \zh{这些年轻人} : le classificateur \zh{个} est obligatoire entre le nombre / démonstratif et le nom. Devant un nombre, utilise \zh{两} (\zh{两个}) et non \zh{二} (\zh{二个}).
\end{enumerate}
\end{attention}

\begin{savaistu}[Regard croisé Cameroun -- Chine : Prévention de la délinquance]
\begin{itemize}
\item \textbf{Cameroun :} Le Code pénal protège les mineurs (peines atténuées, centres de rééducation comme le Centre de Rééducation de Bétamba). Rôle crucial des \emph{Conseils de la Jeunesse} et des associations (ex: \emph{RECAJ}) pour l'encadrement. Problème des \emph{enfants de la rue} (Yaoundé, Douala) lié à la pauvreté et aux conflits (régions NOSO).
\item \textbf{Chine :} Loi sur la protection des mineurs (2021 révisée) et Loi sur la prévention de la délinquance juvénile (2020). Système \emph{École-Famille-Société} (\zh{家校社}) très structuré : \zh{法治副校长} (vice-principaux juridiques, souvent policiers/procureurs) dans chaque école ; \zh{专门学校} (écoles spécialisées de correction) pour les cas graves ; surveillance numérique des cybercafés/jeux vidéo (limitation temps pour mineurs).
\item \textbf{Point commun :} La responsabilité partagée \zh{家庭、学校、社会共同责任} est le principe central dans les deux pays. La prévention par l'éducation (\zh{预防为主}) prime sur la répression.
\end{itemize}
\end{savaistu}

\begin{exercice}[Entraînement autonome -- Devoir / Travail en classe]
\begin{enumerate}
\item \textbf{Traduction (Version).} Traduis en français :\\
\zh{近年来，我们城市的青少年犯罪率上升了。因为家庭教育缺失，而且网络游戏影响大。政府决定建立更多的专门学校，帮助违法青少年改正错误。}
\item \textbf{Traduction (Thème).} Traduis en chinois (caractères + pinyin) :\\
« Ces dernières années, le taux de délinquance juvénile dans notre ville a augmenté. Parce que l'éducation familiale fait défaut et que l'influence des jeux en ligne est grande. Le gouvernement a décidé de créer plus d'écoles spécialisées pour aider les jeunes délinquants à corriger leurs erreurs. »
\item \textbf{Expression écrite.} Rédige un message WeChat / WhatsApp (50-80 caractères) à un ami chinois pour lui expliquer une action de ton lycée (club, conférence, bénévolat) qui aide les jeunes en difficulté. Utilise : \zh{参加}, \zh{活动}, \zh{帮助}, \zh{有意义}.
\item \textbf{Caractères à écrire.} Recopie 5 fois chaque mot en respectant l'ordre des traits : \zh{犯罪}, \zh{帮助}, \zh{预防}, \zh{关爱}, \zh{违法}.
\end{enumerate}
\end{exercice}

\begin{corrige}[Exercice 1 -- Version]
\textbf{Traduction :} « Ces dernières années, le taux de délinquance juvénile dans notre ville a augmenté. Parce que l'éducation familiale fait défaut (est absente), et que l'influence des jeux en ligne est grande. Le gouvernement a décidé d'établir plus d'écoles spécialisées, pour aider les jeunes délinquants à corriger leurs erreurs. »
\textit{Notes :} \zh{犯罪率} (taux de criminalité), \zh{上升} (augmenter), \zh{缺失} (faire défaut/manquer), \zh{网络游戏} (jeux en ligne), \zh{专门学校} (écoles spécialisées de correction), \zh{改正错误} (corriger ses erreurs).
\end{corrige}

\begin{corrige}[Exercice 2 -- Thème]
\zh{近年来，我们城市的青少年犯罪率上升了。因为家庭教育缺失，而且网络游戏影响大。政府决定建立更多的专门学校，帮助违法青少年改正错误。}\\
\emph{Jìnnián lái, wǒmen chéngshì de qīngshàonián fànzuì lǜ shàngshēng le. Yīnwèi jiātíng jiàoyù quēshī, érqiě wǎngluò yóuxì yǐngxiǎng dà. Zhèngfǔ juédìng jiànlì gèng duō de zhuānmén xuéxiào, bāngzhù wéifǎ qīngshàonián gǎizhèng cuòwù.}
\end{corrige}

\begin{corrige}[Exercice 3 -- Expression écrite (Modèle)]
\zh{亲爱的小李：\\
我们学校参加了一个“阳光青年”活动，帮助困难青少年学习技能，很有意义。\\
你有空来看看吗？\\
王明}\\
\emph{Qīn'ài de Xiǎo Lǐ: Wǒmen xuéxiào cānjiā le yí ge "Yángguāng Qīngnián" huódòng, bāngzhù kùnnán qīngnián xuéxí jìnéng, hěn yǒu yìyì. Nǐ yǒukòng lái kànkan ma? Wáng Míng.}
\end{corrige}

\newpage

\coursec{Exercices}

\courssub{Je vérifie mes connaissances}

\begin{exercice}[QCM : le bon sens]
Entoure la bonne réponse.

\begin{enumerate}
\item \zh{犯罪} (fànzuì) signifie :
\begin{itemize}
\item[a)] respecter la loi
\item[b)] commettre un crime
\item[c)] protéger les enfants
\end{itemize}
\item \zh{违法} (wéifǎ) signifie :
\begin{itemize}
\item[a)] enfreindre la loi
\item[b)] étudier le droit
\item[c)] faire ses devoirs
\end{itemize}
\item \zh{惩罚} (chéngfá) signifie :
\begin{itemize}
\item[a)] récompenser
\item[b)] punir
\item[c)] encourager
\end{itemize}
\item \zh{青少年} (qīngshàonián) désigne :
\begin{itemize}
\item[a)] les personnes âgées
\item[b)] les adolescents et les jeunes
\item[c)] les enseignants
\end{itemize}
\end{enumerate}
\end{exercice}

\begin{exercice}[Vrai ou faux ? Justifie]
Réponds par VRAI ou FAUX, puis justifie ta réponse en une phrase en français.

\begin{enumerate}
\item \zh{保护} (bǎohù) veut dire « détruire ».
\item \zh{影响} (yǐngxiǎng) peut être un verbe ou un nom.
\item \zh{解决} (jiějué) s'emploie souvent avec \zh{问题} (wèntí, problème).
\item \zh{教育} (jiàoyù) signifie « punition ».
\end{enumerate}
\end{exercice}

\begin{exercice}[Le mot qui convient]
Complète chaque phrase avec le mot qui convient : \zh{犯罪}, \zh{违法}, \zh{保护}, \zh{惩罚}, \zh{影响}, \zh{解决}.

\begin{enumerate}
\item 偷东西是\zh{＿＿＿＿}的行为。(Tōu dōngxi shì ＿＿＿＿ de xíngwéi.)
\item 父母应该\zh{＿＿＿＿}自己的孩子。(Fùmǔ yīnggāi ＿＿＿＿ zìjǐ de háizi.)
\item 我们要想办法\zh{＿＿＿＿}这个问题。(Wǒmen yào xiǎng bànfǎ ＿＿＿＿ zhège wèntí.)
\item 法律\zh{＿＿＿＿}犯罪的人。(Fǎlǜ ＿＿＿＿ fànzuì de rén.)
\item 坏朋友对青少年有不好的\zh{＿＿＿＿}。(Huài péngyou duì qīngshàonián yǒu bù hǎo de ＿＿＿＿.)
\end{enumerate}
\end{exercice}

\begin{exercice}[Associe le caractère à son pinyin]
Relie chaque caractère à la bonne transcription pinyin.

\begin{center}
\begin{tabularx}{\linewidth}{|X|X|}
\hline
\rowcolor{popblueL} Caractère & Pinyin \\
\hline
1. 罪 & a) jiào \\
\hline
2. 法 & b) hù \\
\hline
3. 犯 & c) zuì \\
\hline
4. 护 & d) fǎ \\
\hline
5. 教 & e) fàn \\
\hline
\end{tabularx}
\end{center}
\end{exercice}

\courssub{Je m'entraîne}

\begin{exercice}[Traduction en français]
Traduis les phrases suivantes en français.

\begin{enumerate}
\item \zh{有的青少年不遵守法律。} (Yǒude qīngshàonián bù zūnshǒu fǎlǜ.)
\item \zh{家庭和学校应该一起保护青少年。} (Jiātíng hé xuéxiào yīnggāi yìqǐ bǎohù qīngshàonián.)
\item \zh{犯罪会影响一个人的未来。} (Fànzuì huì yǐngxiǎng yí gè rén de wèilái.)
\item \zh{政府正在努力解决这个问题。} (Zhèngfǔ zhèngzài nǔlì jiějué zhège wèntí.)
\end{enumerate}
\end{exercice}

\begin{exercice}[Transformation avec 有的……有的……还有的……]
À partir de la liste d'actes suivante : \zh{偷东西} (voler), \zh{打架} (se battre), \zh{逃学} (sécher les cours), écris une phrase complète avec la structure \zh{有的……有的……还有的……}.

Modèle : \zh{有的学生喜欢足球，有的喜欢篮球，还有的喜欢音乐。}
\end{exercice}

\begin{exercice}[Texte à trous : mots-outils]
Complète le texte suivant avec les mots-outils qui conviennent : \zh{的}, \zh{了}, \zh{得}, \zh{在}.

在雅温得，有一些青少年走\zh{＿＿}了错误\zh{＿＿}路。他们\zh{＿＿}街上偷东西，做\zh{＿＿}很不好。政府\zh{＿＿}去年采取\zh{＿＿}新措施，帮助这些年轻人回到学校。

\emph{Indice : il y a six trous ; deux mots-outils sont utilisés deux fois.}
\end{exercice}

\begin{exercice}[Remets les mots en ordre]
Remets les mots dans le bon ordre pour former des phrases correctes, puis traduis-les en français.

\begin{enumerate}
\item 应该 / 关心 / 父母 / 孩子 / 自己的
\item 青少年 / 法律 / 惩罚 / 犯罪的 / 会
\item 一起 / 学校 / 和 / 家庭 / 解决问题 / 应该
\item 影响 / 坏朋友 / 青少年 / 有 / 对 / 不好的
\end{enumerate}
\end{exercice}

\courssub{Je me prépare au Bac}

\begin{exercice}[Compréhension écrite type Bac]
Lis le texte suivant, puis réponds aux questions \textbf{en chinois, par des phrases complètes}.

\begin{textebox}[青少年犯罪的问题]
近年来，青少年犯罪成了很多城市的大问题。\\ (Jìnnián lái, qīngshàonián fànzuì chéngle hěn duō chéngshì de dà wèntí.)\\
有的青少年逃学，有的偷东西，还有的打架。\\ (Yǒude qīngshàonián táoxué, yǒude tōu dōngxi, hái yǒude dǎjià.)\\
专家认为，这个问题有几个原因：有的家庭不关心孩子，\\ (Zhuānjiā rènwéi, zhège wèntí yǒu jǐ ge yuányīn: yǒude jiātíng bù guānxīn háizi,)\\
有的学校不重视教育，还有的年轻人受坏朋友的影响。\\ (yǒude xuéxiào bù zhòngshì jiàoyù, hái yǒude niánqīngrén shòu huài péngyou de yǐngxiǎng.)\\
为了解决这个问题，政府、学校和家庭应该一起努力，\\ (Wèile jiějué zhège wèntí, zhèngfǔ, xuéxiào hé jiātíng yīnggāi yìqǐ nǔlì,)\\
保护青少年，帮助他们走上正确的路。\\ (bǎohù qīngshàonián, bāngzhù tāmen zǒushàng zhèngquè de lù.)\\
\emph{Traduction : Ces dernières années, la délinquance juvénile est devenue un grand problème dans beaucoup de villes. Certains adolescents sèchent les cours, d'autres volent des choses, d'autres encore se battent. Selon les experts, ce problème a plusieurs causes : certaines familles ne se soucient pas de leurs enfants, certaines écoles n'accordent pas d'importance à l'éducation, et certains jeunes subissent l'influence de mauvais amis. Pour résoudre ce problème, le gouvernement, l'école et la famille doivent travailler ensemble, protéger les adolescents et les aider à prendre le bon chemin.}
\end{textebox}

\begin{enumerate}
\item 什么问题成了很多城市的大问题？
\item 青少年做了哪些违法的事？（举三个例子）
\item 专家认为这个问题有几个原因？是什么？
\item 为了解决这个问题，谁应该一起努力？
\item Vrai ou faux ? Justifie en français à partir du texte : « Selon le texte, seule la famille est responsable de la délinquance juvénile. »
\end{enumerate}
\end{exercice}

\begin{exercice}[Thème et version type Bac]
\textbf{A. Thème.} Traduis les phrases suivantes en chinois (caractères simplifiés ; le pinyin n'est pas exigé mais peut être ajouté).

\begin{enumerate}
\item Les parents doivent se soucier de leurs enfants, et l'école doit renforcer l'éducation.
\item Voler est un acte illégal : la loi punit ceux qui commettent des crimes.
\end{enumerate}

\textbf{B. Version.} Traduis en français le paragraphe suivant.

\begin{textebox}[Version]
青少年犯罪的原因很多。\\ (Qīngshàonián fànzuì de yuányīn hěn duō.)\\
有的孩子没有父母的关心，有的年轻人找不到工作，\\ (Yǒude háizi méiyǒu fùmǔ de guānxīn, yǒude niánqīngrén zhǎobudào gōngzuò,)\\
还有的学生不喜欢学习，常常逃学。\\ (hái yǒude xuésheng bù xǐhuan xuéxí, chángcháng táoxué.)\\
我们应该了解这些原因，才能找到好的解决办法。\\ (Wǒmen yīnggāi liǎojiě zhèxiē yuányīn, cái néng zhǎodào hǎo de jiějué bànfǎ.)
\end{textebox}
\end{exercice}

\begin{exercice}[Expression écrite type Bac]
\textbf{Sujet :} \zh{家庭在预防青少年犯罪中的作用} — « Le rôle de la famille dans la prévention de la délinquance juvénile au Cameroun ».

Rédige un court texte en chinois de \textbf{6 à 8 phrases} (environ 60 à 90 caractères). Tu peux t'appuyer sur les pistes suivantes :

\begin{itemize}
\item la situation : la délinquance juvénile existe aussi dans les villes du Cameroun (Yaoundé, Douala) ;
\item le rôle des parents : se soucier de leurs enfants, leur parler, connaître leurs amis ;
\item les moyens : donner une bonne éducation, aider les enfants à rester à l'école ;
\item ta conclusion : famille et école doivent travailler ensemble.
\end{itemize}

\textbf{Contraintes :} utilise au moins une fois la structure \zh{有的……有的……}, le verbe \zh{应该} et deux mots du vocabulaire de la leçon (\zh{保护}, \zh{关心}, \zh{教育}, \zh{预防}…).
\end{exercice}

\newpage

\coursec{Corrigés des exercices}

\begin{corrige}[Exercice 1 — QCM : le bon sens]
\begin{enumerate}
\item \textbf{b)} \zh{犯罪} (fànzuì) signifie « commettre un crime ». \zh{犯} = commettre, enfreindre ; \zh{罪} = crime, péché.
\item \textbf{a)} \zh{违法} (wéifǎ) signifie « enfreindre la loi ». \zh{违} = violer, aller contre ; \zh{法} = loi.
\item \textbf{b)} \zh{惩罚} (chéngfá) signifie « punir ». Attention aux tons : chéng (2\up{e} ton), fá (2\up{e} ton).
\item \textbf{b)} \zh{青少年} (qīngshàonián) désigne les adolescents et les jeunes. \zh{青} évoque la jeunesse, \zh{少} le jeune âge, \zh{年} l'année/l'âge.
\end{enumerate}
\textbf{Point de vigilance :} ne pas confondre \zh{犯罪} (verbe : commettre un crime) et \zh{罪} (nom : crime). On dit \zh{他犯罪了} (tā fànzuì le, il a commis un crime) mais \zh{这是罪} (zhè shì zuì, c'est un crime).
\end{corrige}

\begin{corrige}[Exercice 2 — Vrai ou faux ? Justifie]
\begin{enumerate}
\item \textbf{FAUX.} \zh{保护} (bǎohù) veut dire « protéger », et non « détruire ». C'est le contraire : on protège les enfants, on ne les détruit pas.
\item \textbf{VRAI.} \zh{影响} (yǐngxiǎng) peut être un verbe (\zh{影响未来}, influencer l'avenir) ou un nom (\zh{不好的影响}, une mauvaise influence).
\item \textbf{VRAI.} \zh{解决} (jiějué) s'emploie très souvent avec \zh{问题} : \zh{解决问题} (jiějué wèntí) = « résoudre un problème ». C'est une collocation à mémoriser comme un bloc.
\item \textbf{FAUX.} \zh{教育} (jiàoyù) signifie « éducation » ; « punition » se dit \zh{惩罚} (chéngfá). Ne pas confondre les deux notions.
\end{enumerate}
\end{corrige}

\begin{corrige}[Exercice 3 — Le mot qui convient]
\begin{enumerate}
\item 偷东西是\zh{违法}的行为。(Tōu dōngxi shì wéifǎ de xíngwéi.) « Voler des choses est un acte illégal. » — \zh{违法} est ici adjectif devant \zh{的行为}.
\item 父母应该\zh{保护}自己的孩子。(Fùmǔ yīnggāi bǎohù zìjǐ de háizi.) « Les parents doivent protéger leurs enfants. »
\item 我们要想办法\zh{解决}这个问题。(Wǒmen yào xiǎng bànfǎ jiějué zhège wèntí.) « Nous devons trouver un moyen de résoudre ce problème. »
\item 法律\zh{惩罚}犯罪的人。(Fǎlǜ chéngfá fànzuì de rén.) « La loi punit ceux qui commettent des crimes. »
\item 坏朋友对青少年有不好的\zh{影响}。(Huài péngyou duì qīngshàonián yǒu bù hǎo de yǐngxiǎng.) « Les mauvais amis ont une mauvaise influence sur les adolescents. »
\end{enumerate}
\textbf{Point de vigilance :} \zh{犯罪} n'était pas la bonne réponse à la question 1 : \zh{犯罪的行为} serait redondant (« un acte de commettre-un-crime ») ; on qualifie l'acte avec \zh{违法}.
\end{corrige}

\begin{corrige}[Exercice 4 — Associe le caractère à son pinyin]
\begin{center}
\begin{tabularx}{\linewidth}{|X|X|X|}
\hline
\rowcolor{popblueL} Caractère & Pinyin & Vérification \\
\hline
1. 罪 & c) zuì & 4\up{e} ton, initiale z \\
\hline
2. 法 & d) fǎ & 3\up{e} ton \\
\hline
3. 犯 & e) fàn & 4\up{e} ton, finale -an \\
\hline
4. 护 & b) hù & 4\up{e} ton, initiale h \\
\hline
5. 教 & a) jiào & 4\up{e} ton (dans \zh{教育} jiàoyù) \\
\hline
\end{tabularx}
\end{center}
\textbf{Point de vigilance :} ne pas confondre \zh{犯} (fàn) et \zh{法} (fǎ) : même initiale f, mais finales et tons différents. Le caractère \zh{教} a deux lectures : jiāo (enseigner, \zh{教书}) et jiào (éducation, \zh{教育}) ; ici c'est jiào.
\end{corrige}

\begin{corrige}[Exercice 5 — Traduction en français]
\begin{enumerate}
\item \zh{有的青少年不遵守法律。} « Certains adolescents ne respectent pas la loi. » — \zh{有的} = certains ; \zh{遵守} = respecter, observer.
\item \zh{家庭和学校应该一起保护青少年。} « La famille et l'école doivent protéger ensemble les adolescents. » — \zh{一起} (ensemble) se place avant le verbe.
\item \zh{犯罪会影响一个人的未来。} « Commettre un crime peut affecter l'avenir d'une personne. » — \zh{会} marque ici la probabilité/conséquence ; \zh{一个人} = une personne (classificateur \zh{个}).
\item \zh{政府正在努力解决这个问题。} « Le gouvernement est en train de s'efforcer de résoudre ce problème. » — \zh{正在} marque l'action en cours (présent progressif).
\end{enumerate}
\textbf{Point de vigilance :} en chinois, le sujet de la phrase 3 est \zh{犯罪} (verbe substantivé) ; en français on peut traduire par un infinitif ou un nom (« la délinquance »).
\end{corrige}

\begin{corrige}[Exercice 6 — Transformation avec 有的……有的……还有的……]
\textbf{Réponse attendue (modèle) :}

\zh{有的青少年偷东西，有的打架，还有的逃学。}

\emph{(Yǒude qīngshàonián tōu dōngxi, yǒude dǎjià, hái yǒude táoxué.)}

« Certains adolescents volent des choses, d'autres se battent, d'autres encore sèchent les cours. »

\textbf{Justification grammaticale :} la structure \zh{有的……有的……还有的……} énumère des cas différents à l'intérieur d'un même groupe. Chaque segment reprend un sujet (\zh{有的青少年}, puis \zh{有的} sous-entendu) suivi d'un verbe. \zh{还有的} introduit le dernier élément de la liste. L'ordre des trois actes peut varier ; toute permutation est acceptée.
\end{corrige}

\begin{corrige}[Exercice 7 — Texte à trous : mots-outils]
Texte complété :

在雅温得，有一些青少年走\zh{上}了错误\zh{的}路。他们\zh{在}街上偷东西，做\zh{得}很不好。政府\zh{在}去年采取\zh{了}新措施，帮助这些年轻人回到学校。

\emph{(Zài Yǎwēndé, yǒu yìxiē qīngshàonián zǒushàng le cuòwù de lù. Tāmen zài jiēshang tōu dōngxi, zuò de hěn bù hǎo. Zhèngfǔ zài qùnián cǎiqǔ le xīn cuòshī, bāngzhù zhèxiē niánqīngrén huí dào xuéxiào.)}

« À Yaoundé, certains adolescents ont pris un mauvais chemin. Ils volent des choses dans la rue, ils agissent très mal. Le gouvernement a pris l'année dernière de nouvelles mesures pour aider ces jeunes à retourner à l'école. »

\textbf{Justifications :}
\begin{itemize}
\item trou 1 : \zh{了} (marque de l'accompli après \zh{走上}) — \emph{remarque : la locution \zh{走上了……的路} est correcte ; le mot-outil exigé est \zh{了}} ;
\item trou 2 : \zh{的} (lien entre l'adjectif \zh{错误} et le nom \zh{路}) ;
\item trou 3 : \zh{在} (complément de lieu \zh{在街上} placé avant le verbe) ;
\item trou 4 : \zh{得} (complément de degré : \zh{做得很不好}, « agir très mal ») ;
\item trou 5 : \zh{在} (complément de temps \zh{在去年}) ;
\item trou 6 : \zh{了} (action accomplie : \zh{采取了新措施}).
\end{itemize}
\textbf{Point de vigilance :} \zh{的}, \zh{得}, \zh{地} se prononcent tous « de » mais ne s'emploient pas de la même façon : \zh{的} + nom, \zh{得} + complément de manière après le verbe, \zh{地} + verbe.
\end{corrige}

\begin{corrige}[Exercice 8 — Remets les mots en ordre]
\begin{enumerate}
\item \zh{父母应该关心自己的孩子。} (Fùmǔ yīnggāi guānxīn zìjǐ de háizi.) « Les parents doivent se soucier de leurs enfants. » — ordre : sujet + auxiliaire \zh{应该} + verbe + complément.
\item \zh{法律会惩罚犯罪的青少年。} (Fǎlǜ huì chéngfá fànzuì de qīngshàonián.) « La loi punira les adolescents qui commettent des crimes. » — \zh{犯罪的} qualifie \zh{青少年}.
\item \zh{家庭和学校应该一起解决问题。} (Jiātíng hé xuéxiào yīnggāi yìqǐ jiějué wèntí.) « La famille et l'école doivent résoudre le problème ensemble. » — \zh{一起} se place après \zh{应该}, avant le verbe.
\item \zh{坏朋友对青少年有不好的影响。} (Huài péngyou duì qīngshàonián yǒu bù hǎo de yǐngxiǎng.) « Les mauvais amis ont une mauvaise influence sur les adolescents. » — structure : A + \zh{对} + B + \zh{有} + influence.
\end{enumerate}
\textbf{Point de vigilance :} le complément introduit par \zh{对} se place \emph{avant} le verbe, contrairement au français (« influence \emph{sur} les adolescents »).
\end{corrige}

\begin{corrige}[Exercice 9 — Compréhension écrite type Bac]
\textbf{Réponses rédigées en chinois :}
\begin{enumerate}
\item \zh{青少年犯罪成了很多城市的大问题。} (Qīngshàonián fànzuì chéngle hěn duō chéngshì de dà wèntí.) « La délinquance juvénile est devenue un grand problème dans beaucoup de villes. »
\item \zh{有的青少年逃学，有的偷东西，还有的打架。} « Certains sèchent les cours, d'autres volent des choses, d'autres encore se battent. » (trois exemples exigés)
\item \zh{专家认为这个问题有几个原因：有的家庭不关心孩子，有的学校不重视教育，还有的年轻人受坏朋友的影响。} « Selon les experts, il y a plusieurs causes : certaines familles ne se soucient pas de leurs enfants, certaines écoles négligent l'éducation, et certains jeunes subissent l'influence de mauvais amis. »
\item \zh{为了解决这个问题，政府、学校和家庭应该一起努力。} « Pour résoudre ce problème, le gouvernement, l'école et la famille doivent travailler ensemble. »
\item \textbf{FAUX.} Selon le texte, la responsabilité est partagée : le texte cite trois causes (la famille, l'école, les mauvais amis) et demande au gouvernement, à l'école et à la famille d'agir ensemble. La famille n'est donc pas seule responsable.
\end{enumerate}
\textbf{Barème indicatif (sur 10) :} 2 points par question (1 point pour l'information exacte, 1 point pour la correction de la phrase chinoise). Toute réponse en français aux questions 1 à 4 est pénalisée : la consigne exige le chinois.
\textbf{Points de vigilance :} recopier les structures du texte est autorisé et conseillé ; vérifier les tons de \zh{逃学} (táoxué), \zh{重视} (zhòngshì), \zh{受……的影响} (shòu … de yǐngxiǎng, « subir l'influence de »).
\end{corrige}

\begin{corrige}[Exercice 10 — Thème et version type Bac]
\textbf{A. Thème — propositions de traduction :}
\begin{enumerate}
\item \zh{父母应该关心自己的孩子，学校应该加强教育。} (Fùmǔ yīnggāi guānxīn zìjǐ de háizi, xuéxiào yīnggāi jiāqiáng jiàoyù.) — \zh{加强} (jiāqiáng) = renforcer ; \zh{应该} traduit « doivent ».
\item \zh{偷东西是违法的行为：法律惩罚犯罪的人。} (Tōu dōngxi shì wéifǎ de xíngwéi: fǎlǜ chéngfá fànzuì de rén.) — « ceux qui commettent des crimes » = \zh{犯罪的人} (proposition relative placée avant le nom avec \zh{的}).
\end{enumerate}
\textbf{B. Version — traduction :}

« Les causes de la délinquance juvénile sont nombreuses. Certains enfants ne reçoivent pas l'affection de leurs parents, certains jeunes ne trouvent pas de travail, et certains élèves n'aiment pas étudier et sèchent souvent les cours. Nous devons comprendre ces causes pour pouvoir trouver de bonnes solutions. »

\textbf{Notes de version :} \zh{找不到} (zhǎobudào) = « ne pas arriver à trouver » (complément de résultat négatif) ; \zh{才能} (cái néng) = « pour pouvoir (seulement alors) », à rendre par « pour pouvoir / afin de pouvoir » ; \zh{解决办法} = « solution, moyen de résoudre ».

\textbf{Barème indicatif (sur 10) :} thème 5 points (2,5 par phrase : vocabulaire 1, grammaire 1, caractères 0,5) ; version 5 points (fidélité 3, qualité du français 2).
\textbf{Points de vigilance :} ne pas écrire \zh{犯罪的人的行为} (redondance) ; respecter l'ordre sujet + \zh{应该} + verbe ; en version, ne pas traduire \zh{原因很多} mot à mot par « les raisons beaucoup ».
\end{corrige}

\begin{corrige}[Exercice 11 — Expression écrite type Bac]
\textbf{Proposition de rédaction modèle (8 phrases, environ 85 caractères) :}

\zh{在喀麦隆的城市，比如雅温得和杜阿拉，也有青少年犯罪的问题。有的年轻人逃学，有的偷东西。家庭在预防犯罪中很重要。父母应该关心孩子，常常和他们说话，也应该认识他们的朋友。父母还要给孩子好的教育，帮助他们留在学校。家庭可以保护孩子，让他们不受坏朋友的影响。我认为，家庭和学校应该一起努力。这样，青少年才能走上正确的路。}

\emph{(Zài Kāmàilóng de chéngshì, bǐrú Yǎwēndé hé Dùālā, yě yǒu qīngshàonián fànzuì de wèntí. Yǒude niánqīngrén táoxué, yǒude tōu dōngxi. Jiātíng zài yùfáng fànzuì zhōng hěn zhòngyào. Fùmǔ yīnggāi guānxīn háizi, chángcháng hé tāmen shuōhuà, yě yīnggāi rènshi tāmen de péngyou. Fùmǔ hái yào gěi háizi hǎo de jiàoyù, bāngzhù tāmen liú zài xuéxiào. Jiātíng kěyǐ bǎohù háizi, ràng tāmen bù shòu huài péngyou de yǐngxiǎng. Wǒ rènwéi, jiātíng hé xuéxiào yīnggāi yìqǐ nǔlì. Zhèyàng, qīngshàonián cái néng zǒushàng zhèngquè de lù.)}

« Dans les villes du Cameroun, comme Yaoundé et Douala, le problème de la délinquance juvénile existe aussi. Certains jeunes sèchent les cours, d'autres volent. La famille est très importante dans la prévention de la délinquance. Les parents doivent se soucier de leurs enfants, leur parler souvent, et aussi connaître leurs amis. Les parents doivent encore donner une bonne éducation à leurs enfants et les aider à rester à l'école. La famille peut protéger les enfants et les empêcher de subir l'influence de mauvais amis. Je pense que la famille et l'école doivent travailler ensemble. Ainsi, les adolescents pourront prendre le bon chemin. »

\textbf{Vérification des contraintes :} structure \zh{有的……有的……} (phrase 2) ; verbe \zh{应该} (phrases 4 et 7) ; vocabulaire de la leçon : \zh{预防}, \zh{关心}, \zh{教育}, \zh{保护}, \zh{影响}.

\textbf{Barème indicatif (sur 10) :} respect du sujet et des contraintes 3 points ; correction grammaticale 3 points ; vocabulaire de la leçon 2 points ; caractères et présentation 2 points.

\textbf{Points de vigilance :}
\begin{itemize}
\item \zh{杜阿拉} (Dùālā) est la transcription courante de Douala ; \zh{雅温得} (Yǎwēndé) pour Yaoundé ;
\item ne pas écrire \zh{家庭是很重要在预防中} : le complément \zh{在预防犯罪中} se place avant l'adjectif ;
\item \zh{让他们不受……的影响} : la négation se place avant \zh{受} ;
\item éviter les phrases trop longues : 6 à 8 phrases courtes et correctes valent mieux qu'une phrase complexe fautive.
\end{itemize}
\end{corrige}

\newpage

\coursec{Fiche bilan}

\begin{popfigure}
\begin{tikzpicture}[mindmap, grow cyclic, every node/.style={concept, font=\small}, concept color=popblue, root concept/.append style={concept color=popdark, font=\small\bfseries, text=white}, level 1/.append style={level distance=3.2cm, sibling angle=51.4}, level 2/.append style={level distance=2.1cm, font=\scriptsize}]
\node[root concept]{青少年犯罪 qīngshàonián fànzuì}
  child[concept color=poporange]{ node {Lexique 词汇}
    child { node[align=center]{犯罪 fànzuì \\ 违法 wéifǎ} }
    child { node[align=center]{原因 yuányīn \\ 影响 yǐngxiǎng} }
  }
  child[concept color=popgreen]{ node {Compréhension 理解}
    child { node[align=center]{Qui ? Où ? \\ Pourquoi ?} }
    child { node[align=center]{Repérer les \\ mots-clés} }
  }
  child[concept color=poppurple]{ node {Radicaux 部首}
    child { node {罪 : 罒 + 非} }
    child { node {法 : 氵+ 去} }
  }
  child[concept color=poppink]{ node {Structures 语法}
    child { node {因为…所以…} }
    child { node {越来越 + adj.} }
  }
  child[concept color=popgold]{ node {Méthode 方法}
    child { node {Lire 2 fois} }
    child { node[align=center]{Souligner \\ les réponses} }
  }
  child[concept color=popblue!70]{ node {Culture 文化}
    child { node[align=center]{Cameroun / Chine : \\ école et famille} }
  };
\end{tikzpicture}
\legende{Carte mentale de la leçon : la délinquance juvénile}
\end{popfigure}

\begin{bilan}
\textbf{1. Lexique clé à connaître par cœur}
\begin{center}
\begin{tabularx}{\linewidth}{|X|X|X|X|}
\hline
\rowcolor{popblueL} Caractères & Pinyin & Nature & Traduction \\
\hline
青少年 & qīngshàonián & nom & adolescent, jeune \\
\hline
犯罪 & fànzuì & verbe/nom & commettre un crime ; crime \\
\hline
违法 & wéifǎ & verbe & enfreindre la loi \\
\hline
偷 & tōu & verbe & voler, dérober \\
\hline
原因 & yuányīn & nom & cause, raison \\
\hline
影响 & yǐngxiǎng & nom/verbe & influence ; influencer \\
\hline
家庭 & jiātíng & nom & famille \\
\hline
教育 & jiàoyù & nom/verbe & éducation ; éduquer \\
\hline
帮助 & bāngzhù & verbe/nom & aider ; aide \\
\hline
解决 & jiějué & verbe & résoudre \\
\hline
问题 & wèntí & nom & problème, question \\
\hline
社会 & shèhuì & nom & société \\
\hline
\end{tabularx}
\end{center}

\textbf{2. Règles de grammaire à connaître par cœur}
\begin{center}
\begin{tabularx}{\linewidth}{|X|X|X|}
\hline
\rowcolor{popblueL} Structure & Exemple & Traduction \\
\hline
因为 A，所以 B & 因为缺少家庭教育，所以一些青少年犯罪。 (Yīnwèi quēshǎo jiātíng jiàoyù, suǒyǐ yìxiē qīngshàonián fànzuì.) & Parce qu'il manque d'éducation familiale, certains jeunes délinquent. \\
\hline
越来越 + adjectif & 这个问题越来越严重。 (Zhège wèntí yuèláiyuè yánzhòng.) & Ce problème devient de plus en plus grave. \\
\hline
应该 + verbe & 我们应该帮助他们。 (Wǒmen yīnggāi bāngzhù tāmen.) & Nous devrions les aider. \\
\hline
Sujet + 对 + nom + 有影响 & 家庭对孩子有很大的影响。 (Jiātíng duì háizi yǒu hěn dà de yǐngxiǎng.) & La famille a une grande influence sur l'enfant. \\
\hline
\end{tabularx}
\end{center}

\textbf{3. Méthodes express}
\begin{itemize}
\item \textbf{Lire un texte} : 1\textsuperscript{re} lecture globale (thème général), 2\textsuperscript{e} lecture en soulignant les mots-clés (personnes, lieux, causes, solutions).
\item \textbf{Répondre aux questions} : repérer dans la question le mot-clé (为什么 \emph{pourquoi}, 谁 \emph{qui}, 怎么样 \emph{comment}), le retrouver dans le texte, recopier la phrase-réponse en l'adaptant.
\item \textbf{Traduire} : identifier d'abord sujet + verbe, puis placer les compléments de temps et de lieu \emph{avant} le verbe.
\item \textbf{Mémoriser les caractères} : décomposer en clé + phonétique (罪 = 罒 filet + 非 ; 法 = 氵eau + 去), écrire de haut en bas, de gauche à droite.
\end{itemize}
\end{bilan}

\begin{aretenir}[Checklist avant le Bac]
\begin{itemize}
\item[$\square$] Je sais lire et écrire : 青少年、犯罪、违法、偷、原因、影响、教育、帮助、解决、社会。
\item[$\square$] Je connais les tons exacts de chaque mot du lexique (fànzuì, yuányīn, yǐngxiǎng…).
\item[$\square$] Je sais décomposer 罪 (罒 + 非) et 法 (氵+ 去) et décrire l'ordre des traits.
\item[$\square$] Je maîtrise la structure 因为…所以… pour exprimer la cause.
\item[$\square$] Je sais utiliser 越来越 + adjectif pour exprimer une évolution.
\item[$\square$] Je sais construire une phrase avec 对…有影响 (avoir une influence sur).
\item[$\square$] Je place correctement les compléments de temps et de lieu avant le verbe.
\item[$\square$] Je sais repérer les mots-clés d'une question de compréhension (谁、什么、为什么、怎么样).
\item[$\square$] Je réponds aux questions par des phrases complètes, pas par un seul mot.
\item[$\square$] Je sais comparer la situation des jeunes au Cameroun et en Chine en 3 phrases simples.
\item[$\square$] Je relis ma production : tons, ordre des mots, ponctuation chinoise 。，？
\end{itemize}
\end{aretenir}

\findecours

\end{document}
